\documentclass[thesis]{cluu}

\usepackage[style=cluu]{biblatex}
\usepackage{booktabs}
\addbibresource{pal.bib}

% --- Packages ---------------------------------------------------------------
% NOTE: xeCJK was removed. There is no CJK text in this thesis, and the package
% searches for CJK fonts on every compile. This was almost certainly the single
% largest contributor to slow builds.
\usepackage{fontspec}          % keep if you compile with XeLaTeX/LuaLaTeX
\usepackage{tabularx}
\usepackage{booktabs}
\usepackage{float}   % in your preamble
\usepackage{multirow}
\usepackage{graphicx}
\usepackage{colortbl}          % \cellcolor; xcolor is loaded via tikz
\usepackage{tikz}
\usetikzlibrary{arrows.meta}
% add to main.tex alongside \cg/\crd (same colors, but for inline text spans)
\newcommand{\hlg}[1]{\colorbox{hmgreen!25}{#1}}
\newcommand{\hlr}[1]{\colorbox{hmred!25}{#1}}


% To cache the TikZ pie chart between runs, uncomment the two lines below and
% enable shell-escape in Overleaf (Menu > Compiler settings). The figure will
% then only be rebuilt when its code changes.
% \usetikzlibrary{external}
% \tikzexternalize[prefix=tikz-cache/]

% --- Custom definitions -----------------------------------------------------
\definecolor{hmgreen}{HTML}{2E8B57}
\definecolor{hmred}{HTML}{C0392B}
\newcommand{\cg}[1]{\cellcolor{hmgreen!#1}}
\newcommand{\crd}[1]{\cellcolor{hmred!#1}}
\newcolumntype{Y}{>{\centering\arraybackslash}X}   % centred flexible column

\begin{document}
\title{LLM-generated Swedish: Linguistic Features, Detection, and Human Judgement}
\author{Gina Welsh}
\supervisor{Birger Moëll, Uppsala University}
% Use \supervisors (with s) if you have more than one

\maketitle

\begin{abstract}

This thesis investigates systematic differences between Swedish text generated
by large language models (LLM-Swedish) and Swedish text written by humans (human-Swedish). We generated an LLM-Swedish dataset containing a formal domain (Bachelor's abstracts) and an informal domain (online forum comments). Three different prompts were created, with two of them designed to be AI feature-evasive. We performed a feature analysis comparing human-Swedish and LLM-Swedish using several linguistic measures, including syntactic complexity, lexical diversity, named entity recognition, sentiment analysis, and punctuation. We found that human- and LLM-Swedish differed significantly across several measures in both the formal and informal domain. In addition, we measured the performance of two AI detection methods (embedding-based classification and LLM-as-judge), finding that LLM-as-judge detectors were evaded by surface features that did not fool the embedding-based classifier. Lastly, we ran a human detection study to determine the ability of Swedish speakers in discerning AI-generated Swedish, finding that Swedish speakers distinguished AI-generated forum comments at above chance (63.1\%), but appear to find abstracts more difficult (52.6\%). Overall, the thesis has several contributions as the first study to systematically measure differences between human- and LLM-Swedish, the effect of LLM-Swedish generated by feature-evasive prompts, and the ability of human Swedish speakers to discern LLM-Swedish text.

\end{abstract}

\tableofcontents

\addchap{Preface}
I would like to thank several people who helped contribute to this thesis. Firstly, my supervisor Birger, who was visibly excited by this topic and helped brainstorm directions as well as sources to consult. Secondly, I would like to thank all human participants of my `AI eller inte' quiz, who generously offered their time uncompensated. I would also like to thank my lecturers at Uppsala University, who provided me with the opportunity to conduct real research in this extremely fast-moving academic field. 
\paragraph{}
I would like to thank all my friends in Sweden whom I met during a very transitional time in my life and who made uprooting my life and moving to the other side of the world a joyous endeavour, particularly Irini, Pavlo, Dennis, Sebastian, and Kārlis. I would also like to thank ICA Nära Folkes Livs for providing me with my extremely specific study snacks. There was that one crucial week where none of them were in stock, but I forgive you.
\paragraph{}
Lastly, I would like to thank my family for supporting me. I moved 15,750 kilometers away from you, but you still trusted me to be okay! I am forever grateful to have this insane privilege to live in Europe and study such an incredible and timely topic. I am very much looking forward to my next visit to you all in Australia.

\chapter{Introduction}

Transformer-based Large Language Models (LLMs) can generate text that very closely emulates human writing to the naked eye, and have reshaped our worldwide informational landscape. As LLMs are adopted more widely, it has become important to discern human-written text from LLM-generated text. This discernment informs plagiarism detection, tracking AI-generated propaganda, and our understanding of which AI-generated writing styles and content dissemination are now heavily represented in online media. Feature-based analysis of LLM-generated text is now a topic of interest, with several key studies established  \parencite{liao2023differentiating, guo2023closechatgpthumanexperts, tang2023sciencedetectingllmgeneratedtexts, munoz-ortiz2024contrasting}. Feature analysis studies that focus on surface-level patterns of AI-generated writing can make AI writing characteristics more transparent to a non-technical audience, which in turn can assist in educating the general public as to the nature of AI-generated writing. Alongside the development of AI detection methods, a parallel line of work has produced prompt-based strategies designed to evade detectors \parencite{krishna2023paraphrasing, lu2024sico, cheng2025adversarialparaphrasinguniversalattack}. An additional line of interest is in the abilities of humans to discern between human-written and AI-generated text \parencite{clark2021all, doruetal}. Comparatively little work, however, examines characteristics and detection of LLM-generated text for languages other than English. As a result, this thesis aims to expand these lines of research to the Swedish language.

\section{Research Questions}
This thesis investigates the following research questions:
\begin{itemize}
    
\item RQ 1: How do human-Swedish and LLM-Swedish texts differ in their surface-level features?

\item RQ 2: How does prompt design affect the detectability of Swedish text by AI detectors?

\item RQ 3: How do human Swedish speakers compare to AI detectors in their ability to discern AI-generated Swedish?

\end{itemize}

\section{Experimental-set up of the thesis}

The thesis is comprised of three experiments:

\begin{enumerate}
    \item A feature analysis comparison between human-written and LLM-generated Swedish text for a formal and informal domain,
    \item A classification detection task for two detection paradigms: an embedding-based classifier and an LLM-as-judge group task,
    \item An LLM-Swedish detection task for human Swedish speakers.
\end{enumerate}

The study focuses on Swedish text within a formal domain (Bachelor's academic abstracts) and an informal domain (internet forum comments).

\section{Thesis contributions}

This thesis makes several contributions. Firstly, it provides the first combined feature-based analysis and detection evaluation of LLM-generated Swedish across two registers, addressing a gap in an Anglocentric body of work that has not covered any North Germanic language in detail. Secondly, it introduces a purpose-built Swedish dataset of paired human and LLM texts spanning a formal register (Bachelor's abstracts) and an informal register (forum comments), generated under three prompt conditions and across three models (GPT-5.2, GPT-5.6, and Mistral Small 3.2). Thirdly, it demonstrates that a feature-informed feature-evasive prompt can shift surface features toward human values but tends to over-correct and therefore cautions against feature-targeted evasion as a standalone strategy. Fourthly, our experiment showed a dissociation between two detection paradigms on the same texts, with surface-level evasion succeeding against LLM-as-judge classifiers while failing against an embedding-based classifier reading distributional structure. Finally, the thesis offers a direct human-versus-model comparison on an identical detection task, finding that while Swedish speakers can detect informal LLM comments above chance but not formal abstracts, LLM judges reliably detect LLM-abstracts but struggle with detecting comments. 

\chapter{Background}

This chapter provides an overview of relevant concepts and related work for the experiments conducted in the thesis: multilingual LLMs, textual feature analysis, LLM text detection by AI tools, and LLM text detection by humans.

\section{Transformer-based LLMs}
Transformer-based large language models (LLMs) are built on a neural network architecture with multi-head attention mechanisms, first proposed in \textcite{vaswani2017attention}. LLMs can generate highly fluent and coherent text as a result of emergent abilities that were not evident in earlier language models, performing to the extent that they are comparable, on the surface and at first sight, to human-written text \parencite{mahowald2024dissociating,wei2022emergent}. Examples of LLMs include OpenAI's GPT-3 \parencite{brown2020language} and GPT-4 \parencite{openai2023gpt4}, Anthropic's Claude \parencite{anthropic2024claude}, DeepSeek-R1 \parencite{DeepSeek-R1ai2024DeepSeek-R1v3} and LLaMa \parencite{touvron2023llama}. Compared to language models before the release of GPT-3, large-scale LLMs contain a very large number of parameters, ranging at the time of writing between tens and thousands of billions \parencite{brown2020language, DeepSeek-R1ai2024DeepSeek-R1v3, fedus2022switch}. These parameters are trained on very large amounts of data from a range of sources, which is often pre-filtered in the training process to improve model performance on different user inputs and to mitigate potential bias \parencite{raffel2020exploring, dodge2021documenting, gehman2020realtoxicityprompts, brown2020language}. The scale and composition of this training data, and in particular its language composition, is central to the questions this thesis raises, and it is to the multilingual case that we now turn.

\section{Multilingual LLMs}
LLMs demonstrate emergent abilities for multiple languages. This multilingual ability raises an important question: how well do these models adapt to non-English languages in terms of their linguistic abilities, reasoning abilities, and cultural knowledge?

\subsection{Asymmetric distributions of language data}
Large-scale multilingual LLMs are very often trained in a manner that is English-language-centric, that is, trained on a dataset with a sharp asymmetry between English and all other represented languages. The developer teams behind the latest versions of large-scale proprietary multilingual LLMs (e.g., OpenAI's ChatGPT and Anthropic's Claude) have not recently released publicly-accessible breakdowns of the languages represented in their models' training data, with the exception of GPT-3 (2020). Figure \ref{fig:gpt3-langdist} visualises the language data breakdown of the GPT-3 model release \parencite{brown2020language}, accessible at its supplementary dataset \footnote{https://github.com/openai/gpt-3/tree/master/dataset\_statistics}. GPT-3's training data multilingual distribution is heavily anglocentric, with approximately 93\% of its data word count in English and the next four most represented languages being French, German, Spanish and Italian, bundled at approximately 4.65\% of the data. Swedish represented a very small fraction of the training data, at only 0.12\%. This thesis implements LLM generators and judges released after GPT-3; however, these models do not have a publicly released language diversity breakdown. Nonetheless, we assume that Swedish continues to make up a very low proportion of language data in latest-release models. The question of how this low proportion would affect the nature of generated Swedish text is a key motivation of this study.

\definecolor{cEng}{HTML}{22366B}
\definecolor{cFra}{HTML}{D1495B}
\definecolor{cDeu}{HTML}{15A18C}
\definecolor{cSpa}{HTML}{8338EC}
\definecolor{cIta}{HTML}{F08A24}
\definecolor{cPor}{HTML}{3A86FF}
\definecolor{cNld}{HTML}{2A9D2F}
\definecolor{cSwe}{HTML}{FFD000}   % Swedish highlight
\definecolor{cOth}{HTML}{9AA1A8}
\definecolor{cRem}{HTML}{C4C9CE}   % remainder wedge in pie
\definecolor{cZoom}{HTML}{EDEFF1}  % zoom connector fill

\begin{figure}[t!]
  \centering
  \resizebox{\textwidth}{!}{%
  \begin{tikzpicture}[line join=round]

  % ---- zoom connector (behind) ----
  \fill[cZoom] (3.3927,0.8599) -- (7.0,4.5) -- (7.0,-4.5) -- (3.3927,-0.8599) -- cycle;
  \draw[gray,line width=0.5pt] (3.3927,0.8599) -- (7.0,4.5);
  \draw[gray,line width=0.5pt] (3.3927,-0.8599) -- (7.0,-4.5);

  % ---- pie ----
  \fill[cEng] (0,0) -- (14.2224:3.5) arc[start angle=14.2224, end angle=345.7776, radius=3.5] -- cycle;
  \fill[cRem] (0,0) -- (-14.2224:3.5) arc[start angle=-14.2224, end angle=14.2224, radius=3.5] -- cycle;
  \draw[black,line width=0.8pt] (0,0) circle[radius=3.5];
  \draw[black,line width=0.6pt] (0,0) -- (14.2224:3.5);
  \draw[black,line width=0.6pt] (0,0) -- (-14.2224:3.5);

  \node[font=\small,text=white] at (-1.2,0.0) {English};
  \node[font=\small,text=white] at (-1.2,-0.6) {92.6\%};

  \node[font=\small\bfseries] at (7.75,4.95) {Other languages: 7.4\%};

  % ---- stacked bar (the 7.90\% non-English remainder, expanded) ----
  \fill[cFra] (7.0,2.4722) rectangle (8.5,4.5000);
  \fill[cDeu] (7.0,0.5582) rectangle (8.5,2.4722);
  \fill[cSpa] (7.0,-0.3418) rectangle (8.5,0.5582);
  \fill[cIta] (7.0,-1.0709) rectangle (8.5,-0.3418);
  \fill[cPor] (7.0,-1.6861) rectangle (8.5,-1.0709);
  \fill[cNld] (7.0,-2.0848) rectangle (8.5,-1.6861);
  \fill[cSwe] (7.0,-2.2215) rectangle (8.5,-2.0848);
  \fill[cOth] (7.0,-4.5000) rectangle (8.5,-2.2215);
  \draw[white,line width=0.5pt] (7.0,2.4722) -- (8.5,2.4722) (7.0,0.5582) -- (8.5,0.5582) (7.0,-0.3418) -- (8.5,-0.3418) (7.0,-1.0709) -- (8.5,-1.0709) (7.0,-1.6861) -- (8.5,-1.6861) (7.0,-2.0848) -- (8.5,-2.0848) (7.0,-2.2215) -- (8.5,-2.2215);
  \draw[black,line width=0.9pt] (7.0,4.5) rectangle (8.5,-4.5);

  % ---- leader labels ----
  \node[anchor=west,font=\small] (LFrench) at (9.15,3.4861) {French \textemdash\ 1.81\%};
  \draw[gray,line width=0.5pt] (LFrench.west) -- (8.5,3.4861);
  \node[anchor=west,font=\small] (LGerman) at (9.15,1.5152) {German \textemdash\ 1.46\%};
  \draw[gray,line width=0.5pt] (LGerman.west) -- (8.5,1.5152);
  \node[anchor=west,font=\small] (LSpanish) at (9.15,0.1082) {Spanish \textemdash\ 0.77\%};
  \draw[gray,line width=0.5pt] (LSpanish.west) -- (8.5,0.1082);
  \node[anchor=west,font=\small] (LItalian) at (9.15,-0.7063) {Italian \textemdash\ 0.61\%};
  \draw[gray,line width=0.5pt] (LItalian.west) -- (8.5,-0.7063);
  \node[anchor=west,font=\small] (LPortuguese) at (9.15,-1.3785) {Portuguese \textemdash\ 0.52\%};
  \draw[gray,line width=0.5pt] (LPortuguese.west) -- (8.5,-1.3785);
  \node[anchor=west,font=\small] (LDutch) at (9.15,-1.8985) {Dutch \textemdash\ 0.34\%};
  \draw[gray,line width=0.5pt] (LDutch.west) -- (8.5,-1.8854);
  \node[anchor=west,font=\bfseries\small,fill=cSwe,rounded corners=2pt,inner sep=3pt] (LSwedish) at (9.15,-2.4185) {Swedish \textemdash\ 0.12\%};
  \draw[-{Latex[length=2mm]},semithick,cSwe] (LSwedish.west) -- (8.5,-2.1532);
  \node[anchor=west,font=\small] (LOther) at (9.15,-3.3608) {Other \textemdash\ 1.771.\%};
  \draw[gray,line width=0.5pt] (LOther.west) -- (8.5,-3.3608);

  \end{tikzpicture}%
  }
  \caption[Language distribution of the GPT-3 training data]{%
    Language distribution of GPT-3 training data by word
    count \parencite{brown2020language}.}
  \label{fig:gpt3-langdist}
\end{figure}

\subsection{The curse of multilinguality}
The current linguistic performance of large-scale multilingual LLMs appears to depend on the number of languages that are represented in its training data as well as the model's parameter scale. \textcite{conneau2020unsupervised} named the `curse of multilinguality', a phenomenon where adding more languages to the proportion of training data actually decreases its multilingual abilities by overcrowding' its parameters. The study suggested scaling up of parameters as a way of mitigating this effect, so that proportions of lower-represented languages can reach a satisfactory threshold and improve model performance on those languages. Existing LLM benchmarks have been developed mainly for English datasets, and while some cover a range of high- to mid-resource languages, there is a significant lack of LLM evaluation datasets for other languages, particularly for low-resource ones \parencite{wu2025survey}. An extension of LLM evaluation in other languages could potentially provide insight into the `curse of multilinguality', by investigating the threshold needed for acceptable LLM performance in lower-resource languages. This thesis implements large-scale multilingual LLMs in order to provide some insights as to how its multilinguality could contribute to differences between and detection of human- and LLM-Swedish.

\subsection{Swedish in the LLM landscape}
Several multilingual LLMs have the ability to generate and analyse Swedish. In addition, there are Nordic language LLM resources, notably AI Sweden's LLM GPT-SW3 \parencite{ekgren2024gptsw3}. The SW3-GPT models comprise of decoder-only transformers trained on 320 billion tokens of Swedish, Norwegian, Danish, Icelandic, English and programming code  drawn from the Nordic Pile, a curated 1.2TB corpus spanning the major North Germanic languages and high-quality English data \parencite{ohman2023nordic} Alongside model development, evaluation benchmarks tailored to Swedish include ScandEval (Nielsen, 2023), which offers metrics for Swedish named entity
recognition (NER), sentiment classification, and linguistic
acceptability; a Sweden-related factual-knowledge benchmark by Kunz
(2025); SweSat (Kurfali et al., 2025), derived from the Swedish
university entrance exam (Högskoleprovet); and the Swedish Medical LLM
Benchmark \parencite{moell2025swedishmed}, a framework for evaluating
LLM capabilities on Swedish medical texts. However, while these benchmarks test the capabilities of LLMs in generating and analysing Swedish, there is no current benchmark for the detection and characterisation of LLM-generated Swedish text.
% TODO: add ScandEval, Kunz, SweSat entries to pal.bib

\section{Feature-based analysis of LLM-generated text}
 Linguistic feature analysis, that is, the quantitative analysis of structural properties of text, has its roots in hand-calculated stylometric analysis \parencite{mosteller1963} before progressing to computational methods \parencite{sundararajan-woodard-2018-represents, belvisi2020forensicauthorshipanalysismicroblogging, huang2025authorship}. This thesis focuses solely on computational methods of linguistic feature analysis.

\subsection{Previous work: feature analysis of LLM text}
 In the transformer era, several studies have contributed linguistic feature analysis of LLM-generated text. A key study is \textcite{munoz-ortiz2024contrasting}, who compared the linguistic components of
human and LLM-generated news articles by measuring vocabulary, part-of-speech distributions, dependency arc lengths, constituent lengths, and emotional categories. Across six LLMs, their results showed significant differences between the human-written and LLM-generated text on all textual analysis metrics, finding that human texts have a higher variety of vocabulary, shorter constituents, and more variable sentence length distributions. Human texts portrayed overall a more negative affect than LLM texts, with stronger negative emotions such as fear and disgust. The more positive affect in the LLM texts could be connected inherent positive bias in LLMs \parencite{wu2025survey, giorgi2021characterizing, markowitz2024linguistic}, speculated to be the result of filtration of violent or sexual content from training data \parencite{liao2023differentiating} to align with model safety standards\parencite{ouyang2022training}. Other LLM feature analysis studies include \textcite{tang2023sciencedetectingllmgeneratedtexts}, who found LLM-generated texts to be `less emotional and objective' than human-authored text, as well as more formal and structured. Studies focused on languages other than English include \textcite{park2025katfishnet}, which examines spacing patterns, part-of-speech diversity and comma usage for LLM-generated Korean texts, finding that LLM Korean texts strictly enforce spacing patterns in Korean that humans are looser to apply, and that LLM Korean texts have considerably higher rates of commas and part-of-speech diversity. In addition, \textcite{alshaibani2025arabicaifingerprintstylometric} measured distinctive linguistic patterns differentiating AI-generated Arabic text, finding that LLMs
produce text with different vocabulary diversity and distribution from human texts; and \textcite{milicka2024benchmark}, who benchmarked stylistic variation of LLM-generated Czech, finding systematic differences that distinguish the LLM texts from humans', and finding that these differences shift in similar directions across multiple LLMs, and that these shifts were sharper in Czech than in English. 

Overall, these studies find significant differences between human- and LLM-text across the studied languages, with recurring cross-lingual patterns including differences in vocabulary distribution, punctuation markers, and sentiment affect. These findings indicate that although LLMs generate text that is linguistically sophisticated enough to appear as human-written text to the naked eye, there appear to be systematic differences in the style of the texts cross-linguistically.

\subsection{Register and genre variation}
There is evidence that differences between human and LLM-generated written text also differ across genres. Within the medical domain,  \textcite{liao2023differentiating} compared
human-written and ChatGPT-generated medical texts by their vocabulary
distribution, part-of-speech tag distribution, dependency arcs,
sentiment analysis, and perplexity. Their results found that human-written medical texts contained more concrete information, had higher
vocabulary diversity, and contained information that human
experts would consider more useful. The ChatGPT-generated
medical texts appeared more fluent and logical, but expressed
more generalised terminology rather than specific, effective
information for the specific medical case's context. This could point to the tendency for LLMs to excel at formal competence (that is, linguistic fluency and coherence) rather than functional competence, which includes situational modelling, in order to tailor the output to the context of the medical case \parencite{mahowald2024dissociating}. Genre can also affect detector performance, with \textcite{bao2024fastdetectgpt} finding that supervised detectors performed comparably on in-distribution English news but significantly worse on English scientific writing and German writing. Overall, it appears that register and genre can affect LLM texts significantly, and therefore motivate a cross-register comparison.

\section{Detecting LLM-generated text}
The development of transformer-era LLMs has resulted in numerous studies of the performance of AI detectors as well as attacks to evade them through prompt design. Earlier studies investigated the performance and mechanics of neural detectors for fake news \parencite{zellers2020defendingneuralfakenews, pagnoni-etal-2022-threat} as well as domain-tuned detectors such as RoBERTa for GPT-2 technical texts \parencite{rodriguez-etal-2022-cross}. More recent detection approaches include zero-shot statistical methods  that exploit probability-curvature signals \parencite{mitchell2023detectgpt, bao2024fastdetectgpt} and those that employ LLM-as-judge classification, that is, using LLMs themselves as detectors \parencite{zhu-etal-2023-beat, zheng2023judging}.

\subsection{Supervised and two-sample-test approaches}
Detectability can be measured using supervised classifier approaches that train a classifier on labelled human and machine text. The two-sample test calculates whether two samples are drawn from the same distribution, using the classifier's held-out accuracy as a test statistic, and the null hypothesis is rejected where the classifier accuracy is above chance \parencite{lopezpaz2017revisiting}. In the application of discerning LLM- from human-Swedish, this approach measures how far apart human and LLM Swedish sit in representation space.

\subsection{LLM-as-judge}

An LLM-as-judge detection task deploys reasoning LLM models to classify whether a given text is LLM-generated, with a model's reasoning rationale given as a by-product. This approach originates in general LLM evaluation, where LLMs were used as cost-effective proxies for human raters \textcite{zheng2023judging}. As judges, LLM show evidence of being biased towards their own generated text, such as distinguishing their own output from human text and text generated from other LLMs, as shown in \textcite{panickssery2024llmevaluators}, where models GPT-4 and Llama 2 distinguished their output at above-chance rates and this self-recognition scoring their generations more favourably.

\subsection{Evading detection}

 The emergence of AI detectors has resulted in LLM prompting methods designed to evade such detectors. Several studies detail evasive techniques against AI detection, such as paraphrasing the LLM response \parencite{sadasivan2025aigeneratedtextreliablydetected, krishna2023paraphrasingevadesdetectorsaigenerated, huang-etal-2025-tempparaphraser} and adversarial perturbation \parencite{zhou2024humanizing, cheng2025adversarialparaphrasinguniversalattack}. \textcite{lu2024largelanguagemodelsguided} generated LLM text with linguistic-feature-informed evasive prompts for three genres (essays, open-ended question answering, and fake reviews), with their evasive prompt design achieving strong performance in evading six detectors across the three genres, and significantly outperforming a baseline human-text-paraphrasing, demonstrating that feature-based prompt construction can be a strong feature-evasive technique.

\subsection{Human detection of LLM-generated text}
Several studies have recruited human participants for an AI-generated text detection task to measure their AI-detection abilities. \textcite{waltzer2024spot} recruited university instructors and students to annotate academic texts as AI- or human-written, finding that instructors guessed with 70\% accuracy and students with 60\% accuracy, and that neither LLM usage nor content expertise correlated with higher performance. \textcite{doruetal} conducted an AI detection study for medical texts in German, employing medical experts and non-experts to judge the differences between human-written medical academic texts and texts generated by ChatGPT 3.5. Their participants were more likely to label a text as human-written when they encountered any spelling or grammatical errors, more varied sentence structure, or specific terminology, and more likely to label a text as AI where there were unnecessary details, monotonous sentence structures, and a superficiality or generic tone. \textcite{zaitsu2025stylometry} conducted a study with Japanese speakers and texts generated by multiple versions of GPT, finding that humans had a limited ability to correctly classify texts (overall accuracy = 58\%) and also over-labelled human texts as AI, relying on surface-level impressions like word endings, conjunctions, and use of punctuation marks. \cite{milicka2025humans} conducted the study with Czech speakers, finding a similar level of accuracy (55.4\%) and a usage of incorrect assumptions, such as the assumption that humans write more readable texts than AI. Overall, current research shows that humans do struggle with distinguishing human- and AI-text, and sometimes have uninformed impressions about what AI-generated text looks like.

\section{Swedish-language work and the research gap}
Some existing studies include AI-generated Swedish in their scope, including
research into AI-generated Swedish detection by AI tools
\parencite{landberg2024aiwolf} as well as the shift in the structures of
Swedish language academic writing and the LLM-as-judge approach for Swedish
university-level texts \parencite{hintze2026aiuse}. Landberg (2024) tested
commercial AI detectors and LLM-as-judges on 400 human-written Swedish
news articles and 400 ChatGPT-4 generated texts, alongside an equivalent
English dataset. Their results found that all three chatbots performed at or
below chance on the Swedish data, with ChatGPT 3.5 reaching only 22\% accuracy
and not assigning the AI label at all, and GPTZero labelling all Swedish texts as human written. These results contrasted with English, where AI detectors Copyleaks and Smoding reached very high accuracy levels on the English texts. Overall, this indicates a difference in AI detection performance between English and Swedish (a mid-resource language). It is worth noting that these results are derived from older GPT models, and therefore warrant a replication with more recent models.

Hintze (2026) investigated the use of AI tools in influencing submitted student writing by analysing 17,007 Swedish bachelor's and master's theses
deposited in DiVA between 2020 and 2025, and studying the differences between texts written before November 2022 (the release of ChatGPT) and those written afterwards. Each thesis was scored by an LLM-as-judge
under five dimensions: grammar, vocabulary, academic quality, wordiness, and an AI-likelihood judgement. Their results found statistically significant
post-2022 shifts on all four language-quality dimensions, within all topical clusters, indicating a good chance of LLM influence on Swedish academic texts. As a result, the study provided demonstrated evidence that LLMs have affected the structural properties of texts submitted by Swedish students. As a caveat, the study also noted that the AI-likelihood judgement could be a result of the model's own reasoning process, rather than an actual distinction in the textual properties of LLM texts, and therefore recommends an analysis of
linguistic feature dimensions rather than solely relying on LLM-as-judge reasoning.

Neither of these works characterise the differences between LLM- and human-Swedish on the surface level, the effect of evasive prompt design on AI detection, or the ability of human Swedish speakers to detect AI text. As a result, this thesis is motivated to expand on these works and provide more insight into LLM performance on Swedish.
\chapter{Methodology}

The methodology of this thesis is comprised of three main parts: 1) a feature analysis comparison between human-Swedish and LLM-swedish, 2) a construction of a feature-evasive prompt and their performance in evading AI-detectors, and 3) a comparison between AI tools and human Swedish speakers in detecting LLM-Swedish.

\section{Data collection and creation}

The data amassed for this thesis involved the collection of human-written Swedish and the generation of LLM-Swedish across two registers (formal and informal).

\begin{table}[H]
  \centering
  \small
  \begin{tabular}{llrl}
    \toprule
    Register & Condition & $n$ & Source / model \\
    \midrule
    \multicolumn{4}{l}{\textit{Formal}: Bachelor's abstracts}\\
    \addlinespace[2pt]
    & \texttt{human}  & 170 & DiVA portal \\
    & \texttt{llm\_baseline} & 170 & GPT-5.2 \\
    & \texttt{generic\_evasive} & 170 & GPT-5.2 \\
    & \texttt{feature\_evasive} & 170 & GPT-5.2 \\
    \addlinespace[2pt]
    & \textit{Subtotal} & \textit{680} & \\
    \midrule
    \multicolumn{4}{l}{\textit{Informal}: forum replies}\\
    \addlinespace[2pt]
    & \texttt{human}   & 1{,}149 & Reddit, Flashback \\
    & \texttt{llm\_baseline}& 1{,}149 & GPT-5.2 \\
    & \texttt{generic\_evasive} & 1{,}149 & GPT-5.2 \\
    & \texttt{feature\_evasive} & 1{,perplexity}149 & GPT-5.2 \\
    \addlinespace[2pt]
    & \textit{Subtotal} & \textit{4{,}596} & \\
    \midrule
    & \textbf{Total} & \textbf{5{,}276} & \\
    \bottomrule
  \end{tabular}
    \label{tab:data-overview}
  \caption[Overview of the collected and generated data]{%
    Overview of pairwise comparison data by register and condition.}
  \begin{minipage}{\textwidth}
    \footnotesize
  \end{minipage}
\end{table}

\subsection{Human-written formal data}

For the formal register, we collected 170 Swedish Bachelor's-thesis titles, abstracts and key words from the Digitala Vetenskapliga Arkivet (DiVA) portal, sampled across 30 research topics including humanities, science, law, engineering and agricultural sciences, with a median abstract length of 182 words. All abstracts predated 2017 to avoid abstracts generated by or contaminated with LLM writing assistance.

\subsection{Human-written informal data}

The informal data comprised 136 forum thread questions and 1,149 forum-comment top-level replies (median length = 31 words) across 136 Swedish-language question threads sourced from Reddit and Flashback, each paired with its parent question. Flashback comments spanned 77 threads across a range of topics, from relationships and dating, health and medical, technology and computing, politics and society, everyday questions, and illicit substances. Reddit comments were sourced specifically from four subreddits (\texttt{sweden}, \texttt{svenskpolitik}, \texttt{stockholm}, \texttt{gothenburg}). To mitigate the possibility of LLM influence on the data, only questions and replies with a timestamp before January 2017 were sourced. To ensure that the post's titles were in question form and in Swedish, the script checked for the presence of a question mark (``?'') and used the \texttt{langdetect} Python library to identify the post as being written in Swedish (``sv''). All hyperlinks and embedded images were removed from the comment data, to focus purely on the structural properties and detection of text.

\subsection{LLM data generation}

For the feature analysis, all LLM texts were produced by OpenAI's multilingual large-scale LLM GPT-5.2. The multilingual GPT-5.2 model was chosen since at the time the study commenced, it was a widely-used mainstream model. Our study did not generate text with Scandinavian-based GPT-SW3 due to accessibility issues at the time of writing; however, a Scandinavian-based LLM study would be insightful future work. Two additional smaller sets of LLM-Swedish were generated with GPT-5.6, a later GPT version and the latest model at the time of the study's end, as well as Mistral Small 3.2; both text datasets were generated with the three LLM prompts. All generators were configured with the exact same prompts and a temperature of 1.0, the default temperature setting of non-reasoning GPT-5.2. To test the effect of temperature as a potentially confounding variable, we sampled texts generated with a temperature of 0.5, finding a negligible effect on structural properties (a similar finding also found for Czech in \cite{milicka2024benchmark}). The texts generated by GPT-5.6 and Mistral 3.2 were also used as additional samples to test the robustness of AI tool detection study.

\subsubsection{Evasive prompt designs.}
Our LLM dataset was split into three parts with three prompts: 1) a baseline
LLM prompt that simply instructs the LLM to generate text in the form of a
Bachelor's abstract or forum comment, 2) a `generic' evasive prompt
\texttt{generic\_evasive} that appends an instruction to the baseline to
`write like a human', and 3) an evasive prompt that uses the preliminary
results of the feature analysis as instructions to intentionally evade an AI
detector (\texttt{feature\_evasive}). The \texttt{feature\_evasive} prompt
instructs the model to reuse key terms verbatim rather than varying them, to
write predominantly short sentences with fewer commas and no dashes, to
minimise hedging and epistemic expressions, to prefer short everyday words,
and to name specific entities where possible; the informal version adds
instructions on colloquial forms and discourse particles. Full prompts for
both registers are given in Appendix \ref{app:prompts}
(Tables \ref{tab:prompt-conditions} and \ref{tab:prompt-conditions-informal}),
and examples of generated abstracts and forum comments in Appendix
\ref{app:examples} (Tables \ref{tab:examples-formal} and
\ref{tab:examples-informal}). These prompts were designed to inspect the
structural differences of the LLM-Swedish under different prompt conditions,
and to test their detectability. All prompts were written in Swedish.


%============================================================
% TABLE 3 — Informal register examples
%============================================================


\section{Experiment 1: textual analysis}

\begin{table}[H]
\centering
\setlength{\tabcolsep}{4pt}
\renewcommand{\arraystretch}{1.05}
\caption{Textual analysis features by category.}
\label{tab:textual-analysis-features}
\begin{tabular}{lp{0.62\linewidth}}
\toprule
\textbf{Category} & \textbf{Features} \\
\midrule
Syntactic complexity & Dependency distribution \& distance, parse tree depth, subordinate-clause rate, POS distribution \\
Lexical Diversity     & Type token ratio, MATTR, MTLD, N-gram repetition rate, word length \\
Pragmatic/discourse  & Hedging, negation rate, epistemic markers \\
Punctuation         & Counts: commas, full stops, ellipsis, dashes, colons, semi-colons \\
NER                  & Named entity label distribution \\
Affective (informal only)           & Polarity magnitude, affective extremity, objectivity/subjectivity \\
Significance testing   & Jensen--Shannon divergence (primary), Wasserstein (ordered) \\
\bottomrule
\end{tabular}
\end{table}

Our linguistic feature analysis aimed to measure significant surface-level differences between the human- and LLM-Swedish. The key goal was to determine which features contributed to a high AI surface signal, and which features were altered strongly by different prompt conditions. We compared the human and LLM dataset pairs across 21 features for the formal data, and 30 features for the informal data, including syntactic complexity, lexical diversity, pragmatic markers, named entity recognition, and punctuation.

\subsection{Syntactic Complexity}
For all data points, syntactic complexity was measured using the four measures
below. Part-of-speech and dependency parses were obtained with spaCy's
\texttt{sv\_core\_news\_lg}\footnote{https://github.com/explosion/spacy-models/releases/tag/sv\_core\_news\_lg-3.8.0}
Swedish-language model. Each measure was computed per sentence and averaged over the text.


\subsubsection{Mean dependency distance}
The dependency distance of a token is the absolute difference in linear position
between that token and its syntactic head; the mean dependency distance (MDD) is
the average over all dependency arcs in a sentence \parencite{liu2008dependency}. For
a sentence of $n$ tokens (the root token has no head, giving $n-1$ arcs):
\begin{equation}
\mathrm{MDD} = \frac{1}{n-1}\sum_{i=1}^{n-1}
  \bigl|\,\mathrm{pos}(w_i) - \mathrm{pos}\!\left(\mathrm{head}(w_i)\right)\bigr|,
\end{equation}
where $\mathrm{pos}(\cdot)$ returns the linear index of a token and
$\mathrm{head}(w_i)$ its governor. We chose this measure as larger mean dependency values indicate more non-local dependencies and a greater processing load, measures that indicate higher complexity within a sentence.

\subsubsection{Parse-tree depth}
Parse-tree depth measures how deeply embedded a sentence's syntactic structure is.
Each token in a dependency tree sits at some number of edges below the root, with the
root itself sitting at depth $0$, and a token whose head is the root sitting at depth $1$. Formally, for a token
$w$ in a sentence $S$ with root $r_S$,
\begin{equation}
\mathrm{depth}(w) =
\begin{cases}
0, & \text{if } w = r_S,\\[2pt]
\mathrm{depth}\!\left(\mathrm{head}(w)\right) + 1, & \text{otherwise.}
\end{cases}
\end{equation}
The depth of the sentence is the depth of its deepest token, equivalently the
length in edges of the longest root-to-leaf path:
\begin{equation}
\mathrm{Depth}(S) = \max_{w \in S} \mathrm{depth}(w).
\end{equation}
The measure reported for a document $D$ is the mean over its $|D|$ sentences:
\begin{equation}
\overline{\mathrm{Depth}}(D) = \frac{1}{|D|}\sum_{S \in D} \mathrm{Depth}(S).
\end{equation}

We chose this measure since deeper trees have hierarchical embedding, or nesting, which indicates a more complex sentence. Mean parse-tree depth also complements mean dependency distance.

\subsubsection{Subordinate clause rate}
The subordinate clause rate captures the degree of clausal embedding. Subordinate
clauses were identified from dependency relations introducing dependent clauses
(as marked by Universal Dependency tags \texttt{advcl}, \texttt{acl}, \texttt{ccomp}, \texttt{xcomp}, \texttt{csubj}), and
the rate was reported per sentence:
\begin{equation}
\mathrm{SubRate} = \frac{\#\,\text{subordinate clauses}}{\#\,\text{sentences}}.
\end{equation}

We selected this measure as higher subordinate clause rates correlate with sentence length and complexity.

\subsubsection{Sentence length and token count per sentence}
Tokens per sentence measures the average length of a sentence. Sentence count indicates the number of sentences per document, using spaCy's sentence segmentation. Both measures were calculated in order to separate the effect of producing more sentences from that of producing longer ones. In some cases, a higher mean sentence count can in fact indicate lower sentence complexity if the mean token length is also lower.

\subsection{Lexical Diversity}
Lexical diversity, that is, the richness of vocabulary within a text, was measured using three related but distinct measures: type-token ratio, MTLD and MATTR.

\subsubsection{Type-token ratio (TTR)}
TTR measures the number of unique tokens (types) $V$ divided by the
total number of tokens $N$:
\begin{equation}
\mathrm{TTR} = \frac{V}{N}.
\end{equation}
TTR decreases with text length, so it is comparable only across
texts of similar length. As a result, we also implemented MTLD and MATTR (defined below), in order to track the distribution of unique vocabulary instances within the texts in a way that also controls for text length.

\subsubsection{Measure of Textual Lexical Diversity (MTLD)}
MTLD is the mean length of contiguous token sequences that maintain a running TTR
above a fixed threshold $\tau$ \parencite{mccarthy2010mtld}. While processing the text
token by token, a \emph{factor} is completed each time the running TTR falls to
$\tau = 0.720$, at which point the counter resets. With $F$ as the number of factors, computed once left-to-right and once right-to-left:
\begin{equation}
\mathrm{MTLD} = \frac{1}{2}\left(\frac{N}{F_{\rightarrow}} + \frac{N}{F_{\leftarrow}}\right).
\end{equation}
A higher MTLD score indicates that lexical diversity is sustained over longer spans.

\subsubsection{Moving-Average Type--Token Ratio (MATTR)}
MATTR averages the TTR computed within a fixed-length window that slides across the
text one token at a time, removing TTR's dependence on total length
\parencite{covington2010cutting}. For a window of $w$ tokens over a text of $N$
tokens there are $N-w+1$ windows; let $V_i$ be the number of types in the window
beginning at position $i$:
\begin{equation}
\mathrm{MATTR}(w) = \frac{1}{N-w+1}\sum_{i=1}^{N-w+1}\frac{V_i}{w}.
\end{equation}
We used a window of $w=50$ tokens as a standard window cut-off. Where a text was below 50 tokens, the method fell back to a plain type-token ratio calculation.

\subsection{Perplexity}

Perplexity is a measure of how many next-token candidates a language model finds plausible on average, with a lower perplexity score indicating fewer candidates (correlated with more predictable language) and higher perplexity indicating a wider range of candidates (correlated with less predictable language). Formally, perplexity is a model's average uncertainty about the next token, that is, the exponentiated average per-token cross-entropy of a language model on text. The perplexity equation is defined as:

\begin{equation}
\text{PPL}(W) = \exp\left(-\frac{1}{N}\sum_{i=1}^{N} \log P(w_i \mid w_1, \dots, w_{i-1})\right)
\label{eq:perplexity}
\end{equation}

Where:
\begin{enumerate}
    \item $W = (w_1, w_2, \dots, w_N)$ is the token sequence,
    \item $N$ is the number of tokens,
    \item $P(w_i \mid w_1, \dots, w_{i-1})$ is the model's predicted probability of token $w_i$ given all preceding tokens,
    \item $\log$ is natural log (matching `exp`).
\end{enumerate}
\paragraph{}
We chose perplexity as a comparison of fluency, coherency, and overall predictability between human- and LLM-Swedish. For our perplexity calculation we used AI Sweden's \texttt{gpt-sw3-356m}\footnote{model accessible at https://huggingface.co/AI-Sweden-Models/gpt-sw3-356m}, a causal language model trained natively on a Scandinavian language-majority dataset. We calculated document-level perplexity with a  1024-token sliding-window, as per HuggingFace's perplexity guide.\footnote{guide accessible at https://huggingface.co/docs/transformers/perplexity} In addition, we calculated the burstiness of the text, that is, the standard deviation of per-sentence perplexity within a document, to track the variation of coherence and fluency among different sentences within a text. Sentences were split by a rule-based spaCy sentence tokeniser, and documents with fewer than two usable sentences were excluded (n=6 forum comments excluded from analysis).


\subsection{Pragmatic Markers}
Pragmatic markers were measured through lexicon-based counting: for each category a
curated Swedish word list was matched against the lemmatised text, and the
category rate was normalised per 1000 tokens to allow comparison across texts of
different length:
\begin{equation}
\mathrm{rate}(c) = \frac{\#\,\text{matches of category } c}{N}\times 1000 .
\end{equation}
Five categories were distinguished:

\begin{enumerate}
    \item \textbf{Hedges}, which mark a speaker's commitment to a proposition, e.g. \textit{kanske} (`perhaps'), \textit{ganska} (`quite'),
    \item \textbf{Boosters}, which amplify commitment, force or certainty. Examples include \textit{verkligen} (`really'/`truly'), \textit{absolut} (`absolutely'),
    \item \textbf{Negation}, which marks negated content, e.g. \textit{inte} (`not'), \textit{aldrig} (`never'), \textit{ingen} (`no one'),
    \item \textbf{Epistemic markers}, which mark the speaker's assessment of the likelihood or evidence of a statement, e.g. \textit{troligen} (`credibly'/`plausibly'), \textbf{möjligen} (`possibly'), \textbf{tydligen} (`clearly'
    ), 
    \item \textbf{Discourse markers}, which structure or connect discourse structures, e.g. \textit{alltså} (`So', `well'), \textit{ju} (most closely translated in English to `as you know'), \textbf{liksom} (`like')
\end{enumerate}

The discourse marker rates and types were only measured in the informal forum comment data, as very few of these markers were found in the formal abstract data. 

\subsection{Punctuation}
For punctuation measurement, we counted commas, full stops,
ellipses, dashes, colons, semicolons, question marks, and exclamation marks, each
normalised per 1000 tokens:
\begin{equation}
\mathrm{punct}(p) = \frac{\#\,\text{occurrences of mark } p}{N}\times 1000 .
\end{equation}

\subsection{Named Entity Recognition}
Named entity recognition (NER) identifies and classifies spans of text that refer
to real-world entities. We used it as a proxy for the specificity and concreteness
of a text: a text that names more places, people, organisations, and
quantities are more concrete and specific, and the opposite indicates vagueness. Entities were tagged with spaCy's
\texttt{sv\_core\_news\_lg} model and mapped onto the Stockholm--Ume\aa{} Corpus
tagset \parencite{suc3}. For each category $c$ and for entities overall,
the rate was normalised per 1000 tokens to allow a comparison across texts of
different length:
\begin{equation}
\mathrm{NER}(c) = \frac{\#\,\text{entities of type } c}{N}\times 1000 .
\end{equation}

\begin{table}[h]
\centering
\begin{tabular}{ll}
\toprule
\textbf{Tag} & \textbf{Entity type} \\
\midrule
LOC & Locations \\
MSR & Measure / numerical expression \\
ORG & Organization \\
PRS & Persons \\
TME & Time \\
WRK & Works of art and other artefacts \\
\bottomrule
\end{tabular}
\caption{Named entity categories, following the Swedish NER annotation
guidelines \parencite{sprakbanken2020ner}.}
\label{tab:ner-categories}
\end{table}

\subsection{Sentiment analysis}
Our sentiment analysis was scored with \texttt{KBlab/robust-swedish-sentiment-multiclass}, a Swedish Megatron-BERT-large model fine-tuned for 3-class sentiment classification (POSITIVE / NEUTRAL / NEGATIVE), accessed via the Hugging Face \texttt{transformers} library. Analysis was restricted to the informal register, since preliminary analysis of the formal register was near-uniformly neutral and thus cost-ineffective for our research.

Let $p_{\text{pos}}$, $p_{\text{neu}}$, and $p_{\text{neg}}$ denote the classifier's
predicted probabilities for the POSITIVE, NEUTRAL, and NEGATIVE classes
($p_{\text{pos}} + p_{\text{neu}} + p_{\text{neg}} = 1$), and let $\hat{y}$ denote the
top predicted label. For each document we derived:

\begin{align}
  \text{signed\_polarity} &= p_{\text{pos}} - p_{\text{neg}}, \quad \in [-1, 1] \label{eq:signed-polarity}\\[4pt]
  \text{affective\_extremity} &= 1 - p_{\text{neu}}, \quad \in [0, 1] \label{eq:affective-extremity}\\[4pt]
  \text{is\_subjective} &= \mathbb{1}[\hat{y} \neq \text{NEUTRAL}], \quad \in \{0, 1\} \label{eq:is-subjective}
\end{align}

Signed polarity is the general polarity on a positive–negative spectrum, where +1 is fully positive and −1 fully negative. Affective extremity is the total probability mass the classifier places on the non-neutral (positive or negative) classes, ranging from 0 (fully neutral) to 1 (fully polarised), which captures how emotionally charged a text is regardless of the direction of that emotion. Subjectivity is a binary indicator that is 1 when the top predicted label is not NEUTRAL and 0 otherwise, marking whether a document expresses any sentiment at all.

\subsection{Feature effect size and significance testing}
For feature effect size, each feature was compared between the paired human and
LLM texts. For pair $i$, we define the difference as
\begin{equation}
d_i = x_i^{\mathrm{human}} - x_i^{\mathrm{LLM}},
\end{equation}
so that a positive value indicates a higher feature value in the human text.
Effect size was quantified with Cohen's $d_z$ for paired designs (mean and
standard deviation of paired differences) and the matched-pairs rank-biserial
correlation. Cohen's $d_z$ expresses the mean paired difference in units of its
own standard deviation \parencite{cohen1988, lakens2013}:
\begin{equation}
d_z = \frac{\bar{d}}{s_d}, \qquad
\bar{d} = \frac{1}{n}\sum_{i=1}^{n} d_i, \qquad
s_d = \sqrt{\frac{1}{n-1}\sum_{i=1}^{n}\left(d_i - \bar{d}\right)^2}.
\end{equation}

For significance testing, each feature was tested with a paired Wilcoxon
signed-rank test \parencite{wilcoxon1945}, since the feature distributions were
presumed as not normal ones. We used \texttt{SciPy} to compute the below. For $n$
pairs, let $R_i = \operatorname{rank}(|d_i|)$ be the rank of the absolute
difference, where pairs with $d_i = 0$ were discarded and $n$ reduced
accordingly, and tied values were given their average rank. The test statistic
sums the ranks of the positive differences, and $W^{-}$ correspondingly sums the
ranks of the negative ones:
\begin{equation}
W^{+} = \sum_{i:\,d_i > 0} R_i ,
\qquad
W^{-} = \sum_{i:\,d_i < 0} R_i .
\end{equation}
The matched-pairs rank-biserial correlation $r_{\mathrm{rb}}$ gives a
non-parametric, bounded ($-1 \le r_{\mathrm{rb}} \le 1$) effect size derived from
the signed-rank sums \parencite{kerby2014}:
\begin{equation}
r_{\mathrm{rb}} = \frac{W^{+} - W^{-}}{W^{+} + W^{-}},
\qquad
W^{+} + W^{-} = \frac{n(n+1)}{2}.
\end{equation}
We chose this measure as the sum of the two rank totals is fixed by $n$, which
keeps the ratio bounded and therefore comparable across features. To control the
false discovery rate across the many features tested within each dataset and
condition, $p$-values were adjusted with the Benjamini--Hochberg (BH) procedure
\parencite{benjamini1995}: for $m$ ordered $p$-values
$p_{(1)} \le \dots \le p_{(m)}$ at level $q = 0.05$, the largest $k$ satisfying
\begin{equation}
k = \max\left\{\, j \in \{1,\dots,m\} \;:\; p_{(j)} \le \frac{j}{m}\,q \,\right\}
\end{equation}
is found, and all hypotheses $(1),\dots,(k)$ are rejected; if no such $j$ exists,
then no hypothesis is rejected. Features subject to floor effects in the formal
register (expressive punctuation, boosters) were excluded before correction and
so do not contribute to $m$.


\section{Experiment 2: LLM-Swedish Text Detection by AI tools}

For our text detection component, two AI-detection-by-AI tool methods were
compared. We compared two methods to detect the difference between a detection method on the embedding-space level and a method that examines surface-level patterns in a similar manner to a human reader.

\subsection{Classifier two-sample test}
Our first method was a classifier two-sample test \parencite{lopezpaz2017revisiting}
implemented as a logistic regression trained to separate human and LLM
multilingual sentence embeddings using
\texttt{paraphrase-multilingual-mpnet-base-v2} \parencite{reimers-2019-sentence-bert},
embedding model that is explicitly multilingual and covers Swedish.
Logistic regression was chosen over a higher-capacity classifier as the quantity
of interest is whether the two classes are separable at all and not the maximum
achievable separation, which would require a bigger dataset to achieve. Classifiers were scored by 5-fold cross-validated ROC AUC. Significance was established with permutation testing,
where the class labels were shuffled and the cross-validation repeated to build
an empirical null distribution.

\subsection{LLM-as-judge}
The second method was an LLM-as-judge evaluation task using the reasoning models
Claude Opus 4.8, Gemini 3.1 Pro, and DeepSeek-R1 v4 Flash. All three LLMs
classified the same set of 400 texts as human- or AI-written, balanced across
register and condition. Reasoning models were chosen as the task asks for a
justified judgement about textual features rather than a surface classification.
All prompts were written in Swedish, consistent with the generation prompts, and
the exact prompts are described in Table \ref{tab:detection-prompt-conditions}.

\subsection{Generalisation across generators}
All detection methods used the same dataset of texts. To test the
AI-detectability of our Swedish dataset across other generators, we produced two
further AI Swedish datasets using a newer GPT model (GPT-5.6) and a model from a
different family (Mistral Small 3.2). These datasets were created under the same
conditions as the GPT-5.2 data, with identical prompts and human-written inputs,
and were evaluated with both detectability methods described above, to test whether the detectors differ in their judgements across generators. AI detector prompts can be seen in full in the Appendix. 



\section{Experiment 3: AI Text Detection by Swedish speakers}


\begin{figure}[H]
    \centering
    \includegraphics[width=0.9\linewidth]{abstract_ai_inte_example.png}
    \caption{Excerpt of the forum comments section of the human AI-detection quiz.}
    \label{fig:quiz-example}
\end{figure}


For the human detection experiment, a web quiz was created by the author\footnote{accessible at ginawelsh.github.io/ai\_eller\_inte} that contained a subset of the abstract and comment data for participants to classify as human-written or AI-generated. There were 10 forum thread questions with 4 comment replies per question, and 20 thesis abstract questions, a total of 60 questions. The quiz was shared among WhatsApp student groups, contact referrals, as well as a general Swedish-speaking AI interest Facebook group with around 90 thousand members. Before the game started, the user was prompted to fill in information about their age, their occupation, their level of Swedish, and their estimated hours of LLM tool usage per week. Only participants who self-reported a Swedish proficiency of intermediate level and above could access the quiz. The LLM-generated data points were balanced across prompt condition (LLM baseline, generic-evasive, feature-evasive). As this aspect of the study was public-facing and accessible to anyone, care was taken to exclude any topics or comment replies with highly inflammatory language, sexual references or references to violence. Each participant was free to exit the quiz at any time voluntarily. The comments section of the quiz was shown first before the abstracts section in order to retain engagement from users, as the abstracts section had a higher reading burden than the comments section. At the end of the game session (or upon exit), the user was shown their accuracy results and prompted to fill in an optional short reflection to explain their reasoning process for LLM-detection throughout the game.


\begin{table}[H]
\centering
\begin{tabular}{llrr}
\toprule
\textbf{Dimension} & \textbf{Category} & \textbf{n} & \textbf{Answers} \\
\midrule
\multirow{3}{*}{Language level}
  & Native speaker & 31 & 825 \\
  & Advanced       & 3  & 62  \\
  & Intermediate   & 3  & 54  \\
\midrule
\multirow{6}{*}{Age}
  & younger than 18 & 1  & 8   \\
  & 18--24          & 1  & 36  \\
  & 25--34          & 16 & 437 \\
  & 35--44          & 6  & 71  \\
  & 45--60          & 12 & 351 \\
  & 60+             & 1  & 38  \\
\midrule
\multirow{5}{*}{LLM usage per week}
  & Not at all   & 1  & 7   \\
  & 0--5 hours   & 15 & 421 \\
  & 5--10 hours  & 9  & 242 \\
  & 10--20 hours & 5  & 74  \\
  & 20+ hours    & 7  & 197 \\
\midrule
\textbf{Total} & & \textbf{37} & \textbf{941} \\
\bottomrule
\end{tabular}
\caption{Summary of Swedish speaker participants across language level, age, and
weekly LLM usage.}
\label{tab:participant-summary}
\end{table}

Table \ref{tab:participant-summary} summarises the participant pool for the human detection experiment. The results of each user session were stored in a Google Sheet form with information
about their usage time-stamp, question identification number,
which questions were determined as human, and accuracy scores.
Accuracy, human precision (number of documents labelled as human that were genuinely human) and AI recall (proportion of AI generated documents that were correctly labelled as AI by the participant) were measured as performance scores. Voluntary reasoning submissions were collated into a table organised by theme. 

\chapter{Results \& Discussion}

\section{Experiment 1 Results: Feature analysis}
Table \ref{tab:features-gpt} displays the feature effect sizes of each feature analysed, which shows the features of the LLM-Swedish texts that most considerably diverged from their human-Swedish counterparts.

\begin{table}[h!]
  \centering
  \small
  \setlength{\tabcolsep}{3pt}
  \begin{tabular}{lcccccc}
    & \multicolumn{3}{c}{\textbf{Formal}} & \multicolumn{3}{c}{\textbf{Informal}} \\
    \cmidrule(lr){2-4}\cmidrule(lr){5-7}
    \textbf{Feature} & Base & Hum-like & Det-evasive & Base & Hum-like & Det-evasive \\
    \toprule
    \multicolumn{7}{l}{\textbf{Sentence complexity}} \\
    \addlinespace[2pt]
    Word count                   & \cg{3}$+0.10$* & \cg{29}$+1.16$* & \cg{34}$+1.37$* & \cg{34}$+1.35$* & \cg{38}$+1.53$* & \cg{34}$+1.37$* \\
    Sentence count               & \crd{11}$-0.44$* & \cg{9}$+0.36$* & \cg{49}$+1.94$* & \cg{24}$+0.96$* & \cg{29}$+1.17$* & \cg{41}$+1.62$* \\
    Mean parse-tree depth        & \cg{32}$+1.27$* & \cg{31}$+1.23$* & \crd{33}$-1.33$* & \cg{15}$+0.61$* & \cg{15}$+0.60$* & \crd{15}$-0.60$* \\
    Mean dependency distance     & \cg{26}$+1.02$* & \cg{35}$+1.38$* & \crd{29}$-1.17$* & \cg{22}$+0.87$* & \cg{22}$+0.88$* & \crd{8}$-0.30$* \\
    Subordinate-clause rate      & \cg{12}$+0.48$* & \cg{18}$+0.72$* & \crd{22}$-0.89$* & \cg{6}$+0.24$* & \cg{7}$+0.28$* & \crd{15}$-0.60$* \\
    \addlinespace[4pt]
        \midrule
    \multicolumn{7}{l}{\textbf{Lexical diversity}} \\
    \addlinespace[2pt]
    Lexical diversity (MATTR)    & \cg{31}$+1.22$* & \cg{28}$+1.13$* & \crd{11}$-0.45$* & \crd{16}$-0.65$* & \crd{18}$-0.72$* & \crd{14}$-0.58$* \\
    Lexical diversity (MTLD)     & \cg{36}$+1.43$* & \cg{38}$+1.51$* & \crd{21}$-0.83$* & \cg{19}$+0.75$* & \cg{18}$+0.72$* & \cg{20}$+0.80$* \\
    Type-token ratio             & \cg{14}$+0.54$* & \crd{7}$-0.29$* & \crd{42}$-1.67$* & \crd{26}$-1.02$* & \crd{29}$-1.17$* & \crd{22}$-0.89$* \\
    Bigram repetition rate       & \crd{19}$-0.74$* & \crd{10}$-0.40$* & \cg{33}$+1.31$* & \cg{11}$+0.44$* & \cg{14}$+0.55$* & \cg{6}$+0.22$* \\
    Function-word rate           & \crd{22}$-0.87$* & \crd{13}$-0.51$* & \cg{3}$+0.10$ & \crd{4}$-0.15$* & \crd{2}$-0.09$* & \crd{7}$-0.26$* \\
    \addlinespace[4pt]
        \midrule
   \multicolumn{7}{l}{\textbf{Perplexity}} \\
\addlinespace[2pt]
Perplexity                   & \cg{6}$+0.25$* & \cg{6}$+0.22$* & \crd{0}$-0.00$ & \crd{1}$-0.03$* & \crd{1}$-0.03$* & \crd{1}$-0.03$* \\
Burstiness (sent. PPL SD)    & \cg{2}$+0.09$ & \cg{3}$+0.10$ & \cg{10}$+0.38$* & \cg{2}$+0.08$* & \cg{3}$+0.13$* & \cg{2}$+0.09$* \\
    \addlinespace[4pt]
            \midrule
    \multicolumn{7}{l}{\textbf{Pragmatic markers}} \\
    \addlinespace[2pt]
    Hedge rate                   & \cg{16}$+0.64$* & \cg{23}$+0.92$* & \cg{7}$+0.26$* & \cg{3}$+0.13$* & \cg{5}$+0.19$* & \crd{4}$-0.16$* \\
    Negation rate                & \crd{7}$-0.28$* & \cg{2}$+0.08$* & \crd{5}$-0.18$ & \crd{4}$-0.14$ & \crd{4}$-0.14$ & \cg{3}$+0.11$* \\
    \addlinespace[4pt]
        \midrule
    \multicolumn{7}{l}{\textbf{Punctuation}} \\
    \addlinespace[2pt]
    Comma rate                   & \cg{42}$+1.67$* & \cg{39}$+1.55$* & \crd{7}$-0.27$* & \cg{8}$+0.32$* & \cg{7}$+0.29$* & \crd{6}$-0.23$* \\
    Full-stop rate               & \crd{22}$-0.87$* & \crd{24}$-0.97$* & \cg{58}$+2.32$* & \crd{8}$-0.30$* & \crd{8}$-0.30$* & \cg{5}$+0.20$* \\
    Dash rate                    & \cg{13}$+0.53$* & \cg{18}$+0.71$* & \crd{4}$-0.16$ & \cg{22}$+0.86$* & \cg{19}$+0.77$* & \cg{4}$+0.17$* \\
    \addlinespace[4pt]
    \midrule
    \multicolumn{7}{l}{\textbf{Named entity recognition}} \\
    \addlinespace[2pt]
    Named-entity rate (all)      & \crd{13}$-0.53$* & \crd{14}$-0.56$* & \crd{1}$-0.04$ & \crd{4}$-0.15$* & \crd{4}$-0.14$* & \cg{1}$+0.05$* \\
    \quad location (LOC)         & \crd{8}$-0.30$* & \crd{8}$-0.33$* & \crd{3}$-0.10$ & \crd{2}$-0.09$ & \crd{2}$-0.07$ & \cg{1}$+0.02$* \\
    \quad person (PRS)           & \crd{9}$-0.35$* & \crd{8}$-0.33$* & \cg{1}$+0.04$ & \crd{3}$-0.11$* & \crd{3}$-0.10$* & \crd{0}$-0.01$* \\
    \quad organisation (ORG)     & \crd{8}$-0.33$* & \crd{9}$-0.37$* & \crd{1}$-0.04$ & \crd{1}$-0.04$ & \crd{1}$-0.04$* & \cg{4}$+0.15$* \\
        \midrule
            \multicolumn{7}{l}{\textbf{Discourse particles}} \\
    \addlinespace[2pt]
    Informal-hedge count         & -- & -- & -- & \cg{29}$+1.15$* & \cg{36}$+1.42$* & \cg{18}$+0.73$* \\
    Booster count                & -- & -- & -- & \cg{17}$+0.68$* & \cg{19}$+0.75$* & \crd{3}$-0.13$* \\
    Discourse-particle rate      & -- & -- & -- & \crd{2}$-0.09$ & \crd{3}$-0.10$ & \crd{9}$-0.36$* \\
    \addlinespace[4pt]
        \midrule
    \multicolumn{7}{l}{\textbf{Sentiment analysis}} \\
    \addlinespace[2pt]
    Signed polarity              & --- & --- & --- & \cg{1}$+0.02$* & \cg{1}$+0.04$* & \crd{6}$-0.23$* \\
    Affective extremity          & --- & --- & --- & \crd{4}$-0.15$* & \crd{4}$-0.16$* & \cg{3}$+0.11$* \\
    Subjective rate              & --- & --- & --- & \crd{3}$-0.12$* & \crd{3}$-0.11$* & \crd{0}$+0.00$ \\
    \addlinespace[4pt]
        \midrule
    \multicolumn{7}{l}{\textbf{Profanity (swear-word intensity tiers)}} \\
    \addlinespace[2pt]
    Mild-swear count              & -- & -- & -- & \cg{4}$+0.15$* & \cg{5}$+0.19$* & \cg{7}$+0.29$* \\
    Moderate-swear count          & -- & -- & -- & \crd{1}$-0.02$ & \cg{2}$+0.07$* & \crd{5}$-0.18$* \\
    Strong-swear count            & -- & -- & -- & \crd{2}$-0.08$* & \crd{2}$-0.09$* & \crd{2}$-0.06$* \\
    \bottomrule
  \end{tabular}
  \caption{Selected paired feature differences between human-written and GPT-5.2-generated
  Swedish text (Cohen's $d_z$, LLM-human; paired Wilcoxon signed-rank with
  Benjamini-Hochberg FDR correction). Positive (green) values indicate higher
  feature values in LLM text; negative (red) values indicate higher values in
  human text; shading intensity is proportional to $|d_z|$. * denotes
  significance at $q<0.05$.}
  \label{tab:features-gpt}
\end{table}

\subsection{Sentence complexity}

 
\begin{table}[H]
  \centering
  \footnotesize
  \renewcommand{\arraystretch}{1.3}
  \begin{tabularx}{\linewidth}{@{}>{\bfseries\raggedright}p{2.2cm} >{\raggedright\arraybackslash}X@{}}
    \toprule
    \multicolumn{2}{@{}p{0.95\linewidth}}{\textbf{Thread question}: Ifall det finns no go zoner där polisen inte färdas varför säljer man bara inte sitt knark där?
      \newline{\itshape If there are no-go zones where police don't go, why don't people just sell their drugs there?}} \\
    \midrule
    Human & För att det blir \hlg{chok me knas} om du becknar i de områden där folk redan etablerat sin marknad, \hlg{broder}.
      \newline{\itshape Because you'll get \hlg{shanked with a knife} if you deal in areas where people have already established their market, \hlg{bro}.} \\
    Baseline & ...det finns \hlr{span, civvare, kameror, avlyssning, tillslag} och folk som åker dit ändå.
      \newline{\itshape ...there's \hlr{surveillance, plainclothes cops, cameras, wiretaps, raids}, and people who go there anyway.} \\
    generic-evasive & ...du har \hlr{span, kameror, hemliga tvångsmedel, civvare, tillslag}, och sen är det ofta just där de lägger extra tryck.
      \newline{\itshape ...you've got \hlr{surveillance, cameras, covert measures, plainclothes cops, raids}, and that's often exactly where they apply extra pressure.} \\
    feature-evasive & \hlr{Kameror.} \hlr{Avlyssning.} \hlr{Civilare.} \hlr{Anonyma tips.}
      \newline{\itshape \hlr{Cameras.} \hlr{Wiretaps.} \hlr{Plainclothes.} \hlr{Anonymous tips.}} \\
    \bottomrule
  \end{tabularx}
  \caption{Sentence complexity across matched conditions, informal register.}
  \label{tab:example-complexity-informal}
\end{table}

Sentence complexity increased considerably across almost all LLM prompt conditions, a finding that mirrors both \textcite{guo2023closechatgpthumanexperts} and \textcite{su2024robustdetectionllmgeneratedtext}, where English LLM texts were longer across several domains.  In both registers, the baseline and generic-evasive prompts raised complexity above human text on mean dependency distance, parse-tree depth, and subordinate-clause rate, a finding also seen in \textcite{guo2023closechatgpthumanexperts}, \textcite{su2024robustdetectionllmgeneratedtext} and \textcite{munoz-ortiz2024contrasting}.
 
 
 The feature-evasive prompt reversed all three measures in its output texts and dramatically over-corrected them to the most negative feature effect sizes, particularly for the formal register. The lowered complexity led to shorter sentences, which in turn increased the full-stop rate from -0.87 at baseline to +2.32 under evasion, the highest AI signal across all measurements. In general, the evasive prompt did not neutralise the complexity signal, rather it inverted the signal and created an AI 'tell' in the other direction. It appeared that the feature-evasive prompt instruction `\textit{write predominantly short sentences}' had an unexpectedly drastic influence on its output, suggesting that for our Swedish data, certain evasive instructions may be overly prioritised over others. Other evasive prompt designs in previous English work, such as that of \textcite{lu2024sico} that used an evasive prompt technique focusing heavily perplexity and burstiness, did not appear to have as strong an over-corrective effect as in our study. It is unclear whether the over-correction is due to the Swedish language itself being under-represented in the GPT-5.2 training data or as a model-specific error that could hold across multiple languages.

\subsection{Lexical diversity} 

\begin{table}[H]
  \centering
  \footnotesize
  \renewcommand{\arraystretch}{1.3}
  \begin{tabularx}{\linewidth}{@{}>{\bfseries\raggedright}p{2.2cm} >{\raggedright\arraybackslash}X@{}}
    \toprule
    \multicolumn{2}{@{}p{0.95\linewidth}}{\textbf{Thesis title}: Expertkunskap och metakognition: experters förmåga till metakognitiv övervakning
      \newline{\itshape Expert knowledge and metacognition: experts' capacity for metacognitive monitoring}} \\
    \midrule
    LLM Baseline & ...för att strukturera förståelsen av \hlg{metakognitiv noggrannhet}... exempelvis genom att koppla samman faktisk prestation med skattad säkerhet och därigenom fånga graden av \hlg{kalibrering}.
      \newline{\itshape ...to structure the understanding of \hlg{metacognitive accuracy}... for example by linking actual performance to estimated confidence and thereby capturing the degree of \hlg{calibration}.} \\
    LLM feature-evasive & Den tar upp experters förmåga till \hlr{metakognitiv övervakning}. \hlr{Metakognitiv övervakning} blir då en del av metakognition... Jag kopplar detta till \hlr{metakognitiv övervakning}... (recurs \textbf{9} times verbatim in one abstract)
      \newline{\itshape It addresses experts' capacity for \hlr{metacognitive monitoring}. \hlr{Metacognitive monitoring} then becomes a part of metacognition... I link this to \hlr{metacognitive monitoring}...} \\
    \bottomrule
  \end{tabularx}
  \caption{Lexical variation vs.\ verbatim repetition of the same key phrase. Green marks synonym substitution; red marks unvaried repetition of the identical two-word phrase.}
  \label{tab:example-diversity}
\end{table}

In the formal register, there was a divide between the output of the baseline and generic evasive LLM prompts and the feature-evasive prompt. The baseline and generic-evasive prompt produced considerably lexically more diverse LLM-Swedish than their human counterparts, particularly for MATTR and MTLD (range +1.13 - +1.51), and lower bigram repetitions and function word rates. The higher results here are in keeping with \textcite{Reviriego_2024}, where LLM-generated formal register texts were found to be more lexically diverse than humans'. In a similar way to sentence complexity, the feature-evasive prompt output over-corrected the lexical diversity by decreasing MATTR, MTLD, and type-token ratio to considerably lower than the human texts, as well as drastically increasing the bigram repetition and function word rates. The feature-evasive prompt instruction `\textit{reuse the same key terms and phrases verbatim rather than varying with synonyms}' was possibly too strong an instruction, leading to an over-correction that again raised the AI signal into the opposite direction. Again, this points to a sensitivity of the model (GPT-5.2) to place too much weight into a particular instruction when reading an evasive prompt informed by features. 


For the informal register, lexical diversity was less definitive. While all prompt conditions had measured output in the same direction (lower MATTR, higher MTLD, lower type-token ratio, higher bigram repetition, lower function word rate, this combination of findings is slightly contradictory. The MATTR and type-token ratio being considerably lower than the human control, while the MTLD being higher, suggests that overall, the vocabulary size was not as large among the LLM-Swedish forum comments, but vocabulary density may have clustered into peaks locally within the texts, rather than stretched out over texts as a whole. The higher MLTD could potentially be attributed to the prompt instruction ``Write with dense content and fewer function words'', where the possible outcome was more lexically dense words being put together locally in text. However, our MATTR measurement for the informal data comes with a caveat: the MATTR was calculated over a 50-token fixed sliding window. Where a comment was below 50 tokens, the calculation fell to plain type-token ratio. 807 human comments (approximately 70\%) fell below this 50-token window, and therefore the TTR measurement majorly influenced the MATTR calculation. In contrast, only 2.7\% to 5.31\% (by prompt condition) of LLM comments triggered the TTR fallback. This plausibly explains the mixed result between MATTR and MTLD for the informal register. Therefore, we cannot confirm that the mixed score indicates differences in content density for the informal register.  The feature-evasive prompt for the informal text mitigated the AI signal closer to the human benchmark in MATTR, type-token ratio and bigram repetition rates; in other measures, the direction was reversed slightly but without over-correction. This shows that the feature-evasive prompt did not over-correct in all instances.




\subsection{Punctuation} Formal punctuation was correlated with sentence complexity, with the high-comma, low-full-stop signals of the baseline and generic-evasive outputs inverted under the evasive prompt due to shortened sentences. This was a clear example of a single instruction affecting two features at once, in this case reducing one signal while inadvertently amplifying another. Informal punctuation was largely non-distinctive, the exception being dashes, prominent under the baseline and generic-evasive prompts (+0.86, +0.77) and suppressed by the `do not use dashes' instruction (+0.17) in the feature-evasive prompt.


\subsection{Discourse particles} 

\begin{table}[H]
  \centering
  \footnotesize
  \renewcommand{\arraystretch}{1.3}
  \begin{tabularx}{\linewidth}{@{}>{\bfseries\raggedright}p{2.2cm} >{\raggedright\arraybackslash}X@{}}
    \toprule
    \multicolumn{2}{@{}p{0.95\linewidth}}{\textbf{Thread question}: Vilka enskilda politiker ger bäst intryck? Jag har just nu svårt för alla våra partiledare. Ygeman känns stabil. Vilka favoriter har ni?
      \newline{\itshape Which individual politicians make the best impression? I'm currently struggling with all our party leaders. Ygeman feels solid. What are your favorites?}} \\
    \midrule
    generic-evasive & Sen finns det en del ``bänkpolitiker'' som gör bättre intryck, \hlg{typ} folk som faktiskt svarar på frågor utan att läsa innantill.
      \newline{\itshape Then there are some ``bench politicians'' who make a better impression, \hlg{like} people who actually answer questions instead of reading from a script.} \\
    feature-evasive & \hlr{Fattar känslan.} Partiledarna känns mest PR-maskiner nu.
      \newline{\itshape \hlr{Get the feeling.} The party leaders mostly feel like PR machines now.} \\
    \bottomrule
  \end{tabularx}
  \caption{Discourse particles across matched conditions. Green marks the colloquial particle \textit{typ} (\textit{like}); red marks the evasive passage's complete absence of a discourse particle.}
  \label{tab:example-particles}
\end{table}

\begin{table}[H]
  \centering
  \footnotesize
  \renewcommand{\arraystretch}{1.3}
  \begin{tabularx}{\linewidth}{@{}>{\bfseries\raggedright}p{2.2cm} >{\raggedright\arraybackslash}X@{}}
    \toprule
    \multicolumn{2}{@{}p{0.95\linewidth}}{\textbf{Thread question}: Vilka enskilda politiker ger bäst intryck? Jag har just nu svårt för alla våra partiledare. Ygeman känns stabil. Vilka favoriter har ni?
      \newline{\itshape Which individual politicians make the best impression? I'm currently struggling with all our party leaders. Ygeman feels solid. What are your favorites?}} \\
    \midrule
    LLM Baseline & Sen är det \hlg{ju} lite det där: de som ger bäst intryck är ofta de som inte syns mest.
      \newline{\itshape Then there's, \hlg{you know}, that thing: those who make the best impression are often the ones who are least visible. \emph{(``ju'' is an untranslatable modal particle, roughly ``as you'd expect / after all''.)}} \\
    LLM generic-evasive & Men några som ändå ger rätt stabilt intryck är \hlg{typ} Ygeman som du säger.
      \newline{\itshape But some who still give a fairly solid impression are, \hlg{like}, Ygeman as you say.} \\
    LLM feature-evasive & \hlr{Ygeman håller jag med om.}
      \newline{\itshape \hlr{Ygeman, I agree.} \emph{(lack of hedging via fronted object)}} \\
    \bottomrule
  \end{tabularx}
  \caption{Hedging across matched conditions for one example.}
  \label{tab:example-hedging}
\end{table}


Discourse particles were only measured in the informal comment data. Hedging patterns created an AI signal for the LLM baseline and generic-evasive prompts, particularly for the generic-evasive prompt's usage of informal hedge counts. The feature-evasive prompt mitigated this in the instruction ``\textit{do not overdo hedging words, especially in colloquial ones}''\ldots{} However, there was still a signal in the feature-evasive prompt that meaningfully deviated the comments from human control in a higher amount of hedges. Regardless, the feature-evasive prompt in general successfully mitigated hedging pattern AI signals compared to the LLM baseline and generic-evasive prompts. Baseline and generic-evasive prompts over-produced hedges in both registers (formal +0.64, +0.92; a smaller but significant signal informally). Unlike the complexity and diversity features, the evasive prompt corrected hedging back toward human values with only slight over-correction, under the instruction to ``\textit{keep hedging and epistemic expressions to a minimum and [not] overcompensate by adding colloquial hedges}''.


\subsection{Perplexity and Burstiness}

Text perplexity and burstiness appeared to affect the formal dataset more than the informal dataset. The LLM baseline and generic-evasive prompts produced abstracts that were slightly higher in perplexity than the human abstracts. However, the feature-evasive prompt in fact successfully mitigated this higher perplexity to be null in feature effect size (-0.00). 

\subsection{Named entity recognition} In the formal register, baseline and generic-evasive outputs carried fewer named entities than human text, reducing specificity. The evasive prompt raised NER across all categories back toward human levels, following be concrete and mention specific names where you can. As with hedging, this feature was corrected without over-shooting. Informal NER effects were small throughout, the evasive prompt leaving only a slight deviation from human text.

\begin{table}[H]
  \centering
  \footnotesize
  \renewcommand{\arraystretch}{1.3}
  \begin{tabularx}{\linewidth}{@{}>{\bfseries\raggedright}p{2.2cm} >{\raggedright\arraybackslash}X@{}}
    \toprule
    \multicolumn{2}{@{}p{0.95\linewidth}}{\textbf{Thesis title}: Expertkunskap och metakognition: experters förmåga till metakognitiv övervakning
      \newline{\itshape Expert knowledge and metacognition: experts' capacity for metacognitive monitoring}} \\
    \midrule
    Human & Förmågan till metakognitiv övervakning mättes med hjälp av mätteknikerna Assessment of Cognitive Monitoring Effectiveness \hlg{(ACME)} och Metacognitive Knowledge Monitoring Assessment \hlg{(KMA)}.
      \newline{\itshape The capacity for metacognitive monitoring was measured using the instruments Assessment of Cognitive Monitoring Effectiveness \hlg{(ACME)} and Metacognitive Knowledge Monitoring Assessment \hlg{(KMA)}} \\
    Baseline & Som analytiska verktyg används \hlg{ACME} och \hlg{KMA} för att strukturera förståelsen av metakognitiv noggrannhet...
      \newline{\itshape \hlg{ACME} and \hlg{KMA} are used as analytical tools to structure the understanding of metacognitive accuracy...} \\
    feature-evasive & Jag utgår från John \hlr{Flavell (1979)}. Jag tar också stöd i Ann \hlr{Brown (1987)}... Jag använder Anders \hlr{Ericsson (1993)} och hans arbete om avsiktlig träning.
      \newline{\itshape I draw on John \hlr{Flavell (1979)}. I also draw support from Ann \hlr{Brown (1987)}... I use Anders \hlr{Ericsson (1993)} and his work on deliberate practice.} \\
    \bottomrule
  \end{tabularx}
  \caption{\footnotesize Named entities across matched conditions. Green marks the technical instrument acronyms (ACME/KMA), consistent in every condition. Red marks named researchers cited only by the feature-evasive output -- real scholars invoked as fabricated citations, since they appear in none of the other three conditions.}
  \label{tab:example-ner}
\end{table}


\subsection{Sentiment analysis}

There were more modest (albeit significant) differences between human- and LLM-Swedish. The LLM baseline and generic-evasive prompt outputs had a very slightly higher polarity (+0.02* and +0.04*, respectively), suggesting slightly more positive sentiment, and a lower affective extremity and subjective rate, pointing to more emotionally neutral comments. These findings corroborate previous studies that found LLM text to be more neutral than human text \cite{guo2023closechatgpthumanexperts}. However, the feature-evasive output had considerably lower polarity than the other texts with an effect size of -0.23*, and a higher affective extremity (+0.11), suggesting that the feature-evasive LLM prompt over-corrected texts to lean more negative and more emotional in sentiment. The feature-evasive prompt did not explicitly instruct the LLM to produce texts with more negative sentiment; however, it did instruct to `\textit{feel free to use negated clauses...where it fits}', which may have had an effect on sentiment. In addition, we measured the amount and level of profanity within our informal texts. We found that LLM texts in fact generated a higher level of mild-level profanity (+0.15 - +0.29) but overall a lower level of moderate and strong-level profanity. This suggests that GPT-5.2 was capable of producing profanities in its output, but the profanities were almost always milder and strong profanities were rare. Overall, the LLM baseline and generic-evasive output produced comments with sentiment properties that aligned with previous works, but the feature-evasive comments over-corrected this measure.

\subsection{Connection to related work}
Our study's most salient and consistent AI signals corroborate findings reported for other languages and domains and therefore suggests that Swedish is not an exception to general LLM text generation tendencies. The lower NER rates in both the LLM baseline and generic-evasive prompt output mirrors \textcite{liao2023differentiating}, which found ChatGPT-generated English medical text to rely on more generic terminology and reasoning as opposed to specific, case-tailored information; \textcite{guo2023closechatgpthumanexperts} also characterised ChatGPT output as logical and detailed but prone to generality. Our feature-evasive prompt corrected this signal with subtlety without over-shooting (due to the \textit{be concrete and name specific names} instruction) indicating that this was a stable and directly addressable AI signal in our datasets. Additionally, the LLM Swedish was consistently longer and more sentence-heavy in both registers, in line with the `long and detailed' LLM text profile of \textcite{guo2023closechatgpthumanexperts}. The impersonal, templated character captured by our LLM judges (Table~\ref{tab:judge-rationale-cues}) aligns with \textcite{tang2023sciencedetectingllmgeneratedtexts}, who found LLM text to be more formal and structured. Punctuation also emerged as a distinctive feature, in keeping with Korean study \textcite{park2025katfishnet}. However, we caution against these signal patterns being read as language-independent, as some signals shifted strongly between registers and prompt conditions, indicating that the direction of the signals compared to the human baseline were not consistent.

\section{Experiment 2 Results: AI detection by AI and feature-evasive prompts performance}

\subsection{Embedding-based Classifier AI Detection}

Table~\ref{tab:embedding-auc} sets out embedding classifier results for all domains and prompt conditions. Overall, the informal comments were considerably more separable from human text than the formal abstracts (informal 0.886-0.935 vs. formal 0.643-0.813). This register gap narrowed or reversed for the other two generators. For GPT-5.2, the most detectable condition was the informal feature-evasive data, at a CV ROC AUC of 0.935, a result that runs contrary to the expectation of a detector evasive prompt since evasion raised rather than lowered separability relative to the baseline (0.886) and generic-evasive (0.898) conditions. However, this pattern was specific to GPT-5.2; for Mistral, the evasive prompt was the least detectable informal condition (0.874 vs 0.917 baseline and 0.954 generic-evasive). Across generators, Mistral Small 3.2 was the most separable overall (0.874-0.954) while GPT-5.6 evaded the classifier to the greatest extent (0.728-0.823), suggesting that the later GPT version is less detectable to the embedding-based method.

% GENERATED by src/3_ai_detection_scripts/make_embedding_auc_table_matched.py
% -- do not hand-edit. Re-run that script to refresh.
% Source of truth: 2_text_analysis_scripts/csv_files/embedding_auc_matched_gpt56subset.csv
% All cells computed on the SAME documents (the GPT-5.6 quiz100 subset), so
% generators are comparable; the full-corpus figures are NOT, because a
% cross-validated AUC falls as n falls. See EMBEDDING_AUC_CAVEATS.md.
% Bold = most detectable prompt condition within each generator x register block.
\begin{table}[h!]
  \centering
  \small
  \begin{tabular}{lllccc}
    \toprule
    Generator & Register & Prompt & CV ROC AUC & 95\% CI & Perm.\ $p$ \\
    \midrule
    \multirow{6}{*}{GPT-5.2}
      & \multirow{3}{*}{Formal}
        & Baseline         & \cg{14}0.643 & [0.567, 0.719] & 0.003 \\
      & & generic-evasive  & \cg{14}0.643 & [0.567, 0.719] & 0.005 \\
      & & feature-evasive & \cg{31}\textbf{0.813} & [0.753, 0.873] & $<$0.001 \\
    \cmidrule(l){2-6}
      & \multirow{3}{*}{Informal}
        & Baseline         & \cg{39}0.886 & [0.839, 0.933] & $<$0.001 \\
      & & generic-evasive  & \cg{40}0.898 & [0.853, 0.943] & $<$0.001 \\
      & & feature-evasive & \cg{44}\textbf{0.935} & [0.899, 0.971] & $<$0.001 \\
    \midrule
    \multirow{6}{*}{GPT-5.6}
      & \multirow{3}{*}{Formal}
        & Baseline         & \cg{29}\textbf{0.792} & [0.730, 0.855] & $<$0.001 \\
      & & generic-evasive  & \cg{23}0.728 & [0.658, 0.797] & $<$0.001 \\
      & & feature-evasive & \cg{25}0.747 & [0.679, 0.816] & $<$0.001 \\
    \cmidrule(l){2-6}
      & \multirow{3}{*}{Informal}
        & Baseline         & \cg{29}0.790 & [0.727, 0.852] & $<$0.001 \\
      & & generic-evasive  & \cg{28}0.781 & [0.717, 0.845] & $<$0.001 \\
      & & feature-evasive & \cg{32}\textbf{0.823} & [0.765, 0.881] & $<$0.001 \\
    \midrule
    \multirow{6}{*}{Mistral 3.2}
      & \multirow{3}{*}{Formal}
        & Baseline         & \cg{44}\textbf{0.935} & [0.899, 0.971] & $<$0.001 \\
      & & generic-evasive  & \cg{43}0.925 & [0.887, 0.964] & $<$0.001 \\
      & & feature-evasive & \cg{41}0.909 & [0.867, 0.951] & $<$0.001 \\
    \cmidrule(l){2-6}
      & \multirow{3}{*}{Informal}
        & Baseline         & \cg{42}0.917 & [0.876, 0.957] & $<$0.001 \\
      & & generic-evasive  & \cg{45}\textbf{0.954} & [0.923, 0.984] & $<$0.001 \\
      & & feature-evasive & \cg{37}0.874 & [0.824, 0.924] & $<$0.001 \\
    \bottomrule
  \end{tabular}
  \caption{\footnotesize Embedding-based classifier detectability for all LLM generators,
  registers and prompt conditions. Cross-validated ROC AUC of a
  logistic-regression two-sample test on multilingual sentence embeddings
  (\texttt{paraphrase-multilingual-mpnet-base-v2}, 5-fold CV). AUC${}=0.5$ is
  chance; higher means more separable from human text. Cell shading follows
  Table~\ref{tab:llm-judge-detection-by-condition}: green marks values above
  chance, with intensity proportional to the distance from 0.5. Confidence
  intervals computed by Hanley--McNeil normal approximations.}
  \label{tab:embedding-auc}
\end{table}

\subsection{Experiment 2 results: LLM-as-judge detection}

\begin{table*}[t]
  \centering
  \footnotesize
  \setlength{\tabcolsep}{4pt}
  \begin{tabularx}{\textwidth}{lll YYY YYY YYY}
    \toprule
    & & & \multicolumn{3}{c}{AI recall (\%)}
        & \multicolumn{3}{c}{Human precision (\%)}
        & \multicolumn{3}{c}{Accuracy (\%)} \\
    \cmidrule(lr){4-6} \cmidrule(lr){7-9} \cmidrule(l){10-12}
    Gen. & Register & Condition
      & Cl & Gm & DS & Cl & Gm & DS & Cl & Gm & DS \\
    \midrule
    \multirow{6}{*}{GPT-5.2}
      & \multirow{3}{*}{Formal}
        & Baseline         & \cg{50}100 & \cg{46}96  & \crd{3}47 & \cg{50}100 & \cg{46}96 & \cg{7}57  & \cg{50}100 & \cg{48}98 & \cg{8}58 \\
      & & Human-mim.  & \cg{50}100 & \cg{50}100 & \crd{15}35 & \cg{50}100 & \cg{50}100 & \cg{1}51 & \cg{50}100 & \cg{50}100 & \cg{2}52 \\
      & & Det.-evasive & \cg{50}100 & \cg{50}100 & \cg{42}92 & \cg{50}100 & \cg{50}100 & \cg{40}90 & \cg{50}100 & \cg{50}100 & \cg{31}81 \\
    \cmidrule(l){2-12}
      & \multirow{3}{*}{Informal}
        & Baseline         & \cg{19}69 & \crd{40}10 & \crd{45}5 & \cg{26}76 & \cg{2}52 & \crd{3}47 & \cg{35}85 & \cg{4}54 & \crd{4}46 \\
      & & Human-mim.  & \cg{23}73 & \crd{14}36 & \crd{49}1 & \cg{29}79 & \cg{10}60 & \crd{3}47 & \cg{37}87 & \cg{17}67 & \crd{6}44 \\
      & & Det.-evasive & \crd{12}38 & \crd{42}8  & \crd{50}0 & \cg{12}62 & \cg{1}51 & \crd{4}46 & \cg{19}69 & \cg{3}53 & \crd{7}43 \\
    \midrule
    \multirow{6}{*}{GPT-5.6}
      & \multirow{3}{*}{Formal}
        & Baseline         & \cg{49}99 & -- & \cg{41}91 & \cg{49}99 & -- & \cg{40}90 & \cg{49}99 & -- & \cg{37}87 \\
      & & Human-mim.  & \cg{50}100 & -- & \cg{44}94 & \cg{50}100 & -- & \cg{44}94 & \cg{50}100 & -- & \cg{38}88 \\
      & & Det.-evasive & \cg{48}98 & -- & \cg{43}93 & \cg{48}98 & -- & \cg{42}92 & \cg{49}99 & -- & \cg{37}87 \\
    \cmidrule(l){2-12}
      & \multirow{3}{*}{Informal}
        & Baseline         & \cg{14}64 & \crd{25}25 & \crd{10}40 & \cg{24}74 & \cg{6}56 & \cg{9}59 & \cg{32}82 & \cg{11}61 & \cg{13}63 \\
      & & Human-mim.  & \cg{2}52 & \crd{44}6  & \crd{40}10 & \cg{18}68 & \cg{1}51 & \crd{1}49 & \cg{26}76 & \cg{2}52 & \crd{2}48 \\
      & & Det.-evasive & \crd{2}48 & \crd{48}2 & \crd{40}10 & \cg{16}66 & \cg{0}50 & \crd{1}49 & \cg{24}74 & \cg{0}50 & \crd{2}48 \\
    \midrule
    \multirow{6}{*}{\shortstack[l]{Mistral\\3.2}}
      & \multirow{3}{*}{Formal}
        & Baseline         & \cg{50}100 & -- & \cg{48}98 & \cg{50}100 & -- & \cg{48}98 & \cg{50}100 & -- & \cg{40}90 \\
      & & Human-mim.  & \cg{50}100 & -- & \cg{46}96 & \cg{50}100 & -- & \cg{48}98 & \cg{50}100 & -- & \cg{39}89 \\
      & & Det.-evasive & \cg{50}100 & -- & \cg{50}100 & \cg{50}100 & -- & \cg{50}100 & \cg{50}100 & -- & \cg{41}91 \\
    \cmidrule(l){2-12}
      & \multirow{3}{*}{Informal}
        & Baseline         & \cg{34}84 & \cg{7}57 & \crd{26}24 & \cg{36}86 & \cg{19}69 & \cg{3}53 & \cg{42}92 & \cg{27}77 & \cg{5}55 \\
      & & Human-mim.  & \cg{46}96 & \cg{21}71 & \crd{25}25 & \cg{46}96 & \cg{27}77 & \cg{3}53 & \cg{48}98 & \cg{34}84 & \cg{6}56 \\
      & & Det.-evasive & \cg{28}78 & \cg{2}52 & \crd{38}12 & \cg{32}82 & \cg{17}67 & \crd{1}49 & \cg{39}89 & \cg{25}75 & \crd{1}49 \\
    \bottomrule
  \end{tabularx}
  \caption{\footnotesize LLM-as-judge results by generator, register and prompt condition.
  Judges: Cl = Claude, DS = DeepSeek, Gm = Gemini. AI recall is the share of
  LLM texts the judge correctly flagged as AI. Human precision is the share of
  the items the judge labelled human that really are human, computed by pooling
  that condition's generated items with the register's shared human pool.
  Accuracy is the share of that same pooled set (the condition's generated items
  plus the register's human items) the judge classified correctly overall. Gemini formal-register results for GPT-5.6 and Mistral were unavailable at submission, so their accuracy is marked --.}
  \label{tab:llm-judge-detection-by-condition}
\end{table*}

\subsubsection{}{LLM judge differences} Table~\ref{tab:llm-judge-detection-by-condition}  outlines the LLM-as-judge results. Overall, Claude demonstrated the highest overall performance in detecting LLM-Swedish, displaying the highest accuracy, AI recall and human precision across all generator models and prompt conditions and hitting ceiling performance (100\%) in several conditions. Gemini scored comparably to Claude for the GPT-5.2 and Mistral 3.2 formal texts; however, scores were much lower for the informal domain, particularly for the feature-evasive prompt condition. DeepSeek-R1 demonstrated a mixed performance, scoring highly in some conditions (GPT-5.6 Formal, Mistral 3.2 Formal) but very low in others, especially in the informal domains; for example, it scored  0\% recall on the  GPT-5.2 informal feature-evasive condition, indicating an excessive over-labelling of texts as human-authored. Overall, the informal domain was more difficult across all LLM judges, with the feature-evasive condition the most difficult across generators and judges, indicating strong evasion. Our GPT-5.6 generated texts did not prove to be considerably harder to detect for our LLM judges than our GPT-5.2 pool, with some models actually performing better on certain GPT 5.6 conditions than their GPT 5.2 equivalents; for example, DeepSeek-R1's human precision ranged 51\%-90\% for GPT 5.2 formal texts and 90-94\% for GPT 5.6 formal. For Mistral 3.2 generated texts, the embedding classifier found both registers highly separable (formal AUC 0.909-0.935; informal 0.874-0.954), whereas the LLM judges detected the formal texts reliably but performed markedly worse on informal comments. The discrepancy in performance indicates a dissociation in the embedding-based classifier and LLM-as-judge methods, with the LLM-as-judge being fooled by surface-level patterns that the embedding classifier avoids. The underperformance of the DeepSeek-R1 reasoning model could be attributed to several factors. However, many of the possible factors, such as parameter size, training data language composition, and architecture type are unable to be confirmed due to the proprietary, closed-source nature of these models, the only exception being that DeepSeek has disclosed its R1 reasoning model to contain 671 billion parameters and prioritises English and Chinese language data in its training \parencite{Guo_2025}. The predominance of English and Chinese data could therefore be a possible factor in its underperformance on classifying LLM-Swedish data.

% GENERATED by src/3_ai_detection_scripts/make_judge_cue_table.py -- do not hand-edit.
% Supersedes the hand-authored judge_rationale_cues_table.tex (whose keyword
% lists were never recorded and whose numbers do not reproduce).
\begin{table}[t]
  \centering
  \footnotesize
  \setlength{\tabcolsep}{3.5pt}
  \begin{tabular}{@{}l cc cc cc@{}}
  \toprule
    & \multicolumn{2}{c}{Claude} & \multicolumn{2}{c}{Gemini} & \multicolumn{2}{c}{DeepSeek-R1} \\
    \cmidrule(lr){2-3} \cmidrule(lr){4-5} \cmidrule(l){6-7}
    Cue invoked in rationale & F & I & F & I & F & I \\
    \midrule
    Template structure       & \cg{48}$96$ & \cg{39}$79$ & \cg{45}$90$ & \cg{22}$45$ & \cg{48}$96$ & \cg{14}$29$ \\
    Concrete specifics       & \cg{36}$72$ & \cg{38}$76$ & \cg{24}$48$ & \cg{20}$40$ & \cg{20}$41$ & \cg{38}$76$ \\
    Impersonal tone          & \cg{16}$32$ & \cg{33}$67$ & \cg{10}$21$ & \cg{12}$24$ & \cg{13}$26$ & \cg{43}$87$ \\
    Emotion / subjectivity   & \cg{3}$6$ & \cg{24}$49$ & \cg{1}$2$ & \cg{13}$27$ & \cg{5}$11$ & \cg{22}$45$ \\
    Colloquial register      & \cg{2}$4$ & \cg{37}$75$ & \cg{4}$8$ & \cg{40}$80$ & \cg{4}$9$ & \cg{48}$97$ \\
    Variation / repetition   & \cg{23}$46$ & \cg{2}$5$ & \cg{21}$43$ & \cg{3}$7$ & \cg{46}$92$ & \cg{7}$15$ \\
    Errorless polish         & \cg{21}$42$ & \cg{24}$49$ & \cg{19}$38$ & \cg{10}$21$ & \cg{5}$11$ & \cg{35}$71$ \\
    \bottomrule
  \end{tabular}
  \caption{Share (\%) of each judge's rationales invoking a given cue, by
  register (F~=~formal, I~=~informal). Coded by keyword matching on the
  Swedish \texttt{reasoning} field.}
  \label{tab:judge-rationale-cues}
\end{table}

\subsubsection{LLM reasoning patterns} As each LLM-as-judge call delivered a reasoning string with each classification, we organised the reasoning rationales of each model (see Table~\ref{tab:judge-rationale-cues}). All three judges cited template-like structure and lexical variation in the formal register,  and colloquiality, emotion and impersonal tone for the informal, indicating that different reasoning methods were employed for formal versus informal data. DeepSeek-R1's weighted repetition much more heavily in the formal register and impersonal tone and errorless polish in the informal one. The assumption that AI-generated formal register texts have more repetition may have helped DeepSeek-R1's score in the formal feature-evasive conditions (since our feature-analysis showed lower lexical diversity and higher bigram repetition rates); however, it may have contributed to its lower scores on the baseline and generic-evasive conditions, which had higher lexical diversity and lower bigram rates. Gemini also scored low on the GPT-5.2 and GPT-5.6 informal domains; in its reasoning, Gemini undervalued template structure (45\%) compared to Claude (96\%), with DeepSeek doing the same under-valuing (29\%). Since our feature analysis showed lower lexical diversity in the comment data, possibly making its structure more repetitive and template-like, this was a signal that both LLMs may have failed to exploit template structure for informal reasoning. In summary, we see a wide variation in reasoning patterns across domain and LLM judge, some of which appear to correlate between features and judge rationales.

\subsubsection{Dissociation between feature analysis and AI detection paradigms} Our study revealed a dissociation between the two detection paradigms (AUC vs. LLM-as-judge) across the domain and prompt conditions. Contradictory results such as higher detectability of the feature-evasive prompt on the embedding level as opposed to lower LLM-judge detectability suggests that for our data, the detection tools were not evaded by the same patterns. In contrast with \textcite{lu2024largelanguagemodelsguided}, whose feature-informed evasion held across six detectors, our study showed lowered LLM-judge detectability on our texts, but raised embedding separability. It is possible that the embedding classifier detects distributional, semantic-level regularities that are not removed by surface-feature manipulation by the evasive prompts. In addition, changes to the feature distribution, such as over-correction of text by the feature-evasive prompt, could have accentuated signals into another area of the embedding space that differs significantly from the embeddings of the human text. However, the surface level manipulations by the evasive prompts could indeed stand a better chance of fooling an LLM-as-judge, since they appear to rely on surface tells (lexical diversity, sentence length, hedging, punctuation) that the evasive prompt was designed to mask. Overall, this finding indicates a possible mismatch between the embedding classifier's distributional semantic reading and the surface-cue processing of the reasoning LLMs.

\subsection{Connection to related work}
Our detection results connect to prior work in three ways. Firstly, our LLM
judges performed considerably better on Swedish than the models tested by
\textcite{landberg2024aiwolf}, whose ChatGPT 3.5, Copilot and GPT-SW3 all
scored at or below chance on Swedish news texts in 2024. Our Claude judge reached ceiling
performance on several formal conditions, suggesting that LLM-as-judge has
become a viable detection paradigm for Swedish, at
least for the formal register. Secondly, our finding that the feature-evasive
prompt did not evade both detection paradigms contrasts with
\textcite{lu2024largelanguagemodelsguided}, whose feature-informed evasion held
across six detectors for English. Thirdly, the register-dependence
of our judges' performance is in keeping with \textcite{bao2024fastdetectgpt},
who found that supervised detectors performed comparably on in-distribution
English news but significantly worse on scientific and German writing,
indicating that detector performance is sensitive to both register and
language. Our judges' rationales (Table~\ref{tab:judge-rationale-cues}) were dominated by template structure in
the formal register, a cue that aligns with the more formal and structured LLM
text profile of \textcite{tang2023sciencedetectingllmgeneratedtexts}.

\section{Experiment 3 Results: AI detection by Swedish speakers}

The native and near-native speakers in our human detection study were slightly-above-random detectors for LLM-generated comments, with an accuracy above chance. Section 4.3.1 breaks down results among subgroups (age, language level profession, LLM usage); however, these within-subgroup results remain statistically underpowered and therefore most of the results are preliminary. 

\subsection{Accuracy, precision and recall}

\begin{table}[H]
  \centering
  \small
  \setlength{\tabcolsep}{5pt}
  \renewcommand{\arraystretch}{1.2}
  \begin{tabular}{l c c c c c c c}
    \toprule
    \textbf{Register} & \#P & $N$ & Acc & 95\% CI & $P_{\mathrm{hum}}$ & $R_{\mathrm{AI}}$ & $p$ \\
    \midrule
Informal (comments) & 37 & 827 & \cg{13}63.1 & [59.8, 66.4] & \crd{14}35.8 & \cg{7}57.4 & $.003$* \\
Formal (abstracts)  & 9  & 114 & \cg{3}52.6  & [43.5, 61.8] & \crd{22}27.9 & \crd{4}45.6 & $.460$ \\
    \bottomrule
  \end{tabular}
  \caption{AI-detection performance pooled by register (all values \%, one
  decimal place). \#P~=~participants contributing to that register;
  $N$~=~responses; Acc~=~pooled response accuracy; 95\%~CI~=~normal-approximation
  binomial interval on the pooled proportion, treating responses as independent;
  $P_{\mathrm{hum}}$~=~precision on the human class; $R_{\mathrm{AI}}$~=~recall
  on the AI class; $p$~=~two-sided one-sample $t$-test of participant accuracy
  against chance (50\%).}
  \label{tab:pooled-register}
\end{table}


% ===================== POOLED BY LANGUAGE LEVEL 

In terms of register, our participants reliably and significantly distinguished between human and LLM-Swedish forum comments above chance (63.1\%, CI 59.8\% - 66.4\%, p = .003), indicating that the Swedish speakers in our study were slightly-above-random detectors for LLM-generated Swedish comments. In the informal register, AI text recall was better than human text precision, indicating that Swedish speakers were over-labelling responses as AI-generated rather than detecting genuinely human comments. Formal abstracts appeared more difficult to distinguish, with an accuracy measurement close to chance (52.6\%, CI 43.5\%-61.8\%). The formal abstract responses also indicated lower human text precision than AI recall, suggesting an over-labelling of human abstracts as AI. Overall, this suggested that informal text was significantly easier to judge than formal text (+15.8 points), and that Swedish speakers did demonstrate a biased belief that human Swedish comments were AI-generated.

% ===================== POOLED BY LANGUAGE LEVEL =====================
\begin{table}[h!]
  \centering
  \small
  \setlength{\tabcolsep}{4pt}
  \renewcommand{\arraystretch}{1.15}
  \begin{tabular}{l c c ccc c ccc}
    \toprule
    & & \multicolumn{4}{c}{\textbf{Comments} (informal)}
        & \multicolumn{4}{c}{\textbf{Abstracts} (formal)} \\
    \cmidrule(lr){3-6}\cmidrule(l){7-10}
    \textbf{Language level} & \#P
     & $N$ & Acc & $P_{\mathrm{hum}}$ & $R_{\mathrm{AI}}$
     & $N$ & Acc & $P_{\mathrm{hum}}$ & $R_{\mathrm{AI}}$ \\
    \midrule
Native & 31 & 711 & \cg{12}61.9 & \crd{14}35.7 & \cg{5}55.4 & 114 & \cg{3}52.6 & \crd{22}27.9 & \crd{4}45.6 \\
Advanced & 3 & 62 & \cg{24}74.2 & \crd{12}38.1 & \cg{25}74.5 & 0 & -- & -- & -- \\
Intermediate & 3 & 54 & \cg{17}66.7 & \crd{14}36.0 & \cg{13}62.8 & 0 & -- & -- & -- \\
    \midrule
\textit{Overall} & 37 & 827 & \cg{13}63.1 & \crd{14}35.8 & \cg{7}57.4 & 114 & \cg{3}52.6 & \crd{22}27.9 & \crd{4}45.6 \\
    \bottomrule
  \end{tabular}
  \caption{AI-detection performance pooled by self-reported Swedish language
  level (all values \%, to one decimal place). \#P~=~participants in the group;
  $N$~=~responses given in that register; Acc~=~accuracy;
  $P_{\mathrm{hum}}$~=~precision on the human class; $R_{\mathrm{AI}}$~=~recall
  on the AI class. Cell shading: green marks values above chance (50\%), red
  below, with intensity proportional to the distance from 50. Dashes mark groups
  with no completed abstract responses. \textit{Overall} reports the
  $N$-weighted pooled means.}
  \label{tab:pooled-level}
\end{table}

Our participant pool comprised mostly of native speakers (n=31), who detected
forum comments above chance (accuracy of 61.9\%) and were the only group who
completed any abstract responses. The advanced and intermediate speakers did
show similar patterns to the native group, with an above-chance detection
accuracy of forum comments, and human precision also below AI recall.


% ===================== POOLED BY WEEKLY LLM USE =====================
\begin{table}[H]
  \centering
  \small
  \setlength{\tabcolsep}{4pt}
  \renewcommand{\arraystretch}{1.15}
  \begin{tabular}{l c c ccc c ccc}
    \toprule
    & & \multicolumn{4}{c}{\textbf{Comments} (informal)}
        & \multicolumn{4}{c}{\textbf{Abstracts} (formal)} \\
    \cmidrule(lr){3-6}\cmidrule(l){7-10}
    \textbf{Weekly LLM use} & \#P
     & $N$ & Acc & $P_{\mathrm{hum}}$ & $R_{\mathrm{AI}}$
     & $N$ & Acc & $P_{\mathrm{hum}}$ & $R_{\mathrm{AI}}$ \\
    \midrule
None & 1 & 7 & \crd{7}42.9 & \crd{50}0.0 & \crd{7}42.9 & 0 & -- & -- & -- \\
0--5 h & 15 & 382 & \cg{13}62.8 & \crd{12}37.7 & \cg{5}54.8 & 39 & \cg{12}61.5 & \crd{18}31.8 & \cg{3}53.1 \\
5--10 h & 9 & 182 & \cg{18}67.6 & \crd{11}38.6 & \cg{14}64.1 & 60 & \cg{0}50.0 & \crd{20}29.7 & \crd{8}42.2 \\
10--20 h & 5 & 74 & \cg{8}58.1 & \crd{26}23.5 & \cg{7}57.4 & 0 & -- & -- & -- \\
20+ h & 7 & 182 & \cg{12}62.1 & \crd{15}35.1 & \cg{7}56.7 & 15 & \crd{10}40.0 & \crd{39}11.1 & \crd{12}38.5 \\
    \midrule
\textit{Overall} & 37 & 827 & \cg{13}63.1 & \crd{14}35.8 & \cg{7}57.4 & 114 & \cg{3}52.6 & \crd{22}27.9 & \crd{4}45.6 \\
    \bottomrule
  \end{tabular}
  \caption{AI-detection performance pooled by self-reported weekly LLM use
  (hours). Columns and shading as in Table~\ref{tab:pooled-level}.
  \#P~=~participants; $N$~=~responses per register. Dashes mark groups with no
  completed abstract responses.}
  \label{tab:pooled-llm}
\end{table}

Weekly LLM use showed no particular relationship with detection accuracy on
comments, with the 5-10 h group performing best (67.6\%), while the heaviest users
(20+ h) did not outperform the lightest (62.1\% vs.\ 62.8\%). On abstracts,
accuracy fell as reported use rose, though the 20+ h abstract cell rests on only
15 responses from two participants and should not be interpreted as significant.

\subsection{Human participant reasoning rationales}
% Thematic summary of human participants' AI-detection rationales.
% Uses tabularx/booktabs (already in the thesis preamble).
\begin{table}[H]
  \centering
  \footnotesize
  \setlength{\tabcolsep}{5pt}
  \renewcommand{\arraystretch}{1.2}
  \begin{tabularx}{\linewidth}{@{}>{\raggedright\arraybackslash}p{3.2cm} X@{}}
    \toprule
    \textbf{Cue theme} & \textbf{Representative participant wording (translated)} \\
    \midrule
    Excessive length / over-explaining
      & ``Unnecessarily long comments.'' \\
    \addlinespace[2pt]
    Too-clean grammar \& spelling
      & ``Too-good grammar''; the absence of \emph{särskrivning} (wrongly splitting compound words, a common error in Swedish); the ``well-written'' comments are easy to spot as AI. \\
    \addlinespace[2pt]
    Too formal; no slang or swearing
      & ``No swearing''; people write very casually in comments, so the polished ones stand out. \\
    \addlinespace[2pt]
    Formulaic phrasing \& structure
      & ``Two words of the same class where one would suffice (`a fast and efficient runner')''; the ``not only X but also Y'' pattern; ``likes to list (often three) things''; ``enumerations''. Sentences open with connectives (``In summary\ldots''); ``sentence construction''; ``a clearer structure''. \\
    \addlinespace[2pt]
    Vague filler vs.\ over-factual
      & ``Fluffy and vacuous\ldots\ too many `generic' words'' \emph{vs.}\ ``too fact-based''. \\
    \addlinespace[2pt]
    Dash punctuation (em-dash / hyphen)
      & ``Dashes signal AI''; ``the mark~---''. \\
    \addlinespace[2pt]
    Reader-directed moves
      & ``Asks what others think''\,/\,``What do you all think?''; ``closing calls-to-action''; compares ``this versus something else''. \\
    \bottomrule
  \end{tabularx}
  \caption{\footnotesize Thematic coding of participants' reasoning rationales of how they
  identified AI-generated Swedish text. Participant wording translated from Swedish.}
  \label{tab:human-rationale-cues}
\end{table}
The human participants gave several reflections as to their reasoning for classifying LLM-Swedish, the key themes being writing style that is too formal (formulaic, clearly structure, perfect spelling), specific punctuation markers such as em-dashes, reader-directed moves (e.g. "what do you all think"?), vague and over-factual responses, and excessive lengths. These rationales are in keeping with \textcite{dugan2024detecting}, where participants reported similar heuristics,  and \cite{munoz-ortiz2024contrasting} which did measure more formulaic patterns in its LLM corpus data. That said, certain heuristics reported by the reader may have been a result of the curated dataset used for the public-facing quiz, for example, the comment that there is 'no swearing' in the LLM quiz comments does not hold for the wider corpus. Overall, several heuristics were in line with the feature analysis studies findings, namely the length, formulaic phrasing, and vagueness of texts.

\subsection{Comparison with LLMs as judge}

\begin{table}[h!]
  \centering
  \footnotesize
  \setlength{\tabcolsep}{5pt}
  \renewcommand{\arraystretch}{1.2}
  \begin{tabular}{l l c c c c c c}
    \toprule
    \textbf{Judge} & \textbf{Register} & $N$ & Acc & 95\% CI & $P_{\mathrm{hum}}$ & $R_{\mathrm{AI}}$ & $p$ \\
    \midrule
    \multirow{2}{*}{Claude}
      & Formal   & 20 & \cg{50}100.0 & [83.9, 100.0] & \cg{50}100.0 & \cg{50}100.0 & $<$.001* \\
      & Informal & 40 & \cg{22}72.5  & [57.2, 83.9]  & \crd{2}47.6  & \cg{13}63.3  & $.006$* \\
    \midrule
    \multirow{2}{*}{DeepSeek}
      & Formal   & 20 & \cg{35}85.0  & [64.0, 94.8]  & \cg{17}66.7  & \cg{37}86.7  & $.003$* \\
      & Informal & 40 & \crd{18}32.5 & [20.1, 48.0]  & \crd{23}27.0 & \crd{40}10.0 & $.038$* \\
    \midrule
    \multirow{2}{*}{Gemini}
      & Formal   & 20 & \cg{45}95.0  & [76.4, 99.1]  & \cg{33}83.3  & \cg{43}93.3  & $<$.001* \\
      & Informal & 40 & \crd{5}45.0  & [30.7, 60.2]  & \crd{19}31.2 & \crd{23}26.7 & $.636$ \\
    \bottomrule
  \end{tabular}
  \caption{AI-detection performance pooled by register for each LLM judge on
  the 60-item human quiz (20 formal abstracts, 40 informal comments;
  prompt conditions pooled
  within register). The $95\%$
  CI is a Wilson score interval on that single judge's per-item accuracy. Acc~=~overall accuracy;
  $P_{\mathrm{hum}}$~=~precision on the human class; $R_{\mathrm{AI}}$~=~recall
  on the AI class; $p$~=~two-sided exact binomial test of accuracy against
  chance (50\%). Shading as in Table~\ref{tab:pooled-level}: green above
  chance, red below, intensity $\propto$ distance from 50.}
  \label{tab:llm-judge-pooled-register}
\end{table}

As a like-for-like comparison, the LLMs as judge were employed to classify the exact same quiz as the human participants. Interestingly, while the human subjects seemed to find the informal comments easier to classify as AI than the formal abstracts, this was reversed in the LLMs' classifications -- all judges performed comparatively lower in the informal register. Claude did perform better than the human participants, with 100\% accuracy, human precision and AI recall on the formal register, and 72.5\% accuracy on the informal register. Gemini and DeepSeek did perform better than humans in the formal domain (95\% and 85\% accuracy respectively; however, they performed worse in the informal domain, with DeepSeek achieving the lowest score at 32.5\% (CI 20.1 - 48\%) and Gemini at 45\% (CI 30.7-60.2\%). 

\subsection{Connection to related work}
Our human detection results are broadly in keeping with previous work on human discernment of AI text. \textcite{doruetal} found that German medical experts and non-experts labelled texts as human-written when they contained spelling or grammatical errors and greater variation in sentence structure, and as AI-written when they were monotonous, superficial or generic. Our Swedish participants reported closely comparable heuristics (Table~\ref{tab:human-rationale-cues}), particularly around too-clean grammar and formulaic phrasing, suggesting that speakers of different languages may use similar cues. \textcite{dugan2024detecting} similarly reported that participants relied on surface heuristics of this kind, while \textcite{clark2021all} established that untrained human raters distinguish AI text at close to chance. Our results did not replicate \textcite{russell2025people}, who found frequent LLM users to be accurate detectors of AI-generated English text, though our subgroup results are underpowered and a larger participant pool would be required to test this relationship properly.

\chapter{Limitations}

\section{Textual analysis results}
Two main limitations point to 1) the temporal mismatch between our human- and LLM-written data, and 2) the choice of generator for our feature analysis results.

\subsection{Temporal mismatch}
All human-written data in our study was authored before 2017, whereas the LLM texts were generated in 2026. This gap was a deliberate choice to avoid LLM contamination of the human data; however, it could have introduced era-linked stylistic differences unrelated to authorship, especially for informal forum comment data, such as slang, references, and platform conventions, all of which could not be cleanly isolated in our study. An alternative could have been to get humans to write responses to forum comments, but this may not have been an organic replication of the environment. Due to the anonymity of the forum platforms we used, it is very difficult to detect the extent to which AI use would have affected 2026 comments. Nonetheless, we flag this as a gap in our study and a challenge for this research topic going forward as the gap between the pre-transformer LLM era widens in the future.

\subsection{Single generator for feature analysis}
Secondly, our feature analysis was conducted on a single generator (GPT-5.2). As a result, our feature-evasive prompt was designed with feature results from only one generator. While we did test the robustness of the feature-evasive prompt output by extending our detection experiments to include a later version of GPT and a model from another family (Mistral) as generators, this did not change the prompt design itself. Therefore, the feature-level characterisation of LLM-Swedish rested on only one model and one family, and the over-corrective behaviour we observed under the evasive prompt cannot be claimed as model-independent without further testing. In addition, a Scandinavian-tuned model such as GPT-SW3 was not included as a generator for our feature analysis due to access constraints at the time of writing. However, further work involving a Scandinavian-centric LLM would be a valuable point of comparison between general multilingual and language-specific LLMs.


\section{AI detection results}

The detection experiments are subject to two main limitations: the statistical power of the human study, and the limited architectural transparency for most of the LLMs involved.

\subsubsection{Insufficient statistical power of human study}
The biggest limitation is the statistical underpower of the human detection study, particularly for the formal abstract register results. Many participants answered only the informal comments section and did not continue to the abstracts (likely due to the higher reading burden), and as a result, some quiz items are far better represented than others.  The within-subgroup breakdowns by age, language level, profession, and weekly LLM use are very statistically underpowered, and should therefore be interpreted as preliminary. These results do motivate a larger follow-up study, whether it be a design requiring completion of both registers equally, providing compensation to human participants for their time, or rotating the number of items per participants. A larger study would provide better and more robust insights into the effect of variables such as LLM usage, age, and profession of Swedish speaker participants. 

\subsubsection{Lack of transparency of proprietary models}
Interpretation of the cross-generator detection results is further constrained by the proprietary, closed nature of most of the models involved. Where a judge such as DeepSeek-R1 underperformed on Swedish, we can only speculate about the contributing factors such as parameter scale, training-data language composition, and architecture since these are not disclosed for most of the models used. The one exception is DeepSeek-R1's reported parameter count and its English- and Chinese-dominant training mix, which we tentatively connect to its weaker Swedish performance; but this connection remains a hypothesis rather than a demonstrated cause.

\subsubsection{Partially-edited or LLM-assisted text}

Our study only reports a comparison between fully-human-written Swedish texts and texts that are entirely generated (zero-shot) by LLMs. However, \textcite{wu2024detectrl} reports that after applying LLM to generate a text, humans often edit and modify certain words or passages. Humans may also write an original text themselves before using an LLM to modify its structures. Current mainstream LLM detectors show evidence of degrading in performance when categorising paraphrased text \parencite{wolff2020attacking}. Future work into this topic would benefit from a comparison between fully-human, partially-LLM and fully-LLM-authored works. As our study's LLM-Swedish was generated fully, technique would be less comparable for studies that look at partially LLM-assisted text. This type of study would be interesting to conduct in future research.

\section{Generalisation of patterns to other models and genres}

Our study focused on a specific set of tasks with a specific set of models. In the fast-changing landscape of LLM development, it is difficult to determine if the findings of our study will hold in the future, or whether new architectures and implementations will change in the near-future. As a result, we emphasise our study not as an exhaustive and definitive set of results that characterise LLM-Swedish, but rather a method that can be transferable to future studies with newer models, genres and languages.

\chapter{Conclusion}
This thesis presented a feature-based analysis and AI-detection evaluation of LLM-generated Swedish, with several key findings. Firstly, that LLM Swedish systematically deviated from human Swedish across nearly all features within a formal and informal domain, particularly in sentence complexity, lexical diversity, and punctuation usage. Secondly, the feature-evasive prompt did successfully mitigate some AI signals, namely named-entity-recognition (NER). However, it over-corrected these signals in several measures (sentence complexity, lexical diversity and full stop rate). The over-correction showed that this method of evasive prompting can inadvertently create another AI signal in the opposite direction of the original feature analysis. Thirdly, the two AI detection paradigms diverged on the same texts, with the feature-evasive prompt more successfully evading detection by LLM-as-judge classifiers, particularly for the informal register. The embedding-based classifier, however, raised separability for the evasive prompt for some generators rather than lowering its score uniformly. This indicated that surface-level evasion appears to succeed against judges that attend to particular surface cues, but fail against a classifier reading distributional structure that the prompt could not target. The successful evasion of the judges was strongly register-dependent and largely confined to the informal register. Lastly, Human Swedish speakers were slightly-above-chance detectors of LLM-Swedish forum comments, but performed at chance on formal abstracts (however, this result was preliminary). The detection pattern reversed for the LLM judges, which found the formal abstracts easier than the informal comments. Participants tended to over-label human text as AI, and their reported heuristics (excessive length, too-clean grammar, formulaic phrasing, em-dashes, reader-directed questions) aligned with several of the feature-analysis signals. Human within-subgroup results (age, language level, LLM usage) were statistically underpowered and are reported as preliminary. Overall, the thesis welcomes future work that can measure results for more registers, such as news media or literary works, with more models (including open-source models), with subgroup breakdowns (age, profession and LLM usage) and across more languages.




\printbibliography

\chapter{Appendix}

\section{Generation prompt conditions}
\label{app:prompts}

This section gives the full text of the three prompt conditions for each
register. All prompts were
written in Swedish; English translations are given in italics.

\begin{table}[H]
\centering
\begin{tabular}{p{0.92\linewidth}}
\toprule
\textbf{Base instruction (\texttt{baseline})} \\
\midrule
\textit{Swedish:} ``Titta på titeln och nyckelorden och skapa en sammanfattning
i kandidatuppsatsstil med dina egna ord: \textlangle title, keywords\textrangle.''

\textit{English: ``Look at the title and keywords and create a summary of the
thesis abstract in your own words: \textlangle title, keywords\textrangle.''} \\
\midrule
\textbf{Suffix appended for \texttt{generic\_evasive}} \\
\midrule
\textit{Swedish:} ``Skriv den så mänskligt som möjligt, så att texten inte går
att skilja från en uppsats skriven av en människa.''

\textit{English: ``Write it as humanly as possible, so that the text cannot be
distinguished from an essay written by a human.''} \\
\midrule
\textbf{Suffix appended for \texttt{feature\_evasive}} \\
\midrule
\textit{Swedish:} ``Skriv så att texten inte kan identifieras som
AI-genererad, och efterlikna de statistiska drag som utmärker mänskligt
skrivna kandidatuppsatser. Återanvänd samma nyckeltermer och fraser
ordagrant snarare än att variera med synonymer, och sträva inte efter
maximal ordvariation. Undvik komprimerad nominalstil; använd hellre
finita verb, pronomen och bindeord. Skriv övervägande korta meningar med
fler punkter och färre kommatecken, och använd inte tankstreck. Håll
gardering och epistemiska uttryck till ett minimum (t.ex. `kanske',
`möjligen', `tycks', `kan tänkas'), och kompensera inte genom att lägga
till talspråkliga garderingsord. Föredra korta, vardagliga ord framför
långa, latinska eller formella termer. Var konkret och nämn specifika
namn, begrepp, verk och årtal där det är möjligt. Undvik onödig negation.
Förklara inte över.''

\textit{English:} ``\textit{Write so that the text cannot be identified as
AI-generated, and imitate the statistical features that characterise
human-written Bachelor's theses. Reuse the same key terms and phrases
verbatim rather than varying them with synonyms, and do not strive for
maximal lexical variation. Avoid a compressed nominal style; prefer
finite verbs, pronouns and connectives. Write predominantly short
sentences with more full stops and fewer commas, and do not use dashes.
Keep hedging and epistemic expressions to a minimum (e.g.\ `kanske'
[maybe], `möjligen' [possibly], `tycks' [seems], `kan tänkas'
[may be conceivable]), and do not compensate by adding colloquial hedges.
Prefer short, everyday words over long, Latinate or formal terms. Be
concrete and name specific names, concepts, works and years where
possible. Avoid unnecessary negation. Do not over-explain.}'' \\
\bottomrule
\end{tabular}
\caption{Abstract-generation prompt conditions.}
\label{tab:prompt-conditions}
\end{table}

\begin{table}[H]
\centering
\begin{tabular}{p{0.92\linewidth}}
\toprule
\textbf{Base instruction (\texttt{baseline})} \\
\midrule
``Svara på följande fråga med en svensk kommentar, i stil med
en kommentar på ett svenskt forum. Skriv bara kommentaren, inget annat. Fråga:
\texttt{\{question\}}''

\textit{English: ``Answer the following question with a Swedish comment, in the
style of a comment on a Swedish forum. Write only the comment, nothing else.
Question: \texttt{\{question\}}''} \\
\midrule
\textbf{Suffix appended for \texttt{generic\_evasive}} \\
\midrule
``Skriv så mänskligt som möjligt, så att kommentaren inte går
att skilja från en riktig svensk forumanvändares kommentar.''

\textit{English: ``Write as humanly as possible, so that the comment cannot be
distinguished from a real Swedish forum user's comment.''} \\
\midrule
\textbf{Suffix appended for \texttt{feature\_evasive}} \\
\midrule
``Skriv så att texten inte kan identifieras som AI-genererad,
och efterlikna hur människor skriver i informella kommentarer. Återanvänd inte
samma fraser flera gånger, utan formulera om dig. Skriv innehållstätt med färre
funktionsord; använd gärna kortformer och talspråkliga former. Skriv korta
meningar med fler punkter och färre kommatecken, och använd inte tankstreck.
Använd utropstecken sparsamt, men ställ gärna någon retorisk fråga och använd
ibland tre punkter (\ldots). Överdriv inte med garderingsord (särskilt inte
talspråkliga som `typ', `liksom', `ju' och `väl'), och undvik förstärkningsord
som `verkligen', `absolut' och `väldigt'. Håll epistemiska uttryck på en låg
nivå. Föredra korta, vardagliga ord. Använd gärna nekande satser (med `inte',
`aldrig' osv.) där det passar. Var konkret och nämn specifika namn där det går.
Förklara inte över.''

\textit{English: ``Write so that the text cannot be identified as AI-generated,
and imitate how people write in informal comments. Do not reuse the same phrases
several times; rephrase instead. Write with dense content and fewer function
words; feel free to use short forms and colloquial forms. Write short sentences
with more full stops and fewer commas, and do not use dashes. Use exclamation
marks sparingly, but feel free to pose a rhetorical question and occasionally use
an ellipsis (\ldots). Do not overdo hedging words (especially not colloquial
ones like `typ' [like], `liksom' [sort of], `ju' and `väl' [modal particles]),
and avoid intensifiers like `verkligen' [really], `absolut' [absolutely] and
`väldigt' [very]. Keep epistemic expressions at a low level. Prefer short,
everyday words. Feel free to use negated clauses (with `inte' [not], `aldrig'
[never] etc.) where it fits. Be concrete and name specific names where possible.
Do not over-explain.''} \\
\bottomrule
\end{tabular}
\caption{Forum-comment generation prompt conditions.}
\label{tab:prompt-conditions-informal}
\end{table}

\clearpage
\section{Examples of generated text}
\label{app:examples}

Tables \ref{tab:examples-formal} and \ref{tab:examples-informal} show the
opening of the human-written text alongside the three LLM prompt conditions,
for one example abstract and one example forum question respectively.

\begin{table}[H]
\centering
\begin{tabularx}{\linewidth}{@{}l X@{}}
\multicolumn{2}{@{}p{\linewidth}@{}}{\footnotesize\textbf{Formal register, Bachelor's abstracts example.}} \\
\midrule
Condition & Abstract opening \\
\midrule
\textit{Human} & I ``Samkönade äktenskaps vara eller icke vara inom Svenska kyrkan utförs'', en argumentationsanalys, baserad på hermeneutiks metod och kvalitativ textanalys, av debattartiklar som behandlat frågan som Svenska\,\ldots{}
\newline{\itshape In ``Same-sex marriage's `to be or not to be' within the Church of Sweden'', an argumentation analysis, based on the hermeneutic method and qualitative text analysis, is carried out on opinion pieces that addressed the question of whether the Swedish\,\ldots} \\[5pt]
\midrule
\textit{Baseline} & Uppsatsen undersöker frågan om samkönade äktenskaps ställning inom Svenska kyrkan genom en argumentationsanalys av debattartiklar publicerade i Kyrkans Tidning och Dagens Nyheter. Utifrån ett urval texter\,\ldots{}\newline{\itshape The thesis examines the standing of same-sex marriage within the Church of Sweden through an argumentation analysis of opinion pieces published in Kyrkans Tidning and Dagens Nyheter. From a selection of texts\,\ldots} \\[5pt]
\midrule
\textit{generic-evasive} & Denna studie undersöker frågan om samkönade äktenskaps ställning inom Svenska kyrkan och hur den har motiverats och bestridits i det offentliga samtalet. Utifrån en argumentationsanalys granskas debattartiklar\,\ldots{}\newline{\itshape This study examines the standing of same-sex marriage within the Church of Sweden and how it has been justified and contested in public discourse. Through an argumentation analysis, opinion pieces are examined\,\ldots} \\[5pt]
\midrule
\textit{feature-evasive} & Uppsatsen tar upp frågan om samkönade äktenskaps vara eller icke vara inom Svenska kyrkan. Den fokuserar på hur debatten ser ut i tryckta debattartiklar. Materialet kommer från Kyrkans tidning och Dagens\,\ldots{}\newline{\itshape The thesis takes up the question of same-sex marriage's existence or non-existence within the Church of Sweden. It focuses on what the debate looks like in printed opinion pieces. The material comes from Kyrkans Tidning and Dagens\,\ldots} \\[5pt]
\bottomrule
\end{tabularx}
\caption{Comparison of human and LLM-generated Swedish text, formal domain, for one example abstract.}
\label{tab:examples-formal}
\end{table}

\begin{table}[H]
\centering
\begin{tabularx}{\linewidth}{@{}l X@{}}
\multicolumn{2}{@{}p{\linewidth}@{}}{\footnotesize\textbf{Informal register: forum comments example.}

\textbf{Thread question:} Ingen som vet vart man kan köpa Stenströms under bara skjortor -- på internet? {\footnotesize\itshape (Anyone know where you can buy Stenströms undershirts online?)}} \\
\midrule
Condition & Comment opening \\
\midrule
\textit{Human} & Har lite här i stockholm med.\newline{\footnotesize\itshape Got a few here in Stockholm.} \\[5pt]
\midrule
\textit{Baseline} & Kolla in Stenströms egen hemsida, de har ofta ett bra urval av skjortor där. Annars kan du även kolla på Zalando eller CareOfCarl, de brukar ha ett bra sortiment av deras kläder. Lycka till!\newline{\footnotesize\itshape Check out Stenströms' own website, they often have a good selection of shirts there. Otherwise you can also check Zalando or CareOfCarl, they usually carry a good range of their clothes. Good luck!} \\[5pt]
\midrule
\textit{generic-evasive} & Kolla typ CareOfCarl (brukar ha Stenströms), annars NK/Åhléns online ibland. Jag har även sett deras undertröjor dyka upp på Zalando/Brandos-liknande sidor men utbudet varierar. Ett tips är att söka på\,\ldots{}\newline{\footnotesize\itshape Check out somewhere like CareOfCarl (usually has Stenströms), otherwise NK/Åhléns online sometimes. I've also seen their undershirts show up on Zalando/Brandos-type sites, but the range of stock varies. One trick is to search for\,\ldots} \\[5pt]
\midrule
\textit{feature-evasive} & Kollat Stenströms egna webshop? De har ibland ``undershirt'' under accessoarer. Annars brukar Care of Carl ha Stenströms och ibland undertröjor. NK och Johnells säljer märket online men inte alltid just under.\,\ldots{}\newline{\footnotesize\itshape Have you checked Stenströms' own webshop? They sometimes list ``undershirt'' under accessories. Otherwise Care of Carl usually has Stenströms and sometimes undershirts. NK and Johnells sell the brand online but not always undershirts specifically\,\ldots} \\[5pt]
\bottomrule
\end{tabularx}
\caption{Comparison of human and LLM-generated Swedish text, informal domain, for one example forum question.}
\label{tab:examples-informal}
\end{table}

\begin{table}[h!]
\centering
\footnotesize
\begin{tabular}{p{0.92\linewidth}}
\toprule
\textbf{Detection prompt (\texttt{formal} --- thesis abstract)} \\
\midrule
``Du deltar i en studie som handlar om att identifiera AI-genererad text på
svenska. Analysera följande kandidatuppsatsabstrakt och avgör om det är skrivet av en AI
eller av en människa. Titta på saker som språklig variation, formuleringsval,
meningsbyggnad, och om texten känns autentiskt akademisk eller mallartad.

Abstrakt: \{text\}

Svara med JSON i exakt detta format (ingen annan text utanför JSON-blocket):
\{"classification": "AI" or "MÄNNISKA", "confidence": \textless 0--100\textgreater,
"reasoning": "\textless kort förklaring på 50--100 ord\textgreater"\}''

{\itshape\footnotesize English: ``You are participating in a study about identifying
AI-generated text in Swedish.

Analyze the following bachelor's thesis abstract and determine whether it is
written by an AI or by a human. Look at things like linguistic variation, choice
of phrasing, sentence structure, and whether the text feels authentically
academic or template-like.

Abstract: \{text\}

Answer with JSON in exactly this format (no other text outside the JSON block):
\{"classification": "AI" or "HUMAN", "confidence": \textless 0--100\textgreater,
"reasoning": "\textless short explanation of 50--100 words\textgreater"\}''\par} \\
\midrule
\textbf{Detection prompt (\texttt{informal} --- forum comment)} \\
\midrule
``Du deltar i en studie som handlar om att identifiera AI-genererad text på
svenska.

Analysera följande forumkommentar och avgör om den är skriven av en AI eller av
en människa. Titta på saker som informellt/vardagligt språk, stavfel, personliga
referenser, naturlig röst, och om texten känns genuin eller generisk.

Kommentar: \{text\}

Svara med JSON i exakt detta format (ingen annan text utanför JSON-blocket):
\{"classification": "AI" or "MÄNNISKA", "confidence": \textless 0--100\textgreater,
"reasoning": "\textless kort förklaring på 50--100 ord\textgreater"\}''

{\itshape\footnotesize English: ``You are participating in a study about identifying
AI-generated text in Swedish.

Analyze the following forum comment and determine whether it is written by an AI
or by a human. Look at things like informal/colloquial language, spelling
mistakes, personal references, natural voice, and whether the text feels genuine
or generic.

Comment: \{text\}

Answer with JSON in exactly this format (no other text outside the JSON block):
\{"classification": "AI" or "HUMAN", "confidence": \textless 0--100\textgreater,
"reasoning": "\textless short explanation of 50--100 words\textgreater"\}''\par} \\
\bottomrule
\end{tabular}
\caption{Detection prompt conditions by register.}
\label{tab:detection-prompt-conditions}
\end{table}

\newpage


\section{Full feature analysis}

\begin{table}[H]
  \centering
  \small
  \begin{tabular}{l rr}
    \toprule
    & Formal & Informal \\
    \midrule
    \multicolumn{3}{l}{\textit{Median $|d_z|$ vs.\ human}} \\
    \quad GPT-5.2      & 0.51 & 0.18 \\
    \quad GPT-5.6      & 0.27 & 0.14 \\
    \quad Mistral-3.2  & 0.29 & 0.19 \\
    \addlinespace[3pt]
    \multicolumn{3}{l}{\textit{Median $|d_z|$, LLM vs.\ LLM}} \\
    \quad GPT-5.2 vs.\ GPT-5.6      & 0.33 & 0.20 \\
    \quad GPT-5.2 vs.\ Mistral-3.2  & 0.28 & 0.28 \\
    \quad GPT-5.6 vs.\ Mistral-3.2  & 0.21 & 0.29 \\
    \addlinespace[3pt]
    \multicolumn{3}{l}{\textit{GPT-5.2-significant features also significant in\ldots}} \\
    \quad GPT-5.6 only        & 72\% & 37\% \\
    \quad Mistral-3.2 only    & 86\% & 89\% \\
    \quad both                & 64\% & 33\% \\
    \bottomrule
  \end{tabular}
  \caption{Model-dependence of the feature analysis: GPT-5.2 vs.\ GPT-5.6 vs.\
  Mistral-3.2 (mistral-small-2506), paired Wilcoxon feature comparison across
  baseline / human-like / detector-evasive conditions. GPT-5.6 is a 100-document
  quiz subset; all other cells use the full corpus (170 formal / \textasciitilde{}1{,}150
  informal). See Table~\ref{app:three-way-formal-gpt52-vs-gpt56} et seq.\ for the
  full per-feature breakdown.}
  \label{tab:model-dependence}
\end{table}

% ---- formal register, Baseline condition: all features, human vs 3 models (app:full-three-way-formal-baseline) ----
% GENERATED by src/2_text_analysis_scripts/scripts/make_latex_tables.py
% Source of truth: csv_files/significance_tests_results.csv
% Contents: complete formal feature set (42 features x 3 conditions), paired human-vs-GPT-5.2 Cohen's d_z
% REQUIRES: \usepackage{booktabs}, \usepackage[table]{xcolor}
\providecommand{\cg}[1]{\cellcolor{hmgreen!#1}}
\providecommand{\crd}[1]{\cellcolor{hmred!#1}}
% \definecolor{hmgreen}{HTML}{2E8B57} % uncomment if used standalone
% \definecolor{hmred}{HTML}{C0392B}   % uncomment if used standalone



% Table 7.1 goes here — per-participant human AI-detection results
% (no-feedback set, N=31). Uses existing \cg / \crd macros.

\providecommand{\cg}[1]{\cellcolor{hmgreen!#1}}
\providecommand{\crd}[1]{\cellcolor{hmred!#1}}
% \definecolor{hmgreen}{HTML}{2E8B57} % uncomment if used standalone
% \definecolor{hmred}{HTML}{C0392B}   % uncomment if used standalone

% ---- FULL feature set, formal register (app:full-features-formal) ----
% REQUIRES: booktabs + the \cg / \crd shading macros from the preamble.
% 42 features x 3 conditions.
% If it overflows the page, wrap in longtable or switch to table*.
% NOTE: condition(s) ['detector_aware'] present in the CSV but omitted here (house convention); pass conds= to include them.
\documentclass{article}
\usepackage[margin=1.5cm,top=1.2cm,bottom=1.2cm]{geometry}
\usepackage{booktabs}
\usepackage[table]{xcolor}
\definecolor{hmgreen}{HTML}{2E8B57}
\definecolor{hmred}{HTML}{C0392B}
\newcommand{\cg}[1]{\cellcolor{hmgreen!#1}}
\newcommand{\crd}[1]{\cellcolor{hmred!#1}}
\pagestyle{empty}
\begin{document}

% ---- formal register, Baseline condition: all features, human vs 3 models (app:full-three-way-formal-baseline) ----
% REQUIRES: booktabs + the \cg / \crd shading macros from the preamble.
% 63 features x 3 models.
% Sized to fit one page at \scriptsize with arraystretch 0.9; if it still overflows, wrap in longtable.
\begin{table}[htbp]
  \centering
  \scriptsize
  \setlength{\tabcolsep}{2.5pt}
  \renewcommand{\arraystretch}{0.9}
  \begin{tabular}{l rrrr ccc}
    \toprule
    & \multicolumn{4}{c}{Median} & \multicolumn{3}{c}{Cohen's $d_z$ (vs.\ human)} \\
    \cmidrule(lr){2-5}\cmidrule(l){6-8}
    Feature & H & GPT-5.2 & GPT-5.6 & Mistral & GPT-5.2 & GPT-5.6 & Mistral \\
    \midrule
    \multicolumn{8}{l}{\textit{Stylometric surface} (11)} \\
    \addlinespace[0.5pt]
    \texttt{punct\_comma}                        & 1.227 & 4.722 & 5.698 & 5.534 & \cg{42}$+1.67$* & \cg{48}$+1.92$* & \cg{45}$+1.78$* \\
    \texttt{mtld}                                & 100.081 & 153.897 & 147.404 & 113.300 & \cg{36}$+1.43$* & \cg{34}$+1.36$* & \cg{11}$+0.45$* \\
    \texttt{mattr}                               & 0.827 & 0.867 & 0.873 & 0.849 & \cg{26}$+1.03$* & \cg{28}$+1.13$* & \cg{16}$+0.65$* \\
    \texttt{punct\_period}                       & 4.680 & 3.663 & 4.908 & 4.301 & \crd{22}$-0.87$* & \cg{5}$+0.20$ & \crd{7}$-0.28$* \\
    \texttt{function\_word\_rate}                & 0.424 & 0.377 & 0.347 & 0.359 & \crd{22}$-0.87$* & \crd{36}$-1.42$* & \crd{29}$-1.16$* \\
    \texttt{bigram\_repetition\_rate}            & 0.046 & 0.024 & 0.023 & 0.048 & \crd{19}$-0.74$* & \crd{20}$-0.81$* & \crd{4}$-0.17$ \\
    \texttt{trigram\_repetition\_rate}           & 0.013 & 0.004 & 0.000 & 0.008 & \crd{16}$-0.63$* & \crd{18}$-0.73$* & \crd{7}$-0.28$* \\
    \texttt{ttr}                                 & 0.632 & 0.689 & 0.718 & 0.664 & \cg{14}$+0.54$* & \cg{20}$+0.81$* & \cg{6}$+0.22$* \\
    \texttt{punct\_dash}                         & 0.000 & 0.000 & 0.000 & 0.000 & \cg{13}$+0.53$* & \cg{2}$+0.07$ & \cg{8}$+0.32$* \\
    \texttt{punct\_semicolon}                    & 0.000 & 0.000 & 0.000 & 0.000 & \crd{11}$-0.45$* & \crd{11}$-0.43$* & \crd{12}$-0.46$* \\
    \texttt{punct\_colon}                        & 0.000 & 0.000 & 0.513 & 0.549 & \crd{2}$-0.07$ & \cg{0}$+0.01$ & \cg{6}$+0.24$* \\
    \addlinespace[1.5pt]
    \multicolumn{8}{l}{\textit{Pragmatic markers} (14)} \\
    \addlinespace[0.5pt]
    \texttt{hedge\_neutral\_rate}                & 0.712 & 1.692 & 1.893 & 1.617 & \cg{19}$+0.75$* & \cg{23}$+0.91$* & \cg{16}$+0.65$* \\
    \texttt{hedge\_neutral\_count}               & 1.000 & 3.000 & 3.000 & 3.000 & \cg{17}$+0.68$* & \cg{19}$+0.74$* & \cg{12}$+0.46$* \\
    \texttt{hedge\_rate}                         & 0.815 & 1.852 & 1.893 & 1.720 & \cg{16}$+0.64$* & \cg{19}$+0.75$* & \cg{15}$+0.59$* \\
    \texttt{hedge\_count}                        & 1.000 & 4.000 & 3.000 & 3.000 & \cg{14}$+0.57$* & \cg{14}$+0.56$* & \cg{9}$+0.34$* \\
    \texttt{epistemic\_rate}                     & 0.000 & 0.000 & 0.000 & 0.000 & \crd{12}$-0.47$* & \crd{9}$-0.35$* & \crd{10}$-0.40$* \\
    \texttt{epistemic\_count}                    & 0.000 & 0.000 & 0.000 & 0.000 & \crd{11}$-0.42$* & \crd{10}$-0.41$* & \crd{11}$-0.42$* \\
    \texttt{negation\_rate}                      & 0.347 & 0.000 & 0.000 & 0.000 & \crd{7}$-0.28$* & \crd{4}$-0.14$ & \crd{12}$-0.48$* \\
    \texttt{negation\_count}                     & 1.000 & 0.000 & 0.000 & 0.000 & \crd{6}$-0.22$* & \crd{5}$-0.21$ & \crd{13}$-0.52$* \\
    \texttt{hedge\_informal\_count}              & 0.000 & 0.000 & 0.000 & 0.000 & \crd{3}$-0.13$ & \crd{6}$-0.25$* & \crd{3}$-0.11$ \\
    \texttt{n\_alpha\_tokens}                    & 178.500 & 197.500 & 162.000 & 163.000 & \cg{3}$+0.10$* & \crd{5}$-0.19$ & \crd{6}$-0.22$* \\
    \texttt{negation\_sentence\_rate}            & 0.051 & 0.000 & 0.000 & 0.000 & \crd{2}$-0.08$ & \crd{3}$-0.13$ & \crd{11}$-0.44$* \\
    \texttt{hedge\_informal\_rate}               & 0.000 & 0.000 & 0.000 & 0.000 & \crd{1}$-0.03$ & \crd{6}$-0.23$ & \crd{0}$-0.00$ \\
    \texttt{hedge\_formal\_rate}                 & 0.000 & 0.000 & 0.000 & 0.000 & \crd{1}$-0.02$ & \crd{3}$-0.12$ & \crd{4}$-0.15$ \\
    \texttt{hedge\_formal\_count}                & 0.000 & 0.000 & 0.000 & 0.000 & \crd{0}$+0.00$ & \crd{5}$-0.19$ & \crd{5}$-0.18$* \\
    \addlinespace[1.5pt]
    \multicolumn{8}{l}{\textit{Syntactic complexity} (5)} \\
    \addlinespace[0.5pt]
    \texttt{tree\_depth\_mean}                   & 4.380 & 5.523 & 4.375 & 4.866 & \cg{32}$+1.27$* & \cg{2}$+0.08$ & \cg{15}$+0.59$* \\
    \texttt{dep\_distance\_mean}                 & 3.116 & 3.517 & 3.232 & 3.502 & \cg{26}$+1.02$* & \cg{9}$+0.37$* & \cg{21}$+0.85$* \\
    \texttt{tree\_depth\_max}                    & 7.000 & 8.000 & 7.000 & 7.000 & \cg{13}$+0.51$* & \cg{2}$+0.07$ & \cg{2}$+0.07$ \\
    \texttt{subordinate\_clause\_rate}           & 1.000 & 1.333 & 0.875 & 1.125 & \cg{12}$+0.48$* & \crd{12}$-0.48$* & \cg{3}$+0.10$ \\
    \texttt{n\_sentences}                        & 10.000 & 8.000 & 10.000 & 8.000 & \crd{11}$-0.45$* & \crd{4}$-0.15$ & \crd{11}$-0.42$* \\
    \addlinespace[1.5pt]
    \multicolumn{8}{l}{\textit{Named entities} (9)} \\
    \addlinespace[0.5pt]
    \texttt{ner\_total\_rate}                    & 0.456 & 0.000 & 0.000 & 0.000 & \crd{13}$-0.53$* & \crd{5}$-0.19$* & \crd{6}$-0.23$* \\
    \texttt{ner\_PRS\_rate}                      & 0.000 & 0.000 & 0.000 & 0.000 & \crd{9}$-0.35$* & \crd{3}$-0.11$ & \crd{6}$-0.23$* \\
    \texttt{ner\_ORG\_rate}                      & 0.000 & 0.000 & 0.000 & 0.000 & \crd{8}$-0.33$* & \crd{7}$-0.27$* & \crd{4}$-0.15$ \\
    \texttt{ner\_LOC\_rate}                      & 0.000 & 0.000 & 0.000 & 0.000 & \crd{8}$-0.30$* & \crd{0}$+0.00$ & \crd{1}$-0.03$ \\
    \texttt{ner\_TME\_rate}                      & 0.000 & 0.000 & 0.000 & 0.000 & \crd{6}$-0.22$* & \crd{6}$-0.23$* & \crd{6}$-0.24$* \\
    \texttt{ner\_WRK\_rate}                      & 0.000 & 0.000 & 0.000 & 0.000 & \crd{3}$-0.10$ & \crd{0}$+0.00$ & \crd{2}$-0.06$ \\
    \texttt{ner\_EVN\_rate}                      & 0.000 & 0.000 & 0.000 & 0.000 & \cg{2}$+0.08$ & \crd{0}$+0.00$ & \cg{2}$+0.08$ \\
    \texttt{ner\_OBJ\_rate}                      & 0.000 & 0.000 & 0.000 & 0.000 & \crd{2}$-0.08$ & \crd{0}$+0.00$ & \crd{2}$-0.08$ \\
    \texttt{ner\_MSR\_rate}                      & 0.000 & 0.000 & 0.000 & 0.000 & \crd{2}$-0.08$ & \crd{3}$-0.10$ & \crd{2}$-0.08$ \\
    \addlinespace[1.5pt]
    \multicolumn{8}{l}{\textit{Readability (textstat)} (7)} \\
    \addlinespace[0.5pt]
    \texttt{textstat\_flesch\_kincaid\_grade}    & 12.408 & 16.431 & 14.353 & 15.121 & \cg{43}$+1.71$* & \cg{19}$+0.75$* & \cg{31}$+1.23$* \\
    \texttt{textstat\_smog\_index}               & 14.030 & 17.267 & 15.903 & 16.322 & \cg{40}$+1.59$* & \cg{20}$+0.80$* & \cg{31}$+1.23$* \\
    \texttt{textstat\_gunning\_fog}              & 15.504 & 19.900 & 18.299 & 18.930 & \cg{39}$+1.54$* & \cg{21}$+0.84$* & \cg{32}$+1.28$* \\
    \texttt{textstat\_flesch\_reading\_ease}     & 38.744 & 21.759 & 24.413 & 20.367 & \crd{35}$-1.41$* & \crd{29}$-1.17$* & \crd{37}$-1.49$* \\
    \texttt{rough\_syllable\_count\_sv}          & 391.500 & 472.500 & 404.000 & 405.500 & \cg{7}$+0.26$* & \cg{2}$+0.09$ & \crd{1}$-0.02$ \\
    \texttt{char\_count}                         & 1282.500 & 1542.000 & 1312.000 & 1336.500 & \cg{7}$+0.26$* & \cg{2}$+0.08$ & \crd{1}$-0.02$ \\
    \texttt{word\_count}                         & 182.000 & 199.500 & 162.500 & 165.000 & \cg{2}$+0.09$* & \crd{5}$-0.21$ & \crd{6}$-0.23$* \\
    \addlinespace[1.5pt]
    \multicolumn{8}{l}{\textit{POS proportions} (17)} \\
    \addlinespace[0.5pt]
    \texttt{pos\_CCONJ}                          & 0.045 & 0.080 & 0.068 & 0.062 & \cg{44}$+1.77$* & \cg{31}$+1.25$* & \cg{23}$+0.90$* \\
    \texttt{pos\_PRON}                           & 0.053 & 0.023 & 0.017 & 0.020 & \crd{30}$-1.18$* & \crd{34}$-1.37$* & \crd{33}$-1.33$* \\
    \texttt{pos\_AUX}                            & 0.043 & 0.024 & 0.030 & 0.026 & \crd{21}$-0.83$* & \crd{16}$-0.62$* & \crd{22}$-0.89$* \\
    \texttt{pos\_DET}                            & 0.052 & 0.035 & 0.038 & 0.032 & \crd{20}$-0.78$* & \crd{16}$-0.62$* & \crd{19}$-0.75$* \\
    \texttt{pos\_ADV}                            & 0.045 & 0.060 & 0.044 & 0.039 & \cg{17}$+0.69$* & \crd{0}$+0.00$ & \crd{6}$-0.24$* \\
    \texttt{pos\_ADP}                            & 0.133 & 0.113 & 0.095 & 0.103 & \crd{17}$-0.68$* & \crd{28}$-1.12$* & \crd{25}$-0.99$* \\
    \texttt{pos\_NUM}                            & 0.008 & 0.000 & 0.000 & 0.000 & \crd{16}$-0.63$* & \crd{15}$-0.61$* & \crd{18}$-0.71$* \\
    \texttt{pos\_PUNCT}                          & 0.077 & 0.097 & 0.114 & 0.153 & \cg{13}$+0.53$* & \cg{27}$+1.08$* & \cg{51}$+2.05$* \\
    \texttt{pos\_NOUN}                           & 0.282 & 0.306 & 0.322 & 0.291 & \cg{12}$+0.49$* & \cg{21}$+0.84$* & \cg{7}$+0.27$* \\
    \texttt{pos\_PROPN}                          & 0.007 & 0.000 & 0.000 & 0.006 & \crd{12}$-0.49$* & \crd{3}$-0.11$ & \cg{3}$+0.13$* \\
    \texttt{pos\_ADJ}                            & 0.085 & 0.092 & 0.106 & 0.096 & \cg{9}$+0.35$* & \cg{20}$+0.80$* & \cg{8}$+0.32$* \\
    \texttt{pos\_PART}                           & 0.022 & 0.018 & 0.020 & 0.022 & \crd{7}$-0.28$* & \crd{4}$-0.16$ & \crd{1}$-0.04$ \\
    \texttt{pos\_SYM}                            & 0.000 & 0.000 & 0.000 & 0.000 & \cg{3}$+0.12$ & \crd{7}$-0.28$* & \crd{2}$-0.09$ \\
    \texttt{pos\_INTJ}                           & 0.000 & 0.000 & 0.000 & 0.000 & \crd{2}$-0.08$ & \crd{0}$+0.00$ & \cg{1}$+0.05$ \\
    \texttt{pos\_VERB}                           & 0.109 & 0.108 & 0.103 & 0.097 & \crd{2}$-0.07$ & \crd{7}$-0.29$* & \crd{13}$-0.50$* \\
    \texttt{pos\_SCONJ}                          & 0.023 & 0.023 & 0.021 & 0.012 & \crd{0}$+0.00$ & \crd{5}$-0.18$ & \crd{17}$-0.68$* \\
    \texttt{pos\_X}                              & 0.000 & 0.000 & 0.000 & 0.000 & \crd{0}$+0.00$ & \crd{0}$+0.00$ & \crd{0}$+0.00$ \\
    \bottomrule
  \end{tabular}
  \caption{Formal register, Baseline condition: every feature tested, human vs.\
  each generator (paired Wilcoxon signed-rank, Cohen's $d_z$, Benjamini--Hochberg FDR within each
  generator's family of 63 features). Positive (green) = generator $>$ human, negative (red) =
  human $>$ generator; shading $\propto |d_z|$, * denotes $q<0.05$. GPT-5.2 and Mistral-3.2 use the
  full corpus (170 formal / \textasciitilde{}1{,}150 informal); GPT-5.6 is a 100-document quiz
  subset, so its column has less statistical power. Features are grouped by extractor and ordered by
  $|d_z|$ vs.\ GPT-5.2 at baseline, held fixed across this register's three condition-tables so a row
  tracks the same feature. Source: \texttt{model\_comparison\_three\_way\_results.csv}.}
  \label{app:full-three-way-formal-baseline}
\end{table}

\clearpage

% ---- formal register, Human-like condition: all features, human vs 3 models (app:full-three-way-formal-human_like) ----
% REQUIRES: booktabs + the \cg / \crd shading macros from the preamble.
% 63 features x 3 models.
% Sized to fit one page at \scriptsize with arraystretch 0.9; if it still overflows, wrap in longtable.
\begin{table}[htbp]
  \centering
  \scriptsize
  \setlength{\tabcolsep}{2.5pt}
  \renewcommand{\arraystretch}{0.9}
  \begin{tabular}{l rrrr ccc}
    \toprule
    & \multicolumn{4}{c}{Median} & \multicolumn{3}{c}{Cohen's $d_z$ (vs.\ human)} \\
    \cmidrule(lr){2-5}\cmidrule(l){6-8}
    Feature & H & GPT-5.2 & GPT-5.6 & Mistral & GPT-5.2 & GPT-5.6 & Mistral \\
    \midrule
    \multicolumn{8}{l}{\textit{Stylometric surface} (11)} \\
    \addlinespace[0.5pt]
    \texttt{punct\_comma}                        & 1.227 & 4.285 & 5.264 & 4.285 & \cg{39}$+1.55$* & \cg{49}$+1.96$* & \cg{33}$+1.30$* \\
    \texttt{mtld}                                & 100.081 & 155.790 & 151.157 & 125.044 & \cg{38}$+1.51$* & \cg{38}$+1.53$* & \cg{20}$+0.78$* \\
    \texttt{mattr}                               & 0.827 & 0.864 & 0.875 & 0.860 & \cg{24}$+0.97$* & \cg{29}$+1.15$* & \cg{20}$+0.79$* \\
    \texttt{punct\_period}                       & 4.680 & 3.599 & 4.811 & 4.130 & \crd{24}$-0.97$* & \cg{4}$+0.16$ & \crd{11}$-0.43$* \\
    \texttt{function\_word\_rate}                & 0.424 & 0.398 & 0.365 & 0.399 & \crd{13}$-0.51$* & \crd{29}$-1.16$* & \crd{12}$-0.49$* \\
    \texttt{bigram\_repetition\_rate}            & 0.046 & 0.037 & 0.026 & 0.046 & \crd{10}$-0.40$* & \crd{19}$-0.77$* & \crd{5}$-0.19$* \\
    \texttt{trigram\_repetition\_rate}           & 0.013 & 0.006 & 0.004 & 0.007 & \crd{13}$-0.53$* & \crd{17}$-0.69$* & \crd{9}$-0.36$* \\
    \texttt{ttr}                                 & 0.632 & 0.618 & 0.682 & 0.645 & \crd{7}$-0.29$* & \cg{11}$+0.45$* & \cg{0}$+0.01$ \\
    \texttt{punct\_dash}                         & 0.000 & 0.258 & 0.000 & 0.000 & \cg{18}$+0.71$* & \crd{1}$-0.03$ & \cg{13}$+0.52$* \\
    \texttt{punct\_semicolon}                    & 0.000 & 0.000 & 0.000 & 0.000 & \crd{11}$-0.44$* & \crd{11}$-0.43$* & \crd{12}$-0.46$* \\
    \texttt{punct\_colon}                        & 0.000 & 0.334 & 0.437 & 0.000 & \crd{0}$-0.01$* & \crd{1}$-0.02$* & \crd{4}$-0.17$ \\
    \addlinespace[1.5pt]
    \multicolumn{8}{l}{\textit{Pragmatic markers} (14)} \\
    \addlinespace[0.5pt]
    \texttt{hedge\_neutral\_rate}                & 0.712 & 1.660 & 1.539 & 1.695 & \cg{22}$+0.89$* & \cg{20}$+0.79$* & \cg{22}$+0.86$* \\
    \texttt{hedge\_neutral\_count}               & 1.000 & 5.000 & 3.000 & 3.000 & \cg{31}$+1.25$* & \cg{22}$+0.88$* & \cg{20}$+0.81$* \\
    \texttt{hedge\_rate}                         & 0.815 & 1.939 & 1.744 & 1.932 & \cg{23}$+0.92$* & \cg{17}$+0.69$* & \cg{21}$+0.84$* \\
    \texttt{hedge\_count}                        & 1.000 & 6.000 & 3.000 & 4.000 & \cg{32}$+1.29$* & \cg{19}$+0.77$* & \cg{19}$+0.75$* \\
    \texttt{epistemic\_rate}                     & 0.000 & 0.000 & 0.000 & 0.000 & \crd{10}$-0.38$* & \crd{8}$-0.31$* & \crd{7}$-0.28$* \\
    \texttt{epistemic\_count}                    & 0.000 & 0.000 & 0.000 & 0.000 & \crd{4}$-0.16$ & \crd{7}$-0.27$* & \crd{6}$-0.23$* \\
    \texttt{negation\_rate}                      & 0.347 & 0.504 & 0.487 & 0.149 & \cg{2}$+0.08$* & \cg{5}$+0.19$ & \crd{4}$-0.15$ \\
    \texttt{negation\_count}                     & 1.000 & 2.000 & 1.000 & 0.500 & \cg{13}$+0.51$* & \cg{5}$+0.21$* & \crd{3}$-0.10$ \\
    \texttt{hedge\_informal\_count}              & 0.000 & 0.000 & 0.000 & 0.000 & \cg{9}$+0.34$* & \crd{2}$-0.06$ & \cg{7}$+0.29$* \\
    \texttt{n\_alpha\_tokens}                    & 178.500 & 333.000 & 202.500 & 200.000 & \cg{29}$+1.16$* & \cg{10}$+0.38$* & \cg{3}$+0.12$* \\
    \texttt{negation\_sentence\_rate}            & 0.051 & 0.118 & 0.083 & 0.031 & \cg{9}$+0.37$* & \cg{5}$+0.21$* & \crd{0}$-0.01$ \\
    \texttt{hedge\_informal\_rate}               & 0.000 & 0.000 & 0.000 & 0.000 & \cg{6}$+0.23$* & \crd{1}$-0.04$ & \cg{9}$+0.37$* \\
    \texttt{hedge\_formal\_rate}                 & 0.000 & 0.000 & 0.000 & 0.000 & \crd{0}$-0.00$ & \crd{4}$-0.15$ & \crd{6}$-0.22$* \\
    \texttt{hedge\_formal\_count}                & 0.000 & 0.000 & 0.000 & 0.000 & \cg{4}$+0.14$ & \crd{3}$-0.13$ & \crd{6}$-0.24$* \\
    \addlinespace[1.5pt]
    \multicolumn{8}{l}{\textit{Syntactic complexity} (5)} \\
    \addlinespace[0.5pt]
    \texttt{tree\_depth\_mean}                   & 4.380 & 5.414 & 4.364 & 5.000 & \cg{31}$+1.23$* & \cg{3}$+0.11$ & \cg{19}$+0.75$* \\
    \texttt{dep\_distance\_mean}                 & 3.116 & 3.535 & 3.240 & 3.485 & \cg{35}$+1.38$* & \cg{10}$+0.41$* & \cg{22}$+0.86$* \\
    \texttt{tree\_depth\_max}                    & 7.000 & 8.000 & 8.000 & 7.000 & \cg{19}$+0.75$* & \cg{8}$+0.33$* & \cg{5}$+0.21$* \\
    \texttt{subordinate\_clause\_rate}           & 1.000 & 1.452 & 1.000 & 1.333 & \cg{18}$+0.72$* & \crd{7}$-0.27$ & \cg{14}$+0.54$* \\
    \texttt{n\_sentences}                        & 10.000 & 14.000 & 12.000 & 10.000 & \cg{9}$+0.36$* & \cg{7}$+0.29$* & \crd{6}$-0.25$* \\
    \addlinespace[1.5pt]
    \multicolumn{8}{l}{\textit{Named entities} (9)} \\
    \addlinespace[0.5pt]
    \texttt{ner\_total\_rate}                    & 0.456 & 0.000 & 0.000 & 0.000 & \crd{14}$-0.55$* & \crd{6}$-0.23$* & \crd{9}$-0.37$* \\
    \texttt{ner\_PRS\_rate}                      & 0.000 & 0.000 & 0.000 & 0.000 & \crd{8}$-0.32$* & \crd{3}$-0.13$ & \crd{5}$-0.21$* \\
    \texttt{ner\_ORG\_rate}                      & 0.000 & 0.000 & 0.000 & 0.000 & \crd{9}$-0.37$* & \crd{9}$-0.35$* & \crd{6}$-0.22$* \\
    \texttt{ner\_LOC\_rate}                      & 0.000 & 0.000 & 0.000 & 0.000 & \crd{8}$-0.33$* & \crd{1}$-0.05$ & \crd{6}$-0.23$* \\
    \texttt{ner\_TME\_rate}                      & 0.000 & 0.000 & 0.000 & 0.000 & \crd{5}$-0.21$* & \crd{4}$-0.16$ & \crd{5}$-0.20$* \\
    \texttt{ner\_WRK\_rate}                      & 0.000 & 0.000 & 0.000 & 0.000 & \crd{3}$-0.11$ & \crd{0}$-0.00$ & \crd{3}$-0.13$ \\
    \texttt{ner\_EVN\_rate}                      & 0.000 & 0.000 & 0.000 & 0.000 & \crd{0}$+0.00$ & \crd{0}$+0.00$ & \cg{2}$+0.08$ \\
    \texttt{ner\_OBJ\_rate}                      & 0.000 & 0.000 & 0.000 & 0.000 & \crd{2}$-0.08$ & \crd{0}$+0.00$ & \crd{2}$-0.08$ \\
    \texttt{ner\_MSR\_rate}                      & 0.000 & 0.000 & 0.000 & 0.000 & \crd{2}$-0.08$ & \crd{3}$-0.10$ & \crd{2}$-0.08$ \\
    \addlinespace[1.5pt]
    \multicolumn{8}{l}{\textit{Readability (textstat)} (7)} \\
    \addlinespace[0.5pt]
    \texttt{textstat\_flesch\_kincaid\_grade}    & 12.408 & 15.511 & 13.909 & 14.599 & \cg{30}$+1.18$* & \cg{14}$+0.57$* & \cg{23}$+0.90$* \\
    \texttt{textstat\_smog\_index}               & 14.030 & 16.306 & 15.315 & 15.740 & \cg{27}$+1.06$* & \cg{16}$+0.64$* & \cg{21}$+0.85$* \\
    \texttt{textstat\_gunning\_fog}              & 15.504 & 18.615 & 17.361 & 17.821 & \cg{26}$+1.03$* & \cg{16}$+0.64$* & \cg{20}$+0.81$* \\
    \texttt{textstat\_flesch\_reading\_ease}     & 38.744 & 28.616 & 28.966 & 27.985 & \crd{19}$-0.75$* & \crd{24}$-0.94$* & \crd{22}$-0.87$* \\
    \texttt{rough\_syllable\_count\_sv}          & 391.500 & 751.500 & 487.000 & 462.000 & \cg{32}$+1.26$* & \cg{16}$+0.62$* & \cg{6}$+0.24$* \\
    \texttt{char\_count}                         & 1282.500 & 2448.500 & 1595.500 & 1514.000 & \cg{32}$+1.27$* & \cg{16}$+0.62$* & \cg{6}$+0.23$* \\
    \texttt{word\_count}                         & 182.000 & 333.500 & 203.000 & 202.000 & \cg{29}$+1.16$* & \cg{9}$+0.36$* & \cg{3}$+0.11$* \\
    \addlinespace[1.5pt]
    \multicolumn{8}{l}{\textit{POS proportions} (17)} \\
    \addlinespace[0.5pt]
    \texttt{pos\_CCONJ}                          & 0.045 & 0.068 & 0.062 & 0.065 & \cg{33}$+1.30$* & \cg{27}$+1.08$* & \cg{24}$+0.96$* \\
    \texttt{pos\_PRON}                           & 0.053 & 0.029 & 0.023 & 0.030 & \crd{23}$-0.93$* & \crd{30}$-1.18$* & \crd{22}$-0.88$* \\
    \texttt{pos\_AUX}                            & 0.043 & 0.029 & 0.035 & 0.034 & \crd{15}$-0.60$* & \crd{11}$-0.44$* & \crd{12}$-0.48$* \\
    \texttt{pos\_DET}                            & 0.052 & 0.041 & 0.040 & 0.042 & \crd{13}$-0.51$* & \crd{16}$-0.63$* & \crd{11}$-0.45$* \\
    \texttt{pos\_ADV}                            & 0.045 & 0.070 & 0.048 & 0.050 & \cg{31}$+1.24$* & \cg{6}$+0.23$* & \cg{8}$+0.30$* \\
    \texttt{pos\_ADP}                            & 0.133 & 0.109 & 0.096 & 0.104 & \crd{22}$-0.89$* & \crd{30}$-1.19$* & \crd{28}$-1.12$* \\
    \texttt{pos\_NUM}                            & 0.008 & 0.000 & 0.000 & 0.000 & \crd{16}$-0.64$* & \crd{17}$-0.68$* & \crd{18}$-0.73$* \\
    \texttt{pos\_PUNCT}                          & 0.077 & 0.094 & 0.108 & 0.119 & \cg{13}$+0.53$* & \cg{27}$+1.07$* & \cg{30}$+1.18$* \\
    \texttt{pos\_NOUN}                           & 0.282 & 0.282 & 0.310 & 0.268 & \crd{1}$-0.02$ & \cg{15}$+0.60$* & \crd{5}$-0.20$* \\
    \texttt{pos\_PROPN}                          & 0.007 & 0.000 & 0.004 & 0.004 & \crd{14}$-0.54$* & \crd{5}$-0.21$* & \crd{6}$-0.24$* \\
    \texttt{pos\_ADJ}                            & 0.085 & 0.100 & 0.108 & 0.099 & \cg{13}$+0.50$* & \cg{19}$+0.74$* & \cg{13}$+0.52$* \\
    \texttt{pos\_PART}                           & 0.022 & 0.022 & 0.022 & 0.025 & \crd{2}$-0.07$ & \crd{3}$-0.12$ & \cg{6}$+0.23$* \\
    \texttt{pos\_SYM}                            & 0.000 & 0.000 & 0.000 & 0.000 & \crd{1}$-0.02$ & \crd{7}$-0.28$* & \crd{5}$-0.20$* \\
    \texttt{pos\_INTJ}                           & 0.000 & 0.000 & 0.000 & 0.000 & \cg{1}$+0.04$ & \crd{0}$+0.00$ & \crd{2}$-0.08$ \\
    \texttt{pos\_VERB}                           & 0.109 & 0.110 & 0.106 & 0.105 & \crd{1}$-0.02$ & \crd{4}$-0.17$ & \crd{4}$-0.14$ \\
    \texttt{pos\_SCONJ}                          & 0.023 & 0.031 & 0.027 & 0.016 & \cg{13}$+0.52$* & \cg{6}$+0.24$* & \crd{10}$-0.38$* \\
    \texttt{pos\_X}                              & 0.000 & 0.000 & 0.000 & 0.000 & \crd{0}$+0.00$ & \crd{0}$+0.00$ & \crd{0}$+0.00$ \\
    \bottomrule
  \end{tabular}
  \caption{Formal register, Human-like condition: every feature tested, human vs.\
  each generator (paired Wilcoxon signed-rank, Cohen's $d_z$, Benjamini--Hochberg FDR within each
  generator's family of 63 features). Positive (green) = generator $>$ human, negative (red) =
  human $>$ generator; shading $\propto |d_z|$, * denotes $q<0.05$. GPT-5.2 and Mistral-3.2 use the
  full corpus (170 formal / \textasciitilde{}1{,}150 informal); GPT-5.6 is a 100-document quiz
  subset, so its column has less statistical power. Features are grouped by extractor and ordered by
  $|d_z|$ vs.\ GPT-5.2 at baseline, held fixed across this register's three condition-tables so a row
  tracks the same feature.}
  \label{app:full-three-way-formal-human_like}
\end{table}

\clearpage

% ---- formal register, Detector-evasive condition: all features, human vs 3 models (app:full-three-way-formal-detector_evasive) ----
% REQUIRES: booktabs + the \cg / \crd shading macros from the preamble.
% 63 features x 3 models.
% Sized to fit one page at \scriptsize with arraystretch 0.9; if it still overflows, wrap in longtable.
\begin{table}[htbp]
  \centering
  \scriptsize
  \setlength{\tabcolsep}{2.5pt}
  \renewcommand{\arraystretch}{0.9}
  \begin{tabular}{l rrrr ccc}
    \toprule
    & \multicolumn{4}{c}{Median} & \multicolumn{3}{c}{Cohen's $d_z$ (vs.\ human)} \\
    \cmidrule(lr){2-5}\cmidrule(l){6-8}
    Feature & H & GPT-5.2 & GPT-5.6 & Mistral & GPT-5.2 & GPT-5.6 & Mistral \\
    \midrule
    \multicolumn{8}{l}{\textit{Stylometric surface} (11)} \\
    \addlinespace[0.5pt]
    \texttt{punct\_comma}                        & 1.227 & 0.774 & 3.247 & 0.594 & \crd{7}$-0.27$* & \cg{26}$+1.03$* & \crd{6}$-0.22$* \\
    \texttt{mtld}                                & 100.081 & 76.035 & 107.378 & 88.238 & \crd{21}$-0.83$* & \cg{7}$+0.28$* & \crd{8}$-0.33$* \\
    \texttt{mattr}                               & 0.827 & 0.790 & 0.856 & 0.830 & \crd{16}$-0.64$* & \cg{19}$+0.77$* & \cg{2}$+0.06$ \\
    \texttt{punct\_period}                       & 4.680 & 8.962 & 7.334 & 8.965 & \cg{58}$+2.32$* & \cg{50}$+2.00$* & \cg{45}$+1.80$* \\
    \texttt{function\_word\_rate}                & 0.424 & 0.435 & 0.349 & 0.364 & \cg{3}$+0.10$ & \crd{34}$-1.35$* & \crd{27}$-1.09$* \\
    \texttt{bigram\_repetition\_rate}            & 0.046 & 0.121 & 0.048 & 0.077 & \cg{33}$+1.31$* & \crd{3}$-0.11$ & \cg{11}$+0.44$* \\
    \texttt{trigram\_repetition\_rate}           & 0.013 & 0.043 & 0.011 & 0.020 & \cg{22}$+0.88$* & \crd{3}$-0.13$ & \cg{8}$+0.30$* \\
    \texttt{ttr}                                 & 0.632 & 0.468 & 0.631 & 0.625 & \crd{42}$-1.67$* & \crd{5}$-0.20$ & \crd{5}$-0.19$* \\
    \texttt{punct\_dash}                         & 0.000 & 0.000 & 0.000 & 0.000 & \crd{4}$-0.16$ & \cg{0}$+0.01$ & \crd{2}$-0.09$ \\
    \texttt{punct\_semicolon}                    & 0.000 & 0.000 & 0.000 & 0.000 & \crd{11}$-0.44$* & \crd{5}$-0.18$ & \crd{11}$-0.42$* \\
    \texttt{punct\_colon}                        & 0.000 & 0.000 & 0.480 & 0.000 & \crd{7}$-0.29$* & \cg{2}$+0.09$* & \crd{7}$-0.29$* \\
    \addlinespace[1.5pt]
    \multicolumn{8}{l}{\textit{Pragmatic markers} (14)} \\
    \addlinespace[0.5pt]
    \texttt{hedge\_neutral\_rate}                & 0.712 & 0.858 & 0.583 & 1.354 & \cg{6}$+0.23$* & \cg{3}$+0.13$ & \cg{14}$+0.57$* \\
    \texttt{hedge\_neutral\_count}               & 1.000 & 3.000 & 1.000 & 2.000 & \cg{16}$+0.63$* & \cg{4}$+0.16$ & \cg{8}$+0.32$* \\
    \texttt{hedge\_rate}                         & 0.815 & 1.065 & 0.939 & 1.455 & \cg{7}$+0.26$* & \cg{1}$+0.05$ & \cg{13}$+0.50$* \\
    \texttt{hedge\_count}                        & 1.000 & 4.000 & 2.000 & 2.000 & \cg{17}$+0.68$* & \cg{2}$+0.08$ & \cg{6}$+0.22$* \\
    \texttt{epistemic\_rate}                     & 0.000 & 0.000 & 0.000 & 0.000 & \crd{12}$-0.48$* & \crd{9}$-0.35$* & \crd{8}$-0.30$* \\
    \texttt{epistemic\_count}                    & 0.000 & 0.000 & 0.000 & 0.000 & \crd{7}$-0.27$* & \crd{8}$-0.30$* & \crd{9}$-0.35$* \\
    \texttt{negation\_rate}                      & 0.347 & 0.286 & 0.000 & 0.000 & \crd{5}$-0.18$ & \crd{2}$-0.06$ & \crd{7}$-0.27$* \\
    \texttt{negation\_count}                     & 1.000 & 1.000 & 0.000 & 0.000 & \cg{7}$+0.26$* & \crd{3}$-0.11$ & \crd{9}$-0.35$* \\
    \texttt{hedge\_informal\_count}              & 0.000 & 0.000 & 0.000 & 0.000 & \cg{10}$+0.41$* & \crd{3}$-0.12$ & \crd{4}$-0.14$ \\
    \texttt{n\_alpha\_tokens}                    & 178.500 & 398.500 & 186.500 & 147.500 & \cg{34}$+1.37$* & \cg{5}$+0.20$* & \crd{9}$-0.36$* \\
    \texttt{negation\_sentence\_rate}            & 0.051 & 0.030 & 0.000 & 0.000 & \crd{10}$-0.39$* & \crd{7}$-0.26$* & \crd{13}$-0.51$* \\
    \texttt{hedge\_informal\_rate}               & 0.000 & 0.000 & 0.000 & 0.000 & \cg{6}$+0.24$* & \crd{2}$-0.09$ & \crd{1}$-0.03$ \\
    \texttt{hedge\_formal\_rate}                 & 0.000 & 0.000 & 0.000 & 0.000 & \crd{3}$-0.10$ & \crd{5}$-0.21$ & \crd{4}$-0.16$ \\
    \texttt{hedge\_formal\_count}                & 0.000 & 0.000 & 0.000 & 0.000 & \cg{1}$+0.03$ & \crd{5}$-0.19$ & \crd{5}$-0.19$* \\
    \addlinespace[1.5pt]
    \multicolumn{8}{l}{\textit{Syntactic complexity} (5)} \\
    \addlinespace[0.5pt]
    \texttt{tree\_depth\_mean}                   & 4.380 & 3.144 & 3.382 & 3.197 & \crd{33}$-1.33$* & \crd{27}$-1.07$* & \crd{28}$-1.11$* \\
    \texttt{dep\_distance\_mean}                 & 3.116 & 2.677 & 2.778 & 2.687 & \crd{29}$-1.17$* & \crd{26}$-1.05$* & \crd{21}$-0.82$* \\
    \texttt{tree\_depth\_max}                    & 7.000 & 6.000 & 6.000 & 5.000 & \crd{10}$-0.39$* & \crd{15}$-0.59$* & \crd{26}$-1.04$* \\
    \texttt{subordinate\_clause\_rate}           & 1.000 & 0.650 & 0.525 & 0.500 & \crd{22}$-0.89$* & \crd{27}$-1.07$* & \crd{28}$-1.14$* \\
    \texttt{n\_sentences}                        & 10.000 & 41.000 & 18.000 & 16.000 & \cg{49}$+1.95$* & \cg{29}$+1.16$* & \cg{14}$+0.58$* \\
    \addlinespace[1.5pt]
    \multicolumn{8}{l}{\textit{Named entities} (9)} \\
    \addlinespace[0.5pt]
    \texttt{ner\_total\_rate}                    & 0.456 & 0.437 & 0.000 & 0.000 & \crd{1}$-0.04$ & \cg{3}$+0.13$ & \cg{1}$+0.04$ \\
    \texttt{ner\_PRS\_rate}                      & 0.000 & 0.000 & 0.000 & 0.000 & \cg{1}$+0.04$ & \cg{2}$+0.08$ & \cg{2}$+0.06$ \\
    \texttt{ner\_ORG\_rate}                      & 0.000 & 0.000 & 0.000 & 0.000 & \crd{1}$-0.04$ & \cg{1}$+0.04$ & \cg{0}$+0.01$ \\
    \texttt{ner\_LOC\_rate}                      & 0.000 & 0.000 & 0.000 & 0.000 & \crd{3}$-0.10$ & \cg{3}$+0.12$ & \cg{1}$+0.04$ \\
    \texttt{ner\_TME\_rate}                      & 0.000 & 0.000 & 0.000 & 0.000 & \crd{2}$-0.06$ & \cg{2}$+0.09$ & \crd{2}$-0.09$ \\
    \texttt{ner\_WRK\_rate}                      & 0.000 & 0.000 & 0.000 & 0.000 & \cg{1}$+0.05$ & \cg{3}$+0.11$ & \crd{3}$-0.11$ \\
    \texttt{ner\_EVN\_rate}                      & 0.000 & 0.000 & 0.000 & 0.000 & \cg{2}$+0.08$ & \crd{0}$+0.00$ & \cg{2}$+0.08$ \\
    \texttt{ner\_OBJ\_rate}                      & 0.000 & 0.000 & 0.000 & 0.000 & \crd{2}$-0.06$ & \crd{0}$+0.00$ & \crd{2}$-0.08$ \\
    \texttt{ner\_MSR\_rate}                      & 0.000 & 0.000 & 0.000 & 0.000 & \crd{2}$-0.08$ & \crd{3}$-0.10$ & \crd{2}$-0.08$ \\
    \addlinespace[1.5pt]
    \multicolumn{8}{l}{\textit{Readability (textstat)} (7)} \\
    \addlinespace[0.5pt]
    \texttt{textstat\_flesch\_kincaid\_grade}    & 12.408 & 6.863 & 9.763 & 10.322 & \crd{61}$-2.43$* & \crd{31}$-1.24$* & \crd{20}$-0.81$* \\
    \texttt{textstat\_smog\_index}               & 14.030 & 9.478 & 11.740 & 11.851 & \crd{63}$-2.53$* & \crd{35}$-1.38$* & \crd{28}$-1.14$* \\
    \texttt{textstat\_gunning\_fog}              & 15.504 & 8.987 & 12.552 & 13.827 & \crd{63}$-2.50$* & \crd{31}$-1.23$* & \crd{15}$-0.59$* \\
    \texttt{textstat\_flesch\_reading\_ease}     & 38.744 & 63.827 & 46.185 & 38.863 & \cg{50}$+2.01$* & \cg{14}$+0.58$* & \crd{3}$-0.10$ \\
    \texttt{rough\_syllable\_count\_sv}          & 391.500 & 742.500 & 410.000 & 337.000 & \cg{27}$+1.09$* & \cg{5}$+0.20$* & \crd{7}$-0.29$* \\
    \texttt{char\_count}                         & 1282.500 & 2551.000 & 1373.000 & 1137.500 & \cg{29}$+1.17$* & \cg{6}$+0.25$* & \crd{7}$-0.27$* \\
    \texttt{word\_count}                         & 182.000 & 400.000 & 188.000 & 148.000 & \cg{34}$+1.36$* & \cg{5}$+0.20$* & \crd{9}$-0.37$* \\
    \addlinespace[1.5pt]
    \multicolumn{8}{l}{\textit{POS proportions} (17)} \\
    \addlinespace[0.5pt]
    \texttt{pos\_CCONJ}                          & 0.045 & 0.044 & 0.056 & 0.041 & \crd{1}$-0.05$ & \cg{17}$+0.69$* & \crd{4}$-0.15$ \\
    \texttt{pos\_PRON}                           & 0.053 & 0.080 & 0.029 & 0.031 & \cg{19}$+0.74$* & \crd{25}$-1.01$* & \crd{19}$-0.75$* \\
    \texttt{pos\_AUX}                            & 0.043 & 0.030 & 0.025 & 0.043 & \crd{14}$-0.56$* & \crd{17}$-0.69$* & \cg{2}$+0.06$ \\
    \texttt{pos\_DET}                            & 0.052 & 0.031 & 0.038 & 0.029 & \crd{21}$-0.85$* & \crd{20}$-0.78$* & \crd{23}$-0.93$* \\
    \texttt{pos\_ADV}                            & 0.045 & 0.075 & 0.050 & 0.034 & \cg{31}$+1.23$* & \cg{7}$+0.26$* & \crd{8}$-0.31$* \\
    \texttt{pos\_ADP}                            & 0.133 & 0.096 & 0.091 & 0.090 & \crd{31}$-1.22$* & \crd{34}$-1.34$* & \crd{33}$-1.33$* \\
    \texttt{pos\_NUM}                            & 0.008 & 0.003 & 0.000 & 0.000 & \crd{9}$-0.37$* & \crd{10}$-0.41$* & \crd{15}$-0.60$* \\
    \texttt{pos\_PUNCT}                          & 0.077 & 0.103 & 0.116 & 0.136 & \cg{20}$+0.80$* & \cg{36}$+1.45$* & \cg{29}$+1.15$* \\
    \texttt{pos\_NOUN}                           & 0.282 & 0.253 & 0.323 & 0.281 & \crd{15}$-0.60$* & \cg{24}$+0.94$* & \cg{1}$+0.02$ \\
    \texttt{pos\_PROPN}                          & 0.007 & 0.006 & 0.008 & 0.006 & \crd{0}$+0.00$ & \cg{6}$+0.23$* & \cg{4}$+0.15$ \\
    \texttt{pos\_ADJ}                            & 0.085 & 0.059 & 0.086 & 0.083 & \crd{22}$-0.88$* & \crd{1}$-0.04$ & \crd{1}$-0.03$ \\
    \texttt{pos\_PART}                           & 0.022 & 0.013 & 0.015 & 0.019 & \crd{18}$-0.70$* & \crd{11}$-0.43$* & \crd{5}$-0.21$* \\
    \texttt{pos\_SYM}                            & 0.000 & 0.000 & 0.000 & 0.000 & \crd{3}$-0.13$* & \crd{5}$-0.19$ & \crd{5}$-0.20$* \\
    \texttt{pos\_INTJ}                           & 0.000 & 0.000 & 0.000 & 0.000 & \cg{4}$+0.14$ & \cg{3}$+0.10$ & \cg{2}$+0.06$ \\
    \texttt{pos\_VERB}                           & 0.109 & 0.152 & 0.124 & 0.122 & \cg{37}$+1.49$* & \cg{11}$+0.43$* & \cg{10}$+0.40$* \\
    \texttt{pos\_SCONJ}                          & 0.023 & 0.033 & 0.019 & 0.016 & \cg{14}$+0.56$* & \crd{6}$-0.22$ & \crd{10}$-0.41$* \\
    \texttt{pos\_X}                              & 0.000 & 0.000 & 0.000 & 0.000 & \crd{0}$+0.00$ & \crd{0}$+0.00$ & \crd{0}$+0.00$ \\
    \bottomrule
  \end{tabular}
  \caption{Formal register, Detector-evasive condition: every feature tested, human vs.\
  each generator (paired Wilcoxon signed-rank, Cohen's $d_z$, Benjamini--Hochberg FDR within each
  generator's family of 63 features). Positive (green) = generator $>$ human, negative (red) =
  human $>$ generator; shading $\propto |d_z|$, * denotes $q<0.05$. GPT-5.2 and Mistral-3.2 use the
  full corpus (170 formal / \textasciitilde{}1{,}150 informal); GPT-5.6 is a 100-document quiz
  subset, so its column has less statistical power. Features are grouped by extractor and ordered by
  $|d_z|$ vs.\ GPT-5.2 at baseline, held fixed across this register's three condition-tables so a row
  tracks the same feature.}
  \label{app:full-three-way-formal-detector_evasive}
\end{table}

\clearpage

% ---- informal register, Baseline condition: all features, human vs 3 models (app:full-three-way-informal-baseline) ----
% REQUIRES: booktabs + the \cg / \crd shading macros from the preamble.
% 70 features x 3 models.
% Sized to fit one page at \scriptsize with arraystretch 0.9; if it still overflows, wrap in longtable.
\begin{table}[htbp]
  \centering
  \scriptsize
  \setlength{\tabcolsep}{2.5pt}
  \renewcommand{\arraystretch}{0.9}
  \begin{tabular}{l rrrr ccc}
    \toprule
    & \multicolumn{4}{c}{Median} & \multicolumn{3}{c}{Cohen's $d_z$ (vs.\ human)} \\
    \cmidrule(lr){2-5}\cmidrule(l){6-8}
    Feature & H & GPT-5.2 & GPT-5.6 & Mistral & GPT-5.2 & GPT-5.6 & Mistral \\
    \midrule
    \multicolumn{8}{l}{\textit{Stylometric surface} (14)} \\
    \addlinespace[0.5pt]
    \texttt{ttr}                                 & 0.882 & 0.741 & 0.903 & 0.697 & \crd{19}$-0.76$* & \crd{0}$+0.00$ & \crd{26}$-1.04$* \\
    \texttt{mattr}                               & 0.889 & 0.832 & 0.904 & 0.788 & \crd{11}$-0.44$* & \crd{1}$-0.03$ & \crd{20}$-0.80$* \\
    \texttt{mtld}                                & 71.680 & 93.949 & 97.564 & 73.173 & \cg{8}$+0.33$* & \cg{10}$+0.41$* & \crd{1}$-0.02$ \\
    \texttt{punct\_dash}                         & 0.000 & 0.000 & 0.000 & 0.000 & \cg{8}$+0.31$* & \cg{17}$+0.67$* & \cg{14}$+0.56$* \\
    \texttt{punct\_period}                       & 6.000 & 4.046 & 5.184 & 4.615 & \crd{7}$-0.29$* & \crd{7}$-0.29$* & \crd{6}$-0.25$* \\
    \texttt{bigram\_repetition\_rate}            & 0.000 & 0.014 & 0.000 & 0.037 & \cg{7}$+0.28$* & \crd{5}$-0.20$ & \cg{19}$+0.76$* \\
    \texttt{punct\_question}                     & 0.000 & 0.000 & 0.000 & 0.000 & \crd{7}$-0.26$* & \crd{3}$-0.11$ & \crd{6}$-0.22$* \\
    \texttt{punct\_exclaim}                      & 0.000 & 1.370 & 0.000 & 0.000 & \cg{6}$+0.23$* & \cg{7}$+0.28$* & \crd{0}$-0.01$* \\
    \texttt{punct\_comma}                        & 1.770 & 3.306 & 3.871 & 3.509 & \cg{5}$+0.18$* & \cg{10}$+0.39$* & \cg{5}$+0.21$* \\
    \texttt{punct\_ellipsis}                     & 0.000 & 0.000 & 0.000 & 0.000 & \crd{4}$-0.16$* & \crd{5}$-0.19$ & \crd{4}$-0.14$* \\
    \texttt{function\_word\_rate}                & 0.500 & 0.509 & 0.426 & 0.568 & \cg{4}$+0.15$* & \crd{7}$-0.28$* & \cg{14}$+0.54$* \\
    \texttt{punct\_colon}                        & 0.000 & 0.000 & 0.000 & 0.000 & \crd{3}$-0.13$* & \crd{5}$-0.19$ & \crd{4}$-0.16$* \\
    \texttt{punct\_semicolon}                    & 0.000 & 0.000 & 0.000 & 0.000 & \crd{2}$-0.09$ & \crd{2}$-0.07$ & \crd{3}$-0.11$* \\
    \texttt{trigram\_repetition\_rate}           & 0.000 & 0.000 & 0.000 & 0.000 & \crd{0}$-0.01$ & \crd{2}$-0.07$ & \cg{11}$+0.42$* \\
    \addlinespace[1.5pt]
    \multicolumn{8}{l}{\textit{Pragmatic markers} (18)} \\
    \addlinespace[0.5pt]
    \texttt{hedge\_neutral\_count}               & 1.000 & 3.000 & 1.000 & 4.000 & \cg{23}$+0.90$* & \cg{1}$+0.04$ & \cg{27}$+1.07$* \\
    \texttt{hedge\_count}                        & 1.000 & 4.000 & 1.000 & 6.000 & \cg{22}$+0.87$* & \cg{2}$+0.09$ & \cg{27}$+1.07$* \\
    \texttt{n\_alpha\_tokens}                    & 30.000 & 90.000 & 36.500 & 86.000 & \cg{21}$+0.83$* & \cg{6}$+0.25$* & \cg{20}$+0.81$* \\
    \texttt{booster\_count}                      & 0.000 & 1.000 & 0.000 & 0.000 & \cg{13}$+0.51$* & \cg{5}$+0.18$ & \cg{9}$+0.34$* \\
    \texttt{hedge\_neutral\_rate}                & 1.124 & 3.774 & 2.020 & 5.051 & \cg{11}$+0.43$* & \cg{0}$+0.01$ & \cg{16}$+0.65$* \\
    \texttt{negation\_sentence\_rate}            & 0.148 & 0.125 & 0.200 & 0.250 & \crd{8}$-0.32$* & \crd{2}$-0.06$ & \cg{0}$+0.01$* \\
    \texttt{negation\_rate}                      & 1.053 & 0.826 & 1.476 & 1.961 & \crd{8}$-0.31$* & \crd{2}$-0.07$ & \crd{2}$-0.06$ \\
    \texttt{hedge\_informal\_count}              & 0.000 & 1.000 & 0.000 & 1.000 & \cg{8}$+0.31$* & \cg{2}$+0.09$ & \cg{15}$+0.59$* \\
    \texttt{hedge\_rate}                         & 2.778 & 4.918 & 2.944 & 6.790 & \cg{8}$+0.31$* & \crd{3}$-0.10$ & \cg{14}$+0.58$* \\
    \texttt{epistemic\_count}                    & 0.000 & 0.000 & 0.000 & 0.000 & \cg{6}$+0.23$* & \crd{2}$-0.08$ & \cg{3}$+0.13$* \\
    \texttt{discourse\_particle\_count}          & 0.000 & 1.000 & 0.000 & 2.000 & \cg{4}$+0.16$* & \crd{1}$-0.03$ & \cg{18}$+0.72$* \\
    \texttt{booster\_rate}                       & 0.000 & 0.855 & 0.000 & 0.000 & \cg{4}$+0.14$* & \cg{2}$+0.06$ & \cg{1}$+0.05$* \\
    \texttt{discourse\_particle\_rate}           & 0.000 & 0.813 & 0.000 & 2.326 & \crd{3}$-0.12$ & \crd{1}$-0.03$ & \cg{8}$+0.33$* \\
    \texttt{negation\_count}                     & 1.000 & 1.000 & 1.000 & 2.000 & \crd{1}$-0.05$ & \crd{3}$-0.11$ & \cg{12}$+0.47$* \\
    \texttt{hedge\_formal\_count}                & 0.000 & 0.000 & 0.000 & 0.000 & \cg{1}$+0.04$ & \cg{3}$+0.10$ & \crd{2}$-0.07$* \\
    \texttt{hedge\_informal\_rate}               & 0.000 & 0.855 & 0.000 & 1.493 & \crd{1}$-0.03$* & \crd{3}$-0.12$ & \cg{4}$+0.15$* \\
    \texttt{hedge\_formal\_rate}                 & 0.000 & 0.000 & 0.000 & 0.000 & \crd{1}$-0.03$ & \crd{1}$-0.04$ & \crd{2}$-0.08$* \\
    \texttt{epistemic\_rate}                     & 0.000 & 0.000 & 0.000 & 0.000 & \crd{0}$-0.00$* & \crd{1}$-0.04$ & \crd{1}$-0.04$ \\
    \addlinespace[1.5pt]
    \multicolumn{8}{l}{\textit{Syntactic complexity} (5)} \\
    \addlinespace[0.5pt]
    \texttt{n\_sentences}                        & 3.000 & 6.000 & 3.000 & 7.000 & \cg{19}$+0.77$* & \cg{3}$+0.10$ & \cg{23}$+0.91$* \\
    \texttt{tree\_depth\_max}                    & 4.000 & 6.000 & 4.000 & 5.000 & \cg{17}$+0.68$* & \cg{4}$+0.15$ & \cg{13}$+0.50$* \\
    \texttt{tree\_depth\_mean}                   & 3.000 & 3.667 & 3.450 & 3.143 & \cg{10}$+0.38$* & \cg{2}$+0.09$ & \crd{1}$-0.02$ \\
    \texttt{dep\_distance\_mean}                 & 2.971 & 3.174 & 3.138 & 3.273 & \cg{8}$+0.32$* & \cg{9}$+0.35$* & \cg{10}$+0.40$* \\
    \texttt{subordinate\_clause\_rate}           & 0.750 & 1.143 & 0.600 & 1.000 & \cg{7}$+0.28$* & \crd{4}$-0.17$ & \cg{3}$+0.10$* \\
    \addlinespace[1.5pt]
    \multicolumn{8}{l}{\textit{Named entities} (9)} \\
    \addlinespace[0.5pt]
    \texttt{ner\_total\_rate}                    & 0.000 & 0.000 & 0.000 & 0.000 & \crd{4}$-0.16$* & \cg{3}$+0.10$ & \crd{3}$-0.12$* \\
    \texttt{ner\_PRS\_rate}                      & 0.000 & 0.000 & 0.000 & 0.000 & \crd{3}$-0.11$* & \crd{3}$-0.10$ & \crd{2}$-0.09$* \\
    \texttt{ner\_TME\_rate}                      & 0.000 & 0.000 & 0.000 & 0.000 & \crd{3}$-0.10$* & \cg{1}$+0.05$ & \crd{2}$-0.08$* \\
    \texttt{ner\_LOC\_rate}                      & 0.000 & 0.000 & 0.000 & 0.000 & \crd{3}$-0.10$ & \cg{8}$+0.32$* & \crd{2}$-0.07$ \\
    \texttt{ner\_OBJ\_rate}                      & 0.000 & 0.000 & 0.000 & 0.000 & \cg{1}$+0.05$* & \crd{3}$-0.10$ & \cg{0}$+0.01$ \\
    \texttt{ner\_ORG\_rate}                      & 0.000 & 0.000 & 0.000 & 0.000 & \crd{1}$-0.03$ & \crd{3}$-0.10$ & \crd{0}$-0.00$* \\
    \texttt{ner\_MSR\_rate}                      & 0.000 & 0.000 & 0.000 & 0.000 & \crd{1}$-0.03$ & \crd{3}$-0.10$ & \cg{2}$+0.06$* \\
    \texttt{ner\_EVN\_rate}                      & 0.000 & 0.000 & 0.000 & 0.000 & \crd{1}$-0.03$ & \crd{0}$+0.00$ & \cg{1}$+0.03$ \\
    \texttt{ner\_WRK\_rate}                      & 0.000 & 0.000 & 0.000 & 0.000 & \crd{1}$-0.03$ & \crd{4}$-0.14$ & \crd{1}$-0.04$* \\
    \addlinespace[1.5pt]
    \multicolumn{8}{l}{\textit{Readability (textstat)} (7)} \\
    \addlinespace[0.5pt]
    \texttt{word\_count}                         & 31.000 & 91.000 & 37.000 & 87.000 & \cg{21}$+0.83$* & \cg{7}$+0.26$* & \cg{20}$+0.80$* \\
    \texttt{rough\_syllable\_count\_sv}          & 50.000 & 151.000 & 66.000 & 135.000 & \cg{20}$+0.81$* & \cg{9}$+0.35$* & \cg{18}$+0.73$* \\
    \texttt{char\_count}                         & 178.000 & 524.000 & 228.000 & 473.000 & \cg{20}$+0.80$* & \cg{8}$+0.33$* & \cg{18}$+0.73$* \\
    \texttt{textstat\_smog\_index}               & 8.842 & 10.019 & 10.126 & 9.121 & \cg{11}$+0.45$* & \cg{6}$+0.22$ & \cg{5}$+0.21$* \\
    \texttt{textstat\_gunning\_fog}              & 8.799 & 10.316 & 10.588 & 9.353 & \cg{4}$+0.16$* & \cg{4}$+0.17$ & \cg{1}$+0.02$* \\
    \texttt{textstat\_flesch\_kincaid\_grade}    & 6.790 & 7.882 & 8.667 & 7.064 & \cg{1}$+0.05$* & \cg{2}$+0.07$ & \crd{1}$-0.04$* \\
    \texttt{textstat\_flesch\_reading\_ease}     & 71.916 & 67.668 & 62.774 & 73.481 & \cg{1}$+0.04$* & \crd{1}$-0.02$ & \cg{3}$+0.13$* \\
    \addlinespace[1.5pt]
    \multicolumn{8}{l}{\textit{POS proportions} (17)} \\
    \addlinespace[0.5pt]
    \texttt{pos\_CCONJ}                          & 0.027 & 0.043 & 0.052 & 0.051 & \cg{10}$+0.38$* & \cg{11}$+0.43$* & \cg{17}$+0.66$* \\
    \texttt{pos\_ADJ}                            & 0.061 & 0.082 & 0.083 & 0.074 & \cg{7}$+0.27$* & \cg{3}$+0.11$ & \cg{4}$+0.15$* \\
    \texttt{pos\_PROPN}                          & 0.000 & 0.000 & 0.000 & 0.000 & \crd{6}$-0.23$* & \crd{1}$-0.05$ & \crd{5}$-0.18$* \\
    \texttt{pos\_SCONJ}                          & 0.021 & 0.033 & 0.025 & 0.035 & \cg{6}$+0.22$* & \crd{0}$-0.01$ & \cg{7}$+0.29$* \\
    \texttt{pos\_AUX}                            & 0.053 & 0.062 & 0.047 & 0.068 & \cg{4}$+0.17$* & \crd{4}$-0.14$ & \cg{7}$+0.29$* \\
    \texttt{pos\_NUM}                            & 0.000 & 0.000 & 0.000 & 0.000 & \crd{4}$-0.17$* & \crd{0}$-0.01$ & \crd{6}$-0.23$* \\
    \texttt{pos\_PART}                           & 0.019 & 0.029 & 0.017 & 0.032 & \cg{4}$+0.15$* & \cg{3}$+0.12$ & \cg{6}$+0.25$* \\
    \texttt{pos\_VERB}                           & 0.118 & 0.123 & 0.117 & 0.109 & \cg{4}$+0.14$* & \cg{1}$+0.04$ & \crd{3}$-0.10$* \\
    \texttt{pos\_NOUN}                           & 0.152 & 0.156 & 0.185 & 0.129 & \crd{3}$-0.11$ & \cg{3}$+0.11$ & \crd{9}$-0.36$* \\
    \texttt{pos\_DET}                            & 0.032 & 0.031 & 0.021 & 0.033 & \crd{3}$-0.10$ & \crd{9}$-0.36$* & \crd{2}$-0.06$ \\
    \texttt{pos\_ADV}                            & 0.103 & 0.103 & 0.115 & 0.120 & \crd{1}$-0.05$ & \cg{7}$+0.26$* & \cg{4}$+0.15$* \\
    \texttt{pos\_PRON}                           & 0.113 & 0.114 & 0.069 & 0.118 & \cg{1}$+0.05$* & \crd{10}$-0.39$* & \cg{2}$+0.08$* \\
    \texttt{pos\_PUNCT}                          & 0.105 & 0.104 & 0.112 & 0.125 & \crd{1}$-0.04$ & \cg{4}$+0.14$ & \cg{4}$+0.15$* \\
    \texttt{pos\_ADP}                            & 0.074 & 0.077 & 0.075 & 0.067 & \cg{1}$+0.04$* & \cg{0}$+0.01$ & \crd{3}$-0.12$* \\
    \texttt{pos\_SYM}                            & 0.000 & 0.000 & 0.000 & 0.000 & \crd{0}$-0.01$ & \crd{1}$-0.02$ & \crd{4}$-0.14$* \\
    \texttt{pos\_INTJ}                           & 0.000 & 0.000 & 0.000 & 0.000 & \crd{0}$-0.01$* & \crd{1}$-0.02$ & \crd{1}$-0.02$* \\
    \texttt{pos\_X}                              & 0.000 & 0.000 & 0.000 & 0.000 & \crd{0}$+0.00$ & \crd{0}$+0.00$ & \crd{0}$+0.00$ \\
    \bottomrule
  \end{tabular}
  \caption{Informal register, Baseline condition: every feature tested, human vs.\
  each generator (paired Wilcoxon signed-rank, Cohen's $d_z$, Benjamini--Hochberg FDR within each
  generator's family of 70 features). Positive (green) = generator $>$ human, negative (red) =
  human $>$ generator; shading $\propto |d_z|$, * denotes $q<0.05$. GPT-5.2 and Mistral-3.2 use the
  full corpus (170 formal / \textasciitilde{}1{,}150 informal); GPT-5.6 is a 100-document quiz
  subset, so its column has less statistical power. Features are grouped by extractor and ordered by
  $|d_z|$ vs.\ GPT-5.2 at baseline, held fixed across this register's three condition-tables so a row
  tracks the same feature. Source: \texttt{model\_comparison\_three\_way\_results.csv}.}
  \label{app:full-three-way-informal-baseline}
\end{table}

\clearpage

% ---- informal register, Human-like condition: all features, human vs 3 models (app:full-three-way-informal-human_like) ----
% REQUIRES: booktabs + the \cg / \crd shading macros from the preamble.
% 70 features x 3 models.
% Sized to fit one page at \scriptsize with arraystretch 0.9; if it still overflows, wrap in longtable.
\begin{table}[htbp]
  \centering
  \scriptsize
  \setlength{\tabcolsep}{2.5pt}
  \renewcommand{\arraystretch}{0.9}
  \begin{tabular}{l rrrr ccc}
    \toprule
    & \multicolumn{4}{c}{Median} & \multicolumn{3}{c}{Cohen's $d_z$ (vs.\ human)} \\
    \cmidrule(lr){2-5}\cmidrule(l){6-8}
    Feature & H & GPT-5.2 & GPT-5.6 & Mistral & GPT-5.2 & GPT-5.6 & Mistral \\
    \midrule
    \multicolumn{8}{l}{\textit{Stylometric surface} (14)} \\
    \addlinespace[0.5pt]
    \texttt{ttr}                                 & 0.882 & 0.668 & 0.889 & 0.647 & \crd{29}$-1.17$* & \crd{1}$-0.03$ & \crd{35}$-1.41$* \\
    \texttt{mattr}                               & 0.889 & 0.852 & 0.900 & 0.785 & \crd{7}$-0.28$* & \crd{1}$-0.04$ & \crd{22}$-0.86$* \\
    \texttt{mtld}                                & 71.680 & 118.696 & 98.747 & 72.020 & \cg{18}$+0.72$* & \cg{9}$+0.34$* & \crd{2}$-0.06$ \\
    \texttt{punct\_dash}                         & 0.000 & 0.376 & 0.000 & 0.000 & \cg{19}$+0.77$* & \cg{13}$+0.51$* & \cg{14}$+0.57$* \\
    \texttt{punct\_period}                       & 6.000 & 4.040 & 5.505 & 4.396 & \crd{8}$-0.30$* & \crd{8}$-0.32$* & \crd{7}$-0.27$* \\
    \texttt{bigram\_repetition\_rate}            & 0.000 & 0.025 & 0.000 & 0.053 & \cg{14}$+0.55$* & \crd{3}$-0.13$ & \cg{24}$+0.95$* \\
    \texttt{punct\_question}                     & 0.000 & 0.000 & 0.000 & 0.000 & \crd{7}$-0.29$* & \crd{6}$-0.23$ & \crd{5}$-0.21$* \\
    \texttt{punct\_exclaim}                      & 0.000 & 0.000 & 0.000 & 0.513 & \crd{4}$-0.15$* & \cg{6}$+0.22$ & \crd{0}$+0.00$* \\
    \texttt{punct\_comma}                        & 1.770 & 3.810 & 3.922 & 3.521 & \cg{7}$+0.29$* & \cg{11}$+0.44$* & \cg{5}$+0.21$* \\
    \texttt{punct\_ellipsis}                     & 0.000 & 0.000 & 0.000 & 0.000 & \crd{4}$-0.16$* & \crd{5}$-0.19$ & \crd{4}$-0.14$* \\
    \texttt{function\_word\_rate}                & 0.500 & 0.476 & 0.447 & 0.578 & \crd{2}$-0.09$* & \crd{2}$-0.09$ & \cg{16}$+0.65$* \\
    \texttt{punct\_colon}                        & 0.000 & 0.813 & 0.000 & 0.000 & \crd{1}$-0.03$* & \crd{5}$-0.19$ & \crd{4}$-0.16$* \\
    \texttt{punct\_semicolon}                    & 0.000 & 0.000 & 0.000 & 0.000 & \crd{3}$-0.11$* & \crd{4}$-0.15$ & \crd{3}$-0.12$* \\
    \texttt{trigram\_repetition\_rate}           & 0.000 & 0.000 & 0.000 & 0.010 & \cg{1}$+0.04$* & \crd{1}$-0.05$ & \cg{13}$+0.53$* \\
    \addlinespace[1.5pt]
    \multicolumn{8}{l}{\textit{Pragmatic markers} (18)} \\
    \addlinespace[0.5pt]
    \texttt{hedge\_neutral\_count}               & 1.000 & 4.000 & 1.000 & 6.000 & \cg{23}$+0.93$* & \cg{6}$+0.24$* & \cg{33}$+1.32$* \\
    \texttt{hedge\_count}                        & 1.000 & 8.000 & 2.000 & 8.000 & \cg{35}$+1.40$* & \cg{11}$+0.44$* & \cg{35}$+1.40$* \\
    \texttt{n\_alpha\_tokens}                    & 30.000 & 184.000 & 40.000 & 111.000 & \cg{38}$+1.53$* & \cg{8}$+0.31$* & \cg{31}$+1.23$* \\
    \texttt{booster\_count}                      & 0.000 & 1.000 & 0.000 & 1.000 & \cg{19}$+0.75$* & \cg{8}$+0.30$* & \cg{14}$+0.56$* \\
    \texttt{hedge\_neutral\_rate}                & 1.124 & 2.098 & 2.904 & 5.085 & \crd{1}$-0.04$* & \cg{5}$+0.19$ & \cg{17}$+0.69$* \\
    \texttt{negation\_sentence\_rate}            & 0.148 & 0.250 & 0.250 & 0.250 & \cg{2}$+0.06$* & \cg{1}$+0.02$ & \cg{1}$+0.04$* \\
    \texttt{negation\_rate}                      & 1.053 & 1.667 & 1.695 & 1.961 & \crd{4}$-0.14$ & \cg{2}$+0.06$ & \crd{2}$-0.07$ \\
    \texttt{hedge\_informal\_count}              & 0.000 & 4.000 & 1.000 & 2.000 & \cg{36}$+1.42$* & \cg{10}$+0.41$* & \cg{21}$+0.85$* \\
    \texttt{hedge\_rate}                         & 2.778 & 4.453 & 5.359 & 6.977 & \cg{5}$+0.19$* & \cg{4}$+0.16$* & \cg{16}$+0.65$* \\
    \texttt{epistemic\_count}                    & 0.000 & 0.000 & 0.000 & 0.000 & \cg{8}$+0.33$* & \cg{3}$+0.12$ & \cg{9}$+0.34$* \\
    \texttt{discourse\_particle\_count}          & 0.000 & 2.000 & 0.000 & 3.000 & \cg{16}$+0.65$* & \cg{5}$+0.18$ & \cg{23}$+0.91$* \\
    \texttt{booster\_rate}                       & 0.000 & 0.592 & 0.000 & 0.629 & \cg{1}$+0.03$* & \cg{3}$+0.10$ & \cg{2}$+0.07$* \\
    \texttt{discourse\_particle\_rate}           & 0.000 & 0.955 & 0.000 & 2.459 & \crd{3}$-0.10$ & \cg{4}$+0.16$ & \cg{9}$+0.35$* \\
    \texttt{negation\_count}                     & 1.000 & 3.000 & 1.000 & 2.000 & \cg{21}$+0.82$* & \crd{0}$-0.01$ & \cg{17}$+0.66$* \\
    \texttt{hedge\_formal\_count}                & 0.000 & 0.000 & 0.000 & 0.000 & \cg{5}$+0.21$* & \cg{4}$+0.14$ & \crd{1}$-0.03$ \\
    \texttt{hedge\_informal\_rate}               & 0.000 & 2.055 & 1.515 & 1.695 & \cg{8}$+0.33$* & \cg{2}$+0.07$* & \cg{6}$+0.22$* \\
    \texttt{hedge\_formal\_rate}                 & 0.000 & 0.000 & 0.000 & 0.000 & \crd{0}$-0.00$* & \crd{1}$-0.02$ & \crd{2}$-0.07$* \\
    \texttt{epistemic\_rate}                     & 0.000 & 0.000 & 0.000 & 0.000 & \crd{3}$-0.13$ & \cg{2}$+0.08$ & \cg{0}$+0.01$* \\
    \addlinespace[1.5pt]
    \multicolumn{8}{l}{\textit{Syntactic complexity} (5)} \\
    \addlinespace[0.5pt]
    \texttt{n\_sentences}                        & 3.000 & 10.000 & 3.000 & 8.000 & \cg{29}$+1.17$* & \cg{5}$+0.18$ & \cg{31}$+1.24$* \\
    \texttt{tree\_depth\_max}                    & 4.000 & 6.000 & 5.000 & 5.000 & \cg{26}$+1.05$* & \cg{5}$+0.20$ & \cg{16}$+0.65$* \\
    \texttt{tree\_depth\_mean}                   & 3.000 & 4.000 & 3.500 & 3.300 & \cg{15}$+0.60$* & \cg{3}$+0.13$ & \cg{3}$+0.13$* \\
    \texttt{dep\_distance\_mean}                 & 2.971 & 3.747 & 3.123 & 3.345 & \cg{22}$+0.88$* & \cg{7}$+0.29$* & \cg{12}$+0.48$* \\
    \texttt{subordinate\_clause\_rate}           & 0.750 & 1.167 & 0.667 & 1.143 & \cg{7}$+0.28$* & \crd{2}$-0.06$ & \cg{6}$+0.24$* \\
    \addlinespace[1.5pt]
    \multicolumn{8}{l}{\textit{Named entities} (9)} \\
    \addlinespace[0.5pt]
    \texttt{ner\_total\_rate}                    & 0.000 & 0.000 & 0.000 & 0.000 & \crd{4}$-0.15$* & \cg{1}$+0.03$ & \crd{4}$-0.15$* \\
    \texttt{ner\_PRS\_rate}                      & 0.000 & 0.000 & 0.000 & 0.000 & \crd{3}$-0.10$* & \crd{1}$-0.05$ & \crd{3}$-0.11$* \\
    \texttt{ner\_TME\_rate}                      & 0.000 & 0.000 & 0.000 & 0.000 & \crd{3}$-0.11$* & \crd{0}$+0.00$ & \crd{2}$-0.06$ \\
    \texttt{ner\_LOC\_rate}                      & 0.000 & 0.000 & 0.000 & 0.000 & \crd{2}$-0.08$ & \cg{5}$+0.20$ & \crd{2}$-0.09$ \\
    \texttt{ner\_OBJ\_rate}                      & 0.000 & 0.000 & 0.000 & 0.000 & \crd{0}$-0.01$ & \crd{3}$-0.10$ & \crd{0}$-0.01$ \\
    \texttt{ner\_ORG\_rate}                      & 0.000 & 0.000 & 0.000 & 0.000 & \crd{1}$-0.03$* & \crd{2}$-0.06$ & \crd{1}$-0.03$* \\
    \texttt{ner\_MSR\_rate}                      & 0.000 & 0.000 & 0.000 & 0.000 & \crd{1}$-0.03$ & \crd{3}$-0.10$ & \cg{2}$+0.06$* \\
    \texttt{ner\_EVN\_rate}                      & 0.000 & 0.000 & 0.000 & 0.000 & \cg{1}$+0.05$* & \crd{0}$+0.00$ & \cg{1}$+0.02$ \\
    \texttt{ner\_WRK\_rate}                      & 0.000 & 0.000 & 0.000 & 0.000 & \crd{1}$-0.04$ & \crd{4}$-0.14$ & \crd{1}$-0.03$ \\
    \addlinespace[1.5pt]
    \multicolumn{8}{l}{\textit{Readability (textstat)} (7)} \\
    \addlinespace[0.5pt]
    \texttt{word\_count}                         & 31.000 & 186.000 & 41.500 & 113.000 & \cg{38}$+1.52$* & \cg{8}$+0.32$* & \cg{31}$+1.23$* \\
    \texttt{rough\_syllable\_count\_sv}          & 50.000 & 307.000 & 70.000 & 175.000 & \cg{36}$+1.45$* & \cg{9}$+0.35$* & \cg{28}$+1.11$* \\
    \texttt{char\_count}                         & 178.000 & 1083.000 & 244.000 & 613.000 & \cg{36}$+1.45$* & \cg{9}$+0.35$* & \cg{28}$+1.12$* \\
    \texttt{textstat\_smog\_index}               & 8.842 & 10.290 & 9.516 & 8.842 & \cg{14}$+0.57$* & \cg{1}$+0.03$ & \cg{4}$+0.17$* \\
    \texttt{textstat\_gunning\_fog}              & 8.799 & 11.113 & 9.584 & 9.333 & \cg{8}$+0.33$* & \cg{1}$+0.04$ & \crd{0}$-0.01$* \\
    \texttt{textstat\_flesch\_kincaid\_grade}    & 6.790 & 9.018 & 7.377 & 7.002 & \cg{5}$+0.20$* & \crd{1}$-0.05$ & \crd{2}$-0.06$ \\
    \texttt{textstat\_flesch\_reading\_ease}     & 71.916 & 64.545 & 67.168 & 74.910 & \crd{1}$-0.02$* & \cg{3}$+0.11$ & \cg{4}$+0.16$* \\
    \addlinespace[1.5pt]
    \multicolumn{8}{l}{\textit{POS proportions} (17)} \\
    \addlinespace[0.5pt]
    \texttt{pos\_CCONJ}                          & 0.027 & 0.045 & 0.041 & 0.051 & \cg{12}$+0.49$* & \cg{6}$+0.25$* & \cg{16}$+0.64$* \\
    \texttt{pos\_ADJ}                            & 0.061 & 0.069 & 0.077 & 0.072 & \cg{2}$+0.07$* & \cg{0}$+0.01$ & \cg{3}$+0.11$* \\
    \texttt{pos\_PROPN}                          & 0.000 & 0.006 & 0.000 & 0.004 & \crd{4}$-0.14$ & \crd{4}$-0.15$ & \crd{6}$-0.24$* \\
    \texttt{pos\_SCONJ}                          & 0.021 & 0.034 & 0.029 & 0.038 & \cg{6}$+0.23$* & \crd{1}$-0.03$ & \cg{9}$+0.36$* \\
    \texttt{pos\_AUX}                            & 0.053 & 0.049 & 0.051 & 0.067 & \crd{3}$-0.11$* & \crd{1}$-0.03$ & \cg{8}$+0.31$* \\
    \texttt{pos\_NUM}                            & 0.000 & 0.004 & 0.000 & 0.000 & \crd{1}$-0.05$* & \crd{0}$-0.00$ & \crd{6}$-0.25$* \\
    \texttt{pos\_PART}                           & 0.019 & 0.020 & 0.019 & 0.031 & \crd{3}$-0.11$ & \cg{6}$+0.23$ & \cg{6}$+0.23$* \\
    \texttt{pos\_VERB}                           & 0.118 & 0.112 & 0.116 & 0.114 & \crd{2}$-0.07$* & \cg{3}$+0.11$ & \crd{1}$-0.04$* \\
    \texttt{pos\_NOUN}                           & 0.152 & 0.166 & 0.155 & 0.118 & \crd{0}$-0.01$* & \crd{2}$-0.07$ & \crd{13}$-0.50$* \\
    \texttt{pos\_DET}                            & 0.032 & 0.023 & 0.027 & 0.033 & \crd{7}$-0.26$* & \crd{8}$-0.32$* & \crd{2}$-0.07$ \\
    \texttt{pos\_ADV}                            & 0.103 & 0.116 & 0.143 & 0.127 & \cg{3}$+0.13$* & \cg{12}$+0.47$* & \cg{6}$+0.25$* \\
    \texttt{pos\_PRON}                           & 0.113 & 0.089 & 0.086 & 0.129 & \crd{7}$-0.28$* & \crd{6}$-0.25$* & \cg{6}$+0.22$* \\
    \texttt{pos\_PUNCT}                          & 0.105 & 0.145 & 0.109 & 0.118 & \cg{11}$+0.44$* & \crd{1}$-0.02$ & \cg{2}$+0.06$* \\
    \texttt{pos\_ADP}                            & 0.074 & 0.064 & 0.076 & 0.065 & \crd{5}$-0.20$* & \cg{2}$+0.09$ & \crd{4}$-0.15$* \\
    \texttt{pos\_SYM}                            & 0.000 & 0.016 & 0.000 & 0.000 & \cg{18}$+0.71$* & \crd{1}$-0.02$ & \crd{4}$-0.15$* \\
    \texttt{pos\_INTJ}                           & 0.000 & 0.000 & 0.000 & 0.000 & \crd{2}$-0.07$* & \crd{1}$-0.04$ & \crd{1}$-0.03$* \\
    \texttt{pos\_X}                              & 0.000 & 0.000 & 0.000 & 0.000 & \crd{0}$+0.00$ & \crd{0}$+0.00$ & \crd{0}$+0.00$ \\
    \bottomrule
  \end{tabular}
  \caption{Informal register, Human-like condition: every feature tested, human vs.\
  each generator (paired Wilcoxon signed-rank, Cohen's $d_z$, Benjamini--Hochberg FDR within each
  generator's family of 70 features). Positive (green) = generator $>$ human, negative (red) =
  human $>$ generator; shading $\propto |d_z|$, * denotes $q<0.05$. GPT-5.2 and Mistral-3.2 use the
  full corpus (170 formal / \textasciitilde{}1{,}150 informal); GPT-5.6 is a 100-document quiz
  subset, so its column has less statistical power. Features are grouped by extractor and ordered by
  $|d_z|$ vs.\ GPT-5.2 at baseline, held fixed across this register's three condition-tables so a row
  tracks the same feature. Source: \texttt{model\_comparison\_three\_way\_results.csv}.}
  \label{app:full-three-way-informal-human_like}
\end{table}

\clearpage

% ---- informal register, Detector-evasive condition: all features, human vs 3 models (app:full-three-way-informal-detector_evasive) ----
% REQUIRES: booktabs + the \cg / \crd shading macros from the preamble.
% 70 features x 3 models.
% Sized to fit one page at \scriptsize with arraystretch 0.9; if it still overflows, wrap in longtable.
\begin{table}[htbp]
  \centering
  \scriptsize
  \setlength{\tabcolsep}{2.5pt}
  \renewcommand{\arraystretch}{0.9}
  \begin{tabular}{l rrrr ccc}
    \toprule
    & \multicolumn{4}{c}{Median} & \multicolumn{3}{c}{Cohen's $d_z$ (vs.\ human)} \\
    \cmidrule(lr){2-5}\cmidrule(l){6-8}
    Feature & H & GPT-5.2 & GPT-5.6 & Mistral & GPT-5.2 & GPT-5.6 & Mistral \\
    \midrule
    \multicolumn{8}{l}{\textit{Stylometric surface} (14)} \\
    \addlinespace[0.5pt]
    \texttt{ttr}                                 & 0.882 & 0.717 & 0.930 & 0.651 & \crd{22}$-0.89$* & \cg{5}$+0.21$ & \crd{31}$-1.25$* \\
    \texttt{mattr}                               & 0.889 & 0.863 & 0.932 & 0.773 & \crd{5}$-0.19$* & \cg{5}$+0.18$ & \crd{23}$-0.92$* \\
    \texttt{mtld}                                & 71.680 & 127.764 & 129.227 & 63.093 & \cg{20}$+0.80$* & \cg{12}$+0.46$* & \crd{5}$-0.19$* \\
    \texttt{punct\_dash}                         & 0.000 & 0.000 & 0.000 & 0.000 & \cg{4}$+0.17$* & \cg{4}$+0.14$ & \cg{5}$+0.20$* \\
    \texttt{punct\_period}                       & 6.000 & 11.504 & 14.286 & 10.204 & \cg{5}$+0.20$* & \cg{17}$+0.68$* & \cg{3}$+0.13$* \\
    \texttt{bigram\_repetition\_rate}            & 0.000 & 0.014 & 0.000 & 0.053 & \cg{6}$+0.22$* & \crd{7}$-0.28$* & \cg{19}$+0.74$* \\
    \texttt{punct\_question}                     & 0.000 & 0.535 & 1.802 & 0.909 & \crd{5}$-0.20$ & \cg{6}$+0.25$* & \crd{2}$-0.09$* \\
    \texttt{punct\_exclaim}                      & 0.000 & 0.000 & 0.000 & 0.000 & \crd{4}$-0.16$* & \crd{1}$-0.03$ & \crd{3}$-0.13$* \\
    \texttt{punct\_comma}                        & 1.770 & 1.282 & 1.802 & 1.980 & \crd{6}$-0.23$* & \crd{4}$-0.15$ & \crd{4}$-0.15$ \\
    \texttt{punct\_ellipsis}                     & 0.000 & 0.360 & 1.818 & 0.893 & \cg{0}$+0.01$* & \cg{23}$+0.90$* & \cg{6}$+0.25$* \\
    \texttt{function\_word\_rate}                & 0.500 & 0.444 & 0.386 & 0.578 & \crd{7}$-0.26$* & \crd{17}$-0.66$* & \cg{15}$+0.61$* \\
    \texttt{punct\_colon}                        & 0.000 & 0.000 & 0.000 & 0.000 & \crd{4}$-0.16$* & \crd{5}$-0.18$* & \crd{4}$-0.16$* \\
    \texttt{punct\_semicolon}                    & 0.000 & 0.000 & 0.000 & 0.000 & \crd{3}$-0.12$* & \crd{4}$-0.15$ & \crd{3}$-0.11$* \\
    \texttt{trigram\_repetition\_rate}           & 0.000 & 0.000 & 0.000 & 0.009 & \crd{1}$-0.05$ & \crd{5}$-0.20$ & \cg{10}$+0.40$* \\
    \addlinespace[1.5pt]
    \multicolumn{8}{l}{\textit{Pragmatic markers} (18)} \\
    \addlinespace[0.5pt]
    \texttt{hedge\_neutral\_count}               & 1.000 & 2.000 & 0.000 & 4.000 & \cg{14}$+0.55$* & \crd{5}$-0.20$ & \cg{24}$+0.95$* \\
    \texttt{hedge\_count}                        & 1.000 & 4.000 & 1.000 & 6.000 & \cg{20}$+0.80$* & \crd{1}$-0.04$ & \cg{25}$+1.01$* \\
    \texttt{n\_alpha\_tokens}                    & 30.000 & 150.000 & 37.500 & 101.000 & \cg{34}$+1.37$* & \cg{6}$+0.23$* & \cg{19}$+0.74$* \\
    \texttt{booster\_count}                      & 0.000 & 0.000 & 0.000 & 0.000 & \crd{3}$-0.13$* & \crd{5}$-0.18$ & \cg{2}$+0.08$* \\
    \texttt{hedge\_neutral\_rate}                & 1.124 & 1.504 & 0.000 & 4.545 & \crd{5}$-0.21$* & \crd{6}$-0.25$* & \cg{14}$+0.56$* \\
    \texttt{negation\_sentence\_rate}            & 0.148 & 0.176 & 0.183 & 0.312 & \crd{5}$-0.19$* & \crd{5}$-0.19$ & \cg{5}$+0.21$* \\
    \texttt{negation\_rate}                      & 1.053 & 2.752 & 2.440 & 3.571 & \cg{3}$+0.11$* & \cg{3}$+0.13$ & \cg{7}$+0.28$* \\
    \texttt{hedge\_informal\_count}              & 0.000 & 1.000 & 0.000 & 1.000 & \cg{18}$+0.73$* & \cg{5}$+0.18$ & \cg{16}$+0.64$* \\
    \texttt{hedge\_rate}                         & 2.778 & 2.830 & 2.859 & 6.349 & \crd{4}$-0.16$* & \crd{5}$-0.18$ & \cg{13}$+0.50$* \\
    \texttt{epistemic\_count}                    & 0.000 & 0.000 & 0.000 & 0.000 & \cg{1}$+0.02$ & \crd{4}$-0.15$ & \cg{2}$+0.06$* \\
    \texttt{discourse\_particle\_count}          & 0.000 & 0.000 & 0.000 & 2.000 & \crd{1}$-0.02$ & \crd{6}$-0.25$* & \cg{16}$+0.65$* \\
    \texttt{booster\_rate}                       & 0.000 & 0.000 & 0.000 & 0.000 & \crd{6}$-0.22$* & \crd{6}$-0.22$ & \crd{3}$-0.12$* \\
    \texttt{discourse\_particle\_rate}           & 0.000 & 0.000 & 0.000 & 2.222 & \crd{9}$-0.36$* & \crd{6}$-0.22$* & \cg{8}$+0.32$* \\
    \texttt{negation\_count}                     & 1.000 & 4.000 & 1.000 & 3.000 & \cg{28}$+1.10$* & \cg{1}$+0.05$ & \cg{14}$+0.54$* \\
    \texttt{hedge\_formal\_count}                & 0.000 & 0.000 & 0.000 & 0.000 & \cg{3}$+0.10$* & \cg{4}$+0.17$ & \crd{1}$-0.04$ \\
    \texttt{hedge\_informal\_rate}               & 0.000 & 1.058 & 0.000 & 1.515 & \crd{0}$+0.00$* & \crd{1}$-0.05$ & \cg{3}$+0.12$* \\
    \texttt{hedge\_formal\_rate}                 & 0.000 & 0.000 & 0.000 & 0.000 & \crd{1}$-0.03$ & \cg{1}$+0.03$ & \crd{2}$-0.07$* \\
    \texttt{epistemic\_rate}                     & 0.000 & 0.000 & 0.000 & 0.000 & \crd{5}$-0.21$* & \crd{3}$-0.10$ & \crd{4}$-0.16$* \\
    \addlinespace[1.5pt]
    \multicolumn{8}{l}{\textit{Syntactic complexity} (5)} \\
    \addlinespace[0.5pt]
    \texttt{n\_sentences}                        & 3.000 & 21.000 & 5.000 & 11.000 & \cg{41}$+1.62$* & \cg{23}$+0.93$* & \cg{27}$+1.08$* \\
    \texttt{tree\_depth\_max}                    & 4.000 & 5.000 & 4.000 & 4.000 & \cg{7}$+0.27$* & \crd{5}$-0.18$ & \cg{2}$+0.09$* \\
    \texttt{tree\_depth\_mean}                   & 3.000 & 2.333 & 2.500 & 2.667 & \crd{15}$-0.60$* & \crd{15}$-0.61$* & \crd{10}$-0.38$* \\
    \texttt{dep\_distance\_mean}                 & 2.971 & 2.649 & 2.593 & 2.781 & \crd{8}$-0.30$* & \crd{12}$-0.48$* & \crd{4}$-0.14$* \\
    \texttt{subordinate\_clause\_rate}           & 0.750 & 0.333 & 0.236 & 0.600 & \crd{15}$-0.60$* & \crd{17}$-0.69$* & \crd{8}$-0.31$* \\
    \addlinespace[1.5pt]
    \multicolumn{8}{l}{\textit{Named entities} (9)} \\
    \addlinespace[0.5pt]
    \texttt{ner\_total\_rate}                    & 0.000 & 0.957 & 2.198 & 0.000 & \cg{2}$+0.06$* & \cg{9}$+0.35$* & \crd{2}$-0.08$ \\
    \texttt{ner\_PRS\_rate}                      & 0.000 & 0.000 & 0.000 & 0.000 & \crd{0}$-0.00$* & \cg{0}$+0.01$ & \crd{2}$-0.07$ \\
    \texttt{ner\_TME\_rate}                      & 0.000 & 0.000 & 0.000 & 0.000 & \crd{1}$-0.05$ & \cg{3}$+0.13$ & \crd{1}$-0.03$ \\
    \texttt{ner\_LOC\_rate}                      & 0.000 & 0.000 & 0.000 & 0.000 & \cg{1}$+0.05$* & \cg{9}$+0.36$* & \crd{1}$-0.04$ \\
    \texttt{ner\_OBJ\_rate}                      & 0.000 & 0.000 & 0.000 & 0.000 & \cg{1}$+0.02$* & \crd{2}$-0.08$ & \cg{0}$+0.01$ \\
    \texttt{ner\_ORG\_rate}                      & 0.000 & 0.000 & 0.000 & 0.000 & \cg{4}$+0.15$* & \cg{5}$+0.21$ & \cg{1}$+0.03$* \\
    \texttt{ner\_MSR\_rate}                      & 0.000 & 0.000 & 0.000 & 0.000 & \crd{1}$-0.03$ & \crd{3}$-0.10$ & \crd{0}$-0.01$ \\
    \texttt{ner\_EVN\_rate}                      & 0.000 & 0.000 & 0.000 & 0.000 & \cg{1}$+0.02$ & \crd{0}$+0.00$ & \cg{0}$+0.01$ \\
    \texttt{ner\_WRK\_rate}                      & 0.000 & 0.000 & 0.000 & 0.000 & \crd{1}$-0.02$ & \crd{4}$-0.14$ & \crd{1}$-0.03$ \\
    \addlinespace[1.5pt]
    \multicolumn{8}{l}{\textit{Readability (textstat)} (7)} \\
    \addlinespace[0.5pt]
    \texttt{word\_count}                         & 31.000 & 152.000 & 38.500 & 102.000 & \cg{34}$+1.37$* & \cg{6}$+0.24$* & \cg{19}$+0.74$* \\
    \texttt{rough\_syllable\_count\_sv}          & 50.000 & 243.000 & 64.000 & 155.000 & \cg{32}$+1.29$* & \cg{7}$+0.29$* & \cg{18}$+0.70$* \\
    \texttt{char\_count}                         & 178.000 & 870.000 & 223.000 & 550.000 & \cg{33}$+1.31$* & \cg{8}$+0.31$* & \cg{17}$+0.69$* \\
    \texttt{textstat\_smog\_index}               & 8.842 & 7.264 & 7.554 & 7.032 & \crd{10}$-0.39$* & \crd{12}$-0.48$* & \crd{11}$-0.44$* \\
    \texttt{textstat\_gunning\_fog}              & 8.799 & 5.966 & 6.351 & 5.886 & \crd{16}$-0.62$* & \crd{19}$-0.76$* & \crd{16}$-0.63$* \\
    \texttt{textstat\_flesch\_kincaid\_grade}    & 6.790 & 4.401 & 4.875 & 3.946 & \crd{10}$-0.39$* & \crd{12}$-0.47$* & \crd{11}$-0.44$* \\
    \texttt{textstat\_flesch\_reading\_ease}     & 71.916 & 79.525 & 74.372 & 83.994 & \cg{6}$+0.25$* & \cg{6}$+0.25$* & \cg{8}$+0.32$* \\
    \addlinespace[1.5pt]
    \multicolumn{8}{l}{\textit{POS proportions} (17)} \\
    \addlinespace[0.5pt]
    \texttt{pos\_CCONJ}                          & 0.027 & 0.045 & 0.045 & 0.049 & \cg{12}$+0.48$* & \cg{6}$+0.24$* & \cg{14}$+0.57$* \\
    \texttt{pos\_ADJ}                            & 0.061 & 0.061 & 0.083 & 0.064 & \crd{2}$-0.07$ & \crd{0}$+0.00$ & \cg{0}$+0.01$* \\
    \texttt{pos\_PROPN}                          & 0.000 & 0.018 & 0.032 & 0.009 & \cg{4}$+0.15$* & \cg{5}$+0.21$* & \crd{2}$-0.08$ \\
    \texttt{pos\_SCONJ}                          & 0.021 & 0.027 & 0.014 & 0.031 & \cg{1}$+0.02$* & \crd{8}$-0.33$* & \cg{3}$+0.11$* \\
    \texttt{pos\_AUX}                            & 0.053 & 0.043 & 0.040 & 0.071 & \crd{6}$-0.23$* & \crd{5}$-0.19$ & \cg{9}$+0.35$* \\
    \texttt{pos\_NUM}                            & 0.000 & 0.007 & 0.000 & 0.000 & \cg{2}$+0.07$* & \cg{1}$+0.05$ & \crd{5}$-0.19$* \\
    \texttt{pos\_PART}                           & 0.019 & 0.024 & 0.021 & 0.036 & \cg{0}$+0.01$* & \cg{3}$+0.13$ & \cg{10}$+0.38$* \\
    \texttt{pos\_VERB}                           & 0.118 & 0.125 & 0.118 & 0.112 & \cg{4}$+0.17$* & \cg{2}$+0.06$ & \crd{2}$-0.06$* \\
    \texttt{pos\_NOUN}                           & 0.152 & 0.187 & 0.198 & 0.116 & \cg{5}$+0.20$* & \cg{6}$+0.24$* & \crd{12}$-0.48$* \\
    \texttt{pos\_DET}                            & 0.032 & 0.023 & 0.024 & 0.028 & \crd{7}$-0.26$* & \crd{11}$-0.42$* & \crd{4}$-0.16$* \\
    \texttt{pos\_ADV}                            & 0.103 & 0.099 & 0.114 & 0.121 & \crd{2}$-0.09$* & \cg{7}$+0.26$* & \cg{5}$+0.20$* \\
    \texttt{pos\_PRON}                           & 0.113 & 0.085 & 0.056 & 0.127 & \crd{9}$-0.35$* & \crd{18}$-0.71$* & \cg{5}$+0.19$* \\
    \texttt{pos\_PUNCT}                          & 0.105 & 0.153 & 0.147 & 0.125 & \cg{12}$+0.47$* & \cg{14}$+0.54$* & \cg{4}$+0.17$* \\
    \texttt{pos\_ADP}                            & 0.074 & 0.071 & 0.071 & 0.058 & \crd{1}$-0.05$ & \crd{1}$-0.05$ & \crd{8}$-0.31$* \\
    \texttt{pos\_SYM}                            & 0.000 & 0.000 & 0.000 & 0.000 & \crd{2}$-0.09$ & \crd{2}$-0.09$ & \crd{4}$-0.15$* \\
    \texttt{pos\_INTJ}                           & 0.000 & 0.000 & 0.000 & 0.000 & \crd{2}$-0.08$ & \cg{2}$+0.07$ & \cg{0}$+0.01$* \\
    \texttt{pos\_X}                              & 0.000 & 0.000 & 0.000 & 0.000 & \crd{0}$+0.00$ & \crd{0}$+0.00$ & \crd{0}$+0.00$ \\
    \bottomrule
  \end{tabular}
  \caption{Informal register, Detector-evasive condition: every feature tested, human vs.\
  each generator (paired Wilcoxon signed-rank, Cohen's $d_z$, Benjamini--Hochberg FDR within each
  generator's family of 70 features). Positive (green) = generator $>$ human, negative (red) =
  human $>$ generator; shading $\propto |d_z|$, * denotes $q<0.05$. GPT-5.2 and Mistral-3.2 use the
  full corpus (170 formal / \textasciitilde{}1{,}150 informal); GPT-5.6 is a 100-document quiz
  subset, so its column has less statistical power. Features are grouped by extractor and ordered by
  $|d_z|$ vs.\ GPT-5.2 at baseline, held fixed across this register's three condition-tables so a row
  tracks the same feature. Source: \texttt{model\_comparison\_three\_way\_results.csv}.}
  \label{app:full-three-way-informal-detector_evasive}
\end{table}

\end{document}


\end{document}